\documentclass[11pt,a4paper]{article}
\usepackage[margin=2.3cm]{geometry}
\usepackage{amsmath,amssymb,amsthm}
\usepackage[T1]{fontenc}
\usepackage{lmodern}
\usepackage{booktabs}
\usepackage[hidelinks]{hyperref}
\DeclareMathOperator{\arcosh}{arcosh}
\DeclareMathOperator{\arsinh}{arsinh}
\DeclareMathOperator{\spec}{spec}
\DeclareMathOperator{\spn}{span}
\DeclareMathOperator{\Var}{Var}
\DeclareMathOperator{\dist}{dist}
\newcommand{\one}{\mathbf 1}
\newcommand{\R}{\mathbb R}
\newcommand{\Z}{\mathbb Z}
\newcommand{\E}{\mathbb E}
\newcommand{\PP}{\mathbb P}
\newcommand{\ip}[2]{\langle #1,#2\rangle}
\newcommand{\Xb}{\bar X}
\theoremstyle{plain}
\newtheorem{theorem}{Theorem}[section]
\newtheorem{lemma}[theorem]{Lemma}
\newtheorem{proposition}[theorem]{Proposition}
\newtheorem{corollary}[theorem]{Corollary}
\theoremstyle{definition}
\newtheorem{definition}[theorem]{Definition}
\newtheorem{remark}[theorem]{Remark}
\newtheorem{observation}[theorem]{Numerical observation}
\newtheorem{conjecture}[theorem]{Conjecture}
\numberwithin{equation}{section}
\setlength{\parskip}{0.35em}
\emergencystretch=3em
\relpenalty=150
\binoppenalty=250
\title{Rigorous spectral structure of first-passage times to a point target\\ on the reflecting hypercubic lattice, $d\ge 2$\\[4pt]
\large poles, interlacing, residues, and explicit bounds on the principal killed eigenvalue}
\author{Xiaoxiao Zhouyi\\[2pt]\normalsize Companion note R2}
\date{1 October 2026}
\begin{document}
\setlength{\emergencystretch}{3em}
\maketitle
\medskip\noindent{\small\textit{Note.} This document is the companion note R2 of the article ``Mean and most probable first-passage times of lazy random walks in reflecting hypercubic lattices'' (\texttt{paper.pdf}, Section~7; Supplementary Material, Section~S8), in the repository \texttt{https://\allowbreak github.com/\allowbreak zhouyi-xiaoxiao/\allowbreak master-dissertation}. In the repository, its scripts are in \texttt{code/msc\_rigorous\_spectral/} and its data in \texttt{data/msc\_rigorous\_spectral/}; paths written \texttt{scripts/\dots} and \texttt{data/\dots} below refer to these folders; the paths in the \texttt{\%~src} comments of the \LaTeX{} source are relative to the root of the repository, and the run logs they cite are in \texttt{data/msc\_rigorous\_*/logs/}. Other run logs, the code of the internal checks and working notes mentioned below are not part of the repository; \texttt{notes/VERIFICATION.md} summarises how this note was checked.}\par\medskip


\begin{abstract}
\noindent For the lazy nearest-neighbour walk on $\{1,\dots,N\}^d$ with reflecting (cancelled-move) walls and a single absorbing site, we replace the formal pole analysis of the article by theorems. Part~I (any finite symmetric chain): the first-passage law is a finite signed mixture of exponentials (geometrics) whose rates are exactly the eigenvalues of the killed generator in the cyclic subspace of the target; they strictly interlace the eigenvalues of the reflecting walk that are visible from the target; degeneracies and eigenfunctions vanishing at the target are treated completely; the residues sum to one, the principal residue is positive, its quasi-stationary average is one, there are at least $\mathrm{dist}(x_0,a)-1$ sign changes, and for a start--target pair exchanged by a symmetry the sign of every residue is characterised (in particular $r_1<0$ for opposite corners, all $N$, $d\ge2$; and for opposite corners the residues alternate in sign if and only if $d=1$ or $N=2$). Part~II: explicit two-sided bounds on the principal eigenvalue $\sigma_0$, the second pole $\sigma_1$, the gap, and the residues $r_0,r_1$ in terms of three functionals of the lattice Green's function at the target. Part~III: for opposite corners, $r_0>1$ and $\sigma_0\cdot\mathrm{MFPT}>1$ in every dimension and for every $N$ (monotone coupling), and the functionals are evaluated: in $d=2$ with an explicit $O(1)$ remainder (logarithmic growth) and the limiting constant term, in $d=3$ with the Watson-type constant $C_3=\sum_{y\in\{0,1\}^3}u(y)$, the second-order constant, an explicit $O(1)$ remainder for every $N$ and the limiting constant term, and in every dimension up to explicit constants. Part~IV: a bound on the contribution of all neglected poles, certified tables, and a list of what remains unproved (notably: any statement about the mode itself). Every lemma is accompanied by a numerical check; constants are certified by ball arithmetic.
\end{abstract}

\section*{Summary of results and their status}
\addcontentsline{toc}{section}{Summary of results and their status}
Unit-rate convention: $L=I-P_1$; for activity $q$ the continuous-time rates are $q\sigma_j$ and the discrete-time multipliers $\nu_j=1-q\sigma_j$. $\tau$ is the mean hitting time of the target from a uniform start, $m$ the MFPT from $x_0$, $\Lambda_1$ the smallest non-zero eigenvalue of the reflecting walk visible from the target, $\Xb=\Lambda_1\tau$.
\begin{center}\small
\begin{tabular}{p{0.60\textwidth}p{0.15\textwidth}p{0.19\textwidth}}\toprule
Statement & where & status\\\midrule
$f(t)=\sum_jr_j\sigma_je^{-\sigma_jt}$, $f(t)=\sum_jr_j(1-\nu_j)\nu_j^{t-1}$; the rates $\sigma_j$ are the zeros of $G(s)=\ip{e_a}{(L-s)^{-1}e_a}$ $=$ eigenvalues of $L_a$ on $K_0$ ($\sigma_j$ is absent from the law started at $x$ iff $G_x(\sigma_j)=0$); $\det(L_a-s)=G(s)\det(L-s)$ & Thms \ref{thm:structure}, \ref{thm:fpt} & proved\\
strict interlacing $\Lambda_j<\sigma_j<\Lambda_{j+1}$ in the visible/symmetric sector; multiplicities; equality cases & Thm \ref{thm:interlace}, Cor.\ \ref{cor:box} & proved\\
$\sum r_j=1$; $r_0>0$; $\ge D-1$ sign changes; full alternation in $d=1$; quasi-stationary averages $\sum_x\nu(x)r_j(x)=\delta_{j0}$, $\sum_x\nu(x)m(x)=1/\sigma_0$; sign rule for symmetric pairs; $r_1<0$ (CC, $d\ge2$, all $N$); stationary start: all weights $>0$; $\sum_j|r_j|\le\sqrt{n-1}$ & Thm \ref{thm:residues}, \ref{thm:r1}, Props.\ \ref{prop:uniform}, \ref{prop:qsd} & proved\\
$\tau\le1/\sigma_0\le\tau+1/\Lambda_1$; $\sigma_0\tau=1-\kappa+O(\kappa^2)$ two-sided; bounds for $\sigma_1$, $\sigma_1-\sigma_0$, $r_0$, $r_1$ (general chain) & Thms \ref{thm:sigma0}--\ref{thm:r1size} & proved\\
sign monotonicity across eigenvalues of pure parity; CC: the residues alternate \emph{iff} $d=1$ or $N=2$ (the conjecture ``alternation in $d\ge2$'' is false) & Prop.\ \ref{prop:pure}, Thm \ref{thm:noalt} & proved\\
CC, any $d\ge1$, any $N$: $r_0>1$, $1<\sigma_0m$, $\tau+\frac1{\Lambda_1}\le m$; $\frac13d(N^2-1)\le m-\tau\le\frac{1+\ln W}{\Lambda_1}$ & Lem.\ \ref{lem:arrival}, \ref{lem:coupling}, Thm \ref{thm:r0gt1} & proved\\
CC, any $d\ge2$: $1+\frac{\mu_--1}{\Xb+1}\le\sigma_0m<1+\frac{1+\ln2d}{\Xb}$; all bounds explicit in $\Xb$ & Thm \ref{thm:cc} & proved (uses \eqref{eq:Zvalues} if $Z=Z_d$)\\
$Z_2,Z_3,Z_2',Z_3'$ (Epstein-type sums), $16$ certified decimals & Lem.\ \ref{lem:epstein} & computer-assisted\\
$d=2$: $\tau_2=\frac8\pi N^2\ln N+c_\tau N^2+R$, $-2.14\le R\le3.53$; $m_2-\tau_2=\frac{4\ln2}\pi N^2+R'$ & Thm \ref{thm:d2} & computer-assisted (constants)\\
$d=2$: $R\to c^{(0)}_\tau=0.8830352\ldots$, $R'\to c^{(0)}_{m\tau}=-0.3508943\ldots$ (constant term $C_2^{(0)}=0.5321408\ldots$ of the MFPT) & Thm \ref{thm:d2limits} & proved (limit); values computer-assisted\\
$d=3$: $\tau_3=C_3N^3+c_{\tau3}N^2+R_3$, $-17.4\le R_3\le19.4$, $m_3-\tau_3=c_{m\tau3}N^2+R_3'$, $|R_3'+1|\le81.4$, all $N\ge2$; $C_3=4.320460169373\ldots$ & Thm \ref{thm:d3}(a)--(c) & computer-assisted (constants)\\
$d=3$: $R_3\to1.7573985\ldots$, $R_3'\to-0.2989172\ldots$ (constant term $C_3''=1.4584812\ldots$ of the MFPT) & Thm \ref{thm:d3}(d) & proved (limit); values computer-assisted\\
$N$-explicit bounds on $\Xb$, $\sigma_0,\sigma_1,r_0,r_1$ for $d=2,3$ & Cor.\ \ref{cor:Xbar}, \ref{cor:explicitN} & computer-assisted\\
these bounds with truncated decimal constants are false for $N\ge8872$ ($d=2$) and $N\ge182\,420$ ($d=3$) & Prop.\ \ref{prop:roundedfalse} & computer-assisted (refutation)\\
first-order pole structure (the formal relations of the computations of Sections~5--6 of the article) & Thm \ref{thm:asym} & proved (asymptotic)\\
tail of all neglected poles, continuous and discrete time & Thm \ref{thm:tail}, Cor.\ \ref{cor:tailcc} & proved\\
finite-$N$ enclosures of $\sigma_0,\sigma_1,r_0,r_1,\tau,m$ & Section \ref{sec:cert} & computer-assisted\\
certified sign patterns of all residues for $20$ pairs $(d,N)$ & Prop.\ \ref{prop:signs} & computer-assisted\\
sign statistics for larger $N$; monotonicity of the remainders in $N$ and their rate of convergence & Obs.\ \ref{obs:signs}, \ref{obs:d2}, Rem.\ \ref{rem:d3sharp} & numerical observation\\
location of the mode; unimodality & Section \ref{sec:gaps} & not proved here\\\bottomrule
\end{tabular}
\end{center}

\paragraph{Convention for decimals.} A decimal expansion followed by dots ($\ldots$) is \emph{truncated}: every digit shown is a digit of the exact value (the list of all such strings is checked by script \texttt{21}). Decimals in inequalities are exact rational numbers.

\tableofcontents
\bigskip
\part*{Part I. Structure (any finite symmetric chain)}
\section{Setting, conventions, standing assumptions}\label{sec:setting}

\subsection{An abstract symmetric chain}
Let $\Omega$ be a finite set with $n=|\Omega|\ge 2$ elements and let $L=(L_{xy})_{x,y\in\Omega}$ be a real matrix with
\begin{equation}\label{eq:A}
\begin{gathered}\text{(A1) } L=L^{\mathsf T};\qquad \text{(A2) } L_{xy}\le 0\ (x\ne y),\ \ L\one=0;\\
\text{(A3) the graph } \{x\sim y \iff L_{xy}<0\}\text{ is connected.}\end{gathered}
\end{equation}
Then $\ip f{Lf}=\tfrac12\sum_{x\neq y}|L_{xy}|\,(f(x)-f(y))^2\ge 0$, so $L$ is positive semidefinite, and by (A3) its kernel is spanned by $\one$. We write $e_x$ for the coordinate vectors, $\ip\cdot\cdot$ for the Euclidean inner product on $\R^\Omega$, and
\[
0=\lambda^{(0)}<\lambda^{(1)}<\dots<\lambda^{(M)}
\]
for the \emph{distinct} eigenvalues of $L$, with orthogonal spectral projectors $E_0,\dots,E_M$ and multiplicities $m_i=\operatorname{rank}E_i$; thus $E_0=\frac1n\one\one^{\mathsf T}$, $m_0=1$.

Fix $q>0$. The \emph{continuous-time walk} is the Markov chain with generator $-qL$; the \emph{discrete-time walk} has transition matrix $P=I-qL$, which is a symmetric stochastic matrix as soon as $q\max_x L_{xx}\le 1$. For a target $a\in\Omega$ let $T=T_a$ be the first hitting time of $a$, and write $\Omega'=\Omega\setminus\{a\}$ and
\[
L_a=(L_{xy})_{x,y\in\Omega'}
\]
for the principal submatrix (the \emph{killed generator}). The substochastic matrix of the discrete walk on $\Omega'$ is $Q=I-qL_a$. For $x\in\Omega'$,
\begin{equation}\label{eq:surv}
\begin{aligned}\PP_x(T>t)&=\big(e^{-qtL_a}\one\big)(x)&&(t\ge0,\ \text{continuous time}),\\
\PP_x(T>t)&=\big(Q^{t}\one\big)(x)&&(t=0,1,2,\dots,\ \text{discrete time}),\end{aligned}
\end{equation}
which is the definition $S(t)=e_{x_0}^{\mathsf T}Q^t\one$ of the article. All statements below are written for $L$ (``unit rate''); the activity $q$ only rescales: continuous-time rates are $q\sigma$ and discrete-time multipliers are $\nu=1-q\sigma$.

\subsection{The reflecting box}
For integers $d\ge1$, $N\ge2$ let $\Omega=\{1,\dots,N\}^d$, $n=N^d$, and
\begin{equation}\label{eq:boxL}
(Lf)(x)=\frac1{2d}\sum_{y\sim x}\big(f(x)-f(y)\big),
\end{equation}
the sum being over lattice neighbours $y\in\Omega$ of $x$. Then $P=I-qL$ is exactly the lazy walk of the article: each of the $2d$ directions is chosen with probability $q/(2d)$ and a move leaving $\Omega$ is cancelled. $L$ satisfies \eqref{eq:A}, and $L_{xx}\le1$, so $P$ is stochastic for every $q\in(0,1]$.

\begin{lemma}[spectrum of the reflecting box]\label{lem:boxspec}
For $k\in\{0,\dots,N-1\}^d$ put $\theta_j=\pi j/N$,
\begin{equation}\label{eq:boxeig}
\phi_k(x)=\prod_{i=1}^d \gamma_{k_i}\cos\!\big(\theta_{k_i}(x_i-\tfrac12)\big),\quad \gamma_0=N^{-1/2},\ \gamma_j=(2/N)^{1/2}\ (j\ge1),\qquad
\lambda_k=\frac1d\sum_{i=1}^d\big(1-\cos\theta_{k_i}\big).
\end{equation}
Then $\{\phi_k\}$ is an orthonormal basis of $\R^\Omega$ and $L\phi_k=\lambda_k\phi_k$.
\end{lemma}
\begin{proof}
For $d=1$: $g(x)=\cos(\theta_j(x-\frac12))$ satisfies $g(x-1)+g(x+1)=2\cos\theta_j\,g(x)$ for all real $x$, and $g(0)=g(1)$, $g(N+1)=g(N)$ because $\cos(\theta_j(N+\frac12))=\cos(\pi j+\theta_j/2)=\cos(\pi j-\theta_j/2)$. Hence $(Lg)(x)=\frac12(2g(x)-g(x-1)-g(x+1))=(1-\cos\theta_j)g(x)$ for $1\le x\le N$, the boundary rows being covered by the two reflection identities. The $N$ eigenvalues $1-\cos\theta_j$ are distinct, so the eigenvectors are orthogonal, and $\sum_{x=1}^N\cos^2(\theta_j(x-\frac12))=N/2$ for $j\ge1$ (and $N$ for $j=0$) gives the normalisation. For general $d$, $L=\frac1d\sum_i L^{(i)}$ with $L^{(i)}$ the one-dimensional operator acting on coordinate $i$, and the tensor products of the one-dimensional eigenvectors form an orthonormal eigenbasis.
\end{proof}

The main geometry (``corner to corner'', CC) is $a=(N,\dots,N)$, $x_0=(1,\dots,1)$. For the corner target
\begin{equation}\label{eq:cornerweights}
n\,\phi_k(a)^2=\prod_{i=1}^d\varepsilon_{k_i}\cos^2\frac{\theta_{k_i}}2,\qquad
n\,\phi_k(a)\phi_k(x_0)=(-1)^{|k|}\,n\,\phi_k(a)^2,\qquad \varepsilon_0=1,\ \varepsilon_j=2\ (j\ge1),
\end{equation}
with $|k|=k_1+\dots+k_d$, because $\cos(\theta_j(N-\frac12))=(-1)^j\cos(\theta_j/2)$. In particular $\phi_k(a)\neq0$ for every $k$.

\section{Structure of the killed generator}\label{sec:structure}

Fix a target $a$. Put $w_i=\ip{e_a}{E_ie_a}=\|E_ie_a\|^2\ge0$, so that $\sum_{i=0}^M w_i=1$ and $w_0=1/n$.

\begin{definition}[visible spectrum]\label{def:visible}
An eigenvalue $\lambda^{(i)}$ of $L$ is \emph{visible from $a$} if $w_i>0$, i.e.\ if some eigenfunction of $L$ with that eigenvalue does not vanish at $a$. We list the visible eigenvalues as
$0=\Lambda_0<\Lambda_1<\dots<\Lambda_p$ with weights $w_0=\frac1n,w_1,\dots,w_p>0$ (re-indexed), and put
\[
u_i=\frac{E_ie_a}{\sqrt{w_i}}\quad(0\le i\le p),\qquad K=\spn\{u_0,\dots,u_p\},\qquad K_0=K\cap e_a^{\perp}.
\]
The \emph{secular function} and the \emph{cross function} of a site $x$ are
\begin{equation}\label{eq:G}
G(s)=\ip{e_a}{(L-s)^{-1}e_a}=\sum_{i=0}^p\frac{w_i}{\Lambda_i-s},\qquad
G_x(s)=\ip{e_x}{(L-s)^{-1}e_a}=\sum_{i=0}^p\frac{\sqrt{w_i}\,u_i(x)}{\Lambda_i-s}.
\end{equation}
\end{definition}

$K$ is the cyclic subspace generated by $e_a$ under $L$: the $u_i$ are orthonormal eigenvectors of $L$ and $e_a=\sum_i\sqrt{w_i}\,u_i$. Note $u_0=n^{-1/2}\one$, so $\one\in K$. We identify $\R^{\Omega'}$ with $e_a^\perp=\{v\in\R^\Omega:v(a)=0\}$; under this identification $L_a=\Pi L\Pi|_{e_a^\perp}$ with $\Pi=I-e_ae_a^{\mathsf T}$.

\begin{theorem}[structure of $L_a$]\label{thm:structure}
Assume \eqref{eq:A}. Then:
\begin{enumerate}
\item[(i)] $G$ has exactly $p$ zeros in $\mathbb C$. They are real and simple, and they strictly interlace the visible eigenvalues:
\begin{equation}\label{eq:strict}
0=\Lambda_0<\sigma_0<\Lambda_1<\sigma_1<\Lambda_2<\dots<\sigma_{p-1}<\Lambda_p .
\end{equation}
\item[(ii)] $e_a^\perp=K_0\oplus K^\perp$ orthogonally, $\dim K_0=p$, and both summands are invariant under $L_a$. On $K^\perp$ (the eigenfunctions of $L$ that are orthogonal to $K$; all of them vanish at $a$) $L_a$ coincides with $L$. On $K_0$, $L_a$ has the simple eigenvalues $\sigma_0,\dots,\sigma_{p-1}$ with eigenvectors
\begin{equation}\label{eq:psi}
\psi_j=\sum_{i=0}^p\frac{\sqrt{w_i}}{\Lambda_i-\sigma_j}\,u_i,\qquad \psi_j(x)=G_x(\sigma_j),\qquad \|\psi_j\|^2=G'(\sigma_j),\qquad (L-\sigma_j)\psi_j=e_a .
\end{equation}
\item[(iii)] As a multiset, $\spec(L_a)=\{\sigma_0,\dots,\sigma_{p-1}\}\uplus\{\lambda^{(i)}$ with multiplicity $m_i-1$ if $\lambda^{(i)}$ is visible, $m_i$ if it is not$\}$. Equivalently
\begin{equation}\label{eq:det}
\det(L_a-s)=G(s)\,\det(L-s)\qquad(s\in\mathbb C\setminus\spec L).
\end{equation}
\end{enumerate}
\end{theorem}

\begin{proof}
(i) $\mathcal N(s):=G(s)\prod_{l=0}^p(\Lambda_l-s)=\sum_{i}w_i\prod_{l\ne i}(\Lambda_l-s)$ is a polynomial of degree $p$ (its leading coefficient is $(-1)^p\sum_iw_i=(-1)^p$), so $G$ has at most $p$ zeros in $\mathbb C$. On each interval $(\Lambda_j,\Lambda_{j+1})$, $0\le j\le p-1$, $G$ is continuous with $G'(s)=\sum_iw_i(\Lambda_i-s)^{-2}>0$, $G(s)\to-\infty$ as $s\downarrow\Lambda_j$ and $G(s)\to+\infty$ as $s\uparrow\Lambda_{j+1}$. Hence each of these $p$ intervals contains exactly one zero $\sigma_j$, it is simple, and there are no others.

(ii) $K$ is spanned by eigenvectors of $L$, hence $L$-invariant, and so is $K^\perp$ because $L$ is symmetric. As $e_a\in K$, every $v\in K^\perp$ satisfies $v(a)=\ip v{e_a}=0$, i.e.\ $K^\perp\subset e_a^\perp$. If $v\in e_a^\perp$ and $v=v_K+v_\perp$ with $v_K\in K$, $v_\perp\in K^\perp$, then $v_K(a)=v(a)-v_\perp(a)=0$; this proves the orthogonal decomposition, and $\dim K_0=\dim K-1=p$ since $e_a\in K$ is not orthogonal to $K$. For $v\in K^\perp$: $L_av=\Pi Lv=Lv$ because $Lv\in K^\perp\subset e_a^\perp$. For $v\in K_0$: $Lv\in K$ and $\Pi Lv=Lv-(Lv)(a)\,e_a\in K\cap e_a^\perp=K_0$. So both spaces are $L_a$-invariant.

Now let $\psi_j$ be defined by the first formula in \eqref{eq:psi}. Then $\psi_j\in K$, $\psi_j\ne0$, $(L-\sigma_j)\psi_j=\sum_i\sqrt{w_i}u_i=e_a$, and $\psi_j(a)=\ip{e_a}{\psi_j}=\sum_iw_i/(\Lambda_i-\sigma_j)=G(\sigma_j)=0$, so $\psi_j\in K_0$ and $L_a\psi_j=\Pi(\sigma_j\psi_j+e_a)=\sigma_j\psi_j$. The $p$ numbers $\sigma_j$ are distinct and $\dim K_0=p$, so these are all the eigenvalues of $L_a|_{K_0}$ and they are simple. Finally $\psi_j(x)=\ip{e_x}{\psi_j}=G_x(\sigma_j)$ and $\|\psi_j\|^2=\sum_iw_i(\Lambda_i-\sigma_j)^{-2}=G'(\sigma_j)$ by orthonormality of the $u_i$.

(iii) $L|_{K^\perp}$ has the eigenvalue $\lambda^{(i)}$ with multiplicity $\dim(E_i\R^\Omega\cap K^\perp)=m_i-\dim(E_i\R^\Omega\cap K)$, and $E_i\R^\Omega\cap K$ is spanned by $u_i$ if $\lambda^{(i)}$ is visible and is $\{0\}$ otherwise. Together with (ii) this gives the multiset statement. Both sides of \eqref{eq:det} are then the same polynomial divided by $\prod_i(\lambda^{(i)}-s)^{m_i}$: indeed $\det(L_a-s)=\prod_j(\sigma_j-s)\prod_{i\,\mathrm{vis}}(\lambda^{(i)}-s)^{m_i-1}\prod_{i\,\mathrm{invis}}(\lambda^{(i)}-s)^{m_i}$ and $\mathcal N(s)=\prod_j(\sigma_j-s)$ (same zeros, same leading coefficient $(-1)^p$).
(\eqref{eq:det} is also Cramer's rule for the $(a,a)$ entry of $(L-s)^{-1}$.)
\end{proof}

\begin{remark}[degeneracies]\label{rem:degenerate}
Nothing in Theorem~\ref{thm:structure} requires simple eigenvalues. Two kinds of coincidence can occur and are harmless.
(a) A visible eigenvalue of multiplicity $m_i\ge2$ remains an eigenvalue of $L_a$ with multiplicity $m_i-1$; its eigenfunctions vanish at $a$.
(b) A zero $\sigma_j$ of $G$ may coincide with an \emph{invisible} eigenvalue $\lambda^{(i)}$ of $L$; then $\sigma_j$ is an eigenvalue of $L_a$ of multiplicity $m_i+1$, but the eigenvector $\psi_j\in K_0$ is still canonically defined and orthogonal to the other $m_i$.
In both cases the extra eigenfunctions lie in $K^\perp$ and are invisible to the first-passage law (Theorem~\ref{thm:fpt}).
\end{remark}

\section{Spectral representation of the first-passage law}\label{sec:fpt}

\begin{theorem}[first-passage representation]\label{thm:fpt}
Assume \eqref{eq:A}, let $a$ be the target and $x\in\Omega'$ the start. Define the \emph{residues}
\begin{equation}\label{eq:res}
r_j(x)=\frac{\psi_j(x)\,\ip{\psi_j}{\one}}{\|\psi_j\|^2}=-\frac{G_x(\sigma_j)}{\sigma_j\,G'(\sigma_j)},\qquad j=0,\dots,p-1 .
\end{equation}
Then, for the walk with activity $q$,
\begin{align}
&\text{continuous time:}&&\PP_x(T>t)=\sum_{j=0}^{p-1}r_j(x)\,e^{-q\sigma_jt},&&f(t)=\sum_{j=0}^{p-1}r_j(x)\,q\sigma_j\,e^{-q\sigma_jt}\quad(t\ge0);\label{eq:cont}\\
&\text{discrete time:}&&\PP_x(T>t)=\sum_{j=0}^{p-1}r_j(x)\,\nu_j^{\,t},&&f(t)=\sum_{j=0}^{p-1}r_j(x)\,(1-\nu_j)\,\nu_j^{\,t-1}\quad(t\ge1),\label{eq:disc}
\end{align}
with $\nu_j=1-q\sigma_j$ (and the convention $0^0=1$). Moreover, for every $s\in\mathbb C$ that is not a zero of $G$,
\begin{equation}\label{eq:renewal}
\sum_{j=0}^{p-1}\frac{r_j(x)\,\sigma_j}{\sigma_j-s}=\frac{G_x(s)}{G(s)} ,
\end{equation}
so that $\E_x e^{-uT}=G_x(-u/q)/G(-u/q)$ in continuous time and $\E_x z^{T}=G_x(\zeta)/G(\zeta)$ with $\zeta=-(1-z)/(qz)$ in discrete time.
Only the zeros of $G$ occur: eigenvalues of $L_a$ belonging to $K^\perp$ carry no weight, and $\sigma_j$ is absent from the law started at $x$ if and only if $G_x(\sigma_j)=0$.
\end{theorem}

\begin{proof}
The vector $\one_{\Omega'}=\one-e_a$ lies in $K$ (both terms do) and vanishes at $a$, so $\one_{\Omega'}\in K_0$. By Theorem~\ref{thm:structure}(ii) the $\psi_j$ form an orthogonal basis of $K_0$, hence $\one_{\Omega'}=\sum_j\frac{\ip{\psi_j}{\one_{\Omega'}}}{\|\psi_j\|^2}\psi_j$. Since $\psi_j(a)=0$ and $\one=\sqrt n\,u_0$,
\[
\ip{\psi_j}{\one_{\Omega'}}=\ip{\psi_j}{\one}=\sqrt n\cdot\frac{\sqrt{w_0}}{\Lambda_0-\sigma_j}=-\frac1{\sigma_j},
\]
which with \eqref{eq:psi} gives the second expression in \eqref{eq:res}. Applying $e^{-qtL_a}$, respectively $(I-qL_a)^t$, to this expansion and using \eqref{eq:surv} gives the survival formulae; $f=-\frac{d}{dt}\PP_x(T>t)$, respectively $f(t)=\PP_x(T>t-1)-\PP_x(T>t)$, gives the densities.

For \eqref{eq:renewal}: for $x\ne a$, $G_x(s)=-\frac1s\ip{e_x}{(1-L/s)^{-1}e_a}=O(s^{-2})$ as $s\to\infty$, while $G(s)=-1/s+O(s^{-2})$, so the rational function $G_x/G$ vanishes at infinity; its poles are the simple zeros $\sigma_j$ of $G$ (the poles $\Lambda_i$ of $G_x$ are poles of $G$ of the same order and cancel), with residues $G_x(\sigma_j)/G'(\sigma_j)=-r_j\sigma_j$. Hence $G_x(s)/G(s)=\sum_j(-r_j\sigma_j)/(s-\sigma_j)$. Finally $\E_xe^{-uT}=\int_0^\infty e^{-ut}f(t)\,dt=\sum_jr_jq\sigma_j/(u+q\sigma_j)$, and $\E_xz^T=\sum_jr_j(1-\nu_j)z/(1-\nu_jz)=\sum_j r_j\sigma_j/(\sigma_j-\zeta)$.
\end{proof}

\begin{remark}[relation to the renewal ratio of the article]
With
\[\widehat P(y,u|x)=\int_0^\infty e^{-ut}p_t(x,y)\,dt=\ip{e_x}{(u+qL)^{-1}e_y}\]
for the walk without target, $G_x(-u/q)=q\,\widehat P(a,u|x)$, so \eqref{eq:renewal} is the renewal identity $\widehat F(u)=\widehat P(a,u|x_0)/\widehat P(a,u|a)$, and Theorem~\ref{thm:fpt} proves that its poles are exactly $u=-q\sigma_j$ with $\sigma_j$ the eigenvalues of the killed generator in the sector $K_0$. The $a_j,\nu_j$ of the computations of Sections~5--6 of the article are $a_j=r_j\sigma_j$, $\nu_j=\sigma_j$.
\end{remark}

\section{Interlacing}\label{sec:interlacing}

\begin{theorem}[interlacing]\label{thm:interlace}
Assume \eqref{eq:A}. Let $\lambda_1\le\dots\le\lambda_n$ and $\mu_1\le\dots\le\mu_{n-1}$ be the eigenvalues of $L$ and of $L_a$, counted with multiplicity.
\begin{enumerate}
\item[(a)] (Cauchy) $\lambda_k\le\mu_k\le\lambda_{k+1}$ for $1\le k\le n-1$.
\item[(b)] (strict interlacing in the visible sector) The eigenvalues of $L_a$ that carry first-passage weight are the zeros of $G$, and they satisfy \eqref{eq:strict}: between two consecutive \emph{visible} eigenvalues of $L$ there is exactly one, and it is different from both.
\item[(c)] (equality cases) An eigenvalue $\mu$ of $L_a$ is also an eigenvalue of $L$ if and only if $L$ has an eigenfunction with eigenvalue $\mu$ vanishing at $a$; the multiplicity of $\mu$ in $\spec(L_a)$ is then given by Theorem~\ref{thm:structure}(iii) and Remark~\ref{rem:degenerate}.
\item[(d)] (symmetry sector) Let $\Gamma_a$ be the group of permutations $g$ of $\Omega$ with $gLg^{-1}=L$ and $ga=a$, and $\mathcal F=\{f: f\circ g=f\ \forall g\in\Gamma_a\}$. Then $K\subset\mathcal F$. Hence the $\sigma_j$ are eigenvalues of the killed walk restricted to the $\Gamma_a$-symmetric sector, and they strictly interlace the visible distinct eigenvalues of the reflecting walk restricted to that sector.
\end{enumerate}
\end{theorem}
\begin{proof}
(a) is the Cauchy interlacing theorem for a principal submatrix of a real symmetric matrix (R.~A.~Horn and C.~R.~Johnson, \emph{Matrix Analysis}, Cambridge Univ.\ Press: Theorem~4.3.17 in the 2nd edition, 2013; Theorem~4.3.8 in the 1st edition, 1985). It also follows from the Courant--Fischer min--max principle, since $L_a$ is the compression of $L$ to the hyperplane $e_a^\perp$. (b) is Theorem~\ref{thm:structure}(i) together with Theorem~\ref{thm:fpt}. (c): by Theorem~\ref{thm:structure}(ii)--(iii), $\spec(L_a)=\{\sigma_j\}\uplus\spec(L|_{K^\perp})$ and no $\sigma_j$ is a visible eigenvalue of $L$; so $\mu\in\spec(L_a)\cap\spec(L)$ iff $\mu\in\spec(L|_{K^\perp})$ (using Remark~\ref{rem:degenerate}(b) when $\mu=\sigma_j$ is an invisible eigenvalue), iff $L$ has an eigenfunction for $\mu$ in $K^\perp$; and an eigenfunction $v$ of $L$ with eigenvalue $\lambda^{(i)}$ lies in $K^\perp$ iff $0=\ip v{u_i}\sqrt{w_i}=\ip v{E_ie_a}=\ip{v}{e_a}=v(a)$. (d) $ge_a=e_a$ and $g$ commutes with every $E_i$, so $gE_ie_a=E_ie_a$.
\end{proof}

\begin{corollary}[the reflecting box]\label{cor:box}
Let $L$ be \eqref{eq:boxL} and $a\in\Omega$ arbitrary.
\begin{enumerate}
\item[(a)] An eigenvalue $\lambda$ of $L$ is visible from $a$ iff there is $k$ with $\lambda_k=\lambda$ and $\prod_i\cos(\theta_{k_i}(a_i-\frac12))\ne0$, i.e.\ iff for no $i$ the number $k_i(2a_i-1)/N$ is an odd integer.
\item[(b)] For a corner target every eigenvalue of $L$ is visible. Hence the number of poles of the first-passage law is $p=\#\{\lambda_k:k\in\{0,\dots,N-1\}^d\}-1$ (number of distinct values), the poles satisfy
$\Lambda_j<\sigma_j<\Lambda_{j+1}$ for the increasing list $\Lambda_0<\dots<\Lambda_p$ of \emph{all} distinct eigenvalues of $L$, and $\Gamma_a\supseteq S_d$ (coordinate permutations), so all $\psi_j$ are symmetric functions of the coordinates. A distinct eigenvalue of multiplicity $m$ (permutation images and accidental degeneracies alike) stays in $\spec(L_a)$ with multiplicity $m-1$ and without first-passage weight.
\item[(c)] For $d\ge 2$, $N\ge2$ and the corner target: $\Lambda_1=\frac1d(1-\cos\frac\pi N)$ with multiplicity $d$ (the modes $k=e_i$), $\Lambda_2=2\Lambda_1$ (the modes $k=e_i+e_j$, $i\ne j$), and
\begin{equation}\label{eq:W}
W:=n\,w_1=2d\cos^2\frac{\pi}{2N},\qquad
0<\sigma_0<\Lambda_1<\sigma_1<2\Lambda_1<\sigma_j\ \ (j\ge2) .
\end{equation}
\end{enumerate}
\end{corollary}
\begin{proof}
(a),(b): $w(\lambda)=\sum_{k:\lambda_k=\lambda}\phi_k(a)^2$ and Lemma~\ref{lem:boxspec}; for the corner use \eqref{eq:cornerweights}.
(c) The smallest non-zero value of $\lambda_k$ is attained exactly at $k=e_i$, $1\le i\le d$. Any other $k\ne0$ has either at least two non-zero components, giving $\lambda_k\ge\frac2d(1-\cos\frac\pi N)=2\Lambda_1$ with equality iff $k=e_i+e_j$ ($i\ne j$), or exactly one non-zero component $k_i\ge2$ (which requires $N\ge3$), giving
$\lambda_k\ge\frac1d(1-\cos\frac{2\pi}N)=\frac8d\sin^2\frac\pi{2N}\cos^2\frac\pi{2N}=2\Lambda_1\cdot2\cos^2\frac{\pi}{2N}>2\Lambda_1$ because $2\cos^2\frac\pi{2N}>1$ for $N\ge3$. So $\Lambda_2=2\Lambda_1$, and its eigenspace is spanned by the $\phi_{e_i+e_j}$. The weight is $n\sum_i\phi_{e_i}(a)^2=d\cdot2\cos^2\frac\pi{2N}$ by \eqref{eq:cornerweights}.
\end{proof}

\section{Residues: normalisation, moments, signs}\label{sec:residues}

Throughout this section $x\in\Omega'$ is the start and $r_j=r_j(x)$.

\begin{theorem}[residues]\label{thm:residues}
Assume \eqref{eq:A}.
\begin{enumerate}
\item[(a)] $\sum_{j=0}^{p-1}r_j=1$.
\item[(b)] In continuous time $\E_x[T^\ell]=\ell!\,q^{-\ell}\sum_jr_j\sigma_j^{-\ell}$ for $\ell\ge1$; in both time conventions $\E_xT=\frac1q\sum_jr_j/\sigma_j$.
\item[(c)] If the graph of $L_a$ on $\Omega'$ is connected, then $\sigma_0$ is the smallest eigenvalue of $L_a$, it is a simple eigenvalue of $L_a$, $\psi_0$ has a strict sign on $\Omega'$, and
\[
r_0(y)=\frac{\psi_0(y)\sum_{z}\psi_0(z)}{\sum_z\psi_0(z)^2}>0\qquad\text{for every }y\in\Omega' .
\]
\item[(d)] Let $D=\dist(x,a)\ge1$ be the graph distance in the graph of $L$. Then
\begin{equation}\label{eq:vanish}
\sum_jr_j\sigma_j^{\,i}=0\quad(1\le i\le D-1),\qquad (-1)^{D-1}\sum_jr_j\sigma_j^{\,D}>0 .
\end{equation}
Consequently the sequence $(r_0,r_1,\dots,r_{p-1})$ (zeros deleted) has at least $D-1$ sign changes; in particular at least $D$ residues are non-zero, $p\ge D$, and under the hypothesis of (c) at least $\lceil (D-1)/2\rceil$ residues are negative.
\end{enumerate}
\end{theorem}

\begin{proof}
(a) is \eqref{eq:cont} at $t=0$. (b) follows by integrating $t^{\ell-1}\PP_x(T>t)$ over $[0,\infty)$, respectively summing $\PP_x(T>t)=\sum_jr_j\nu_j^t$ over $t\ge0$, which gives $\sum_jr_j/(1-\nu_j)$ (here $|\nu_j|<1$ since $0<q\sigma_j<2$: $\spec L\subset[0,2\max_xL_{xx}]$ by Gershgorin and $q\max_x L_{xx}\le1$, and $\sigma_j<\Lambda_p$).

(c) If $|\Omega'|=1$ (i.e.\ $n=2$, $\Omega'=\{y\}$), then $L_a$ is the $1\times1$ matrix $(L_{yy})$ with $L_{yy}=-L_{ya}>0$, $p=\dim K_0=1$, so $\sigma_0=L_{yy}$ is the only eigenvalue of $L_a$, $\psi_0\ne0$ is a non-zero multiple of $e_y$, and $r_0(y)=1$ by (a); all claims hold trivially. Let now $|\Omega'|\ge2$ and $c=\max_xL_{xx}$. The matrix $A=cI-L_a$ is real symmetric, entrywise non-negative and of order $\ge2$, and it is irreducible because its off-diagonal zero pattern is that of $L_a$, whose graph is connected. By the Perron--Frobenius theorem (Horn--Johnson, Theorem~8.4.4: for an irreducible non-negative matrix of order at least $2$ the spectral radius $\rho(A)$ is positive, is an algebraically simple eigenvalue, and has an entrywise positive eigenvector $\psi$) and because the eigenvalues of the symmetric matrix $A$ are real and bounded in modulus by $\rho(A)$, the largest eigenvalue of $A$ is $\rho(A)$ and it is simple; i.e.\ the smallest eigenvalue $\mu_1=c-\rho(A)$ of $L_a$ is simple and $\psi>0$ on $\Omega'$. As $\mu_1$ is simple and $K_0$, $K^\perp$ are $L_a$-invariant, $\psi$ belongs to one of them; but $\ip\psi{\one_{\Omega'}}>0$ and $\one_{\Omega'}\in K_0$, so $\psi\notin K^\perp$. Hence $\psi\in K_0$ is a multiple of some $\psi_j$ and $\mu_1=\sigma_j$; being the minimum, $\mu_1=\sigma_0$. The formula for $r_0$ is \eqref{eq:res}.

(d) Let $\beta=L_a\one_{\Omega'}$. Since the rows of $L$ sum to zero, $\beta(y)=-L_{ya}\ge0$ for $y\in\Omega'$, and $\beta$ is supported on the neighbours of $a$. By the expansion in the proof of Theorem~\ref{thm:fpt}, $\sum_jr_j\sigma_j^{\,i}=(L_a^{\,i}\one_{\Omega'})(x)=(L_a^{\,i-1}\beta)(x)=\sum_y(L_a^{\,i-1})_{xy}\beta(y)$ for $i\ge1$. The entry $(L_a^{\,i-1})_{xy}$ is a sum over sequences $x=z_0,z_1,\dots,z_{i-1}=y$ in $\Omega'$, with consecutive points equal or adjacent, of $\prod_lL_{z_lz_{l+1}}$; it vanishes if $\dist(x,y)>i-1$. Every neighbour $y$ of $a$ has $\dist(x,y)\ge D-1$, so the sum is zero for $i\le D-1$. For $i=D$ only geodesics of length $D-1$ from $x$ to neighbours $y$ of $a$ contribute; each contributes a product of $D-1$ negative numbers times $\beta(y)>0$, and at least one exists (cut the last step off a geodesic from $x$ to $a$; its vertices other than $a$ lie in $\Omega'$). This proves \eqref{eq:vanish}.

Suppose the non-zero residues had only $S\le D-2$ sign changes. Choose $c_1<\dots<c_S$, none equal to a $\sigma_j$, separating consecutive maximal blocks of equal sign, and put $P(\sigma)=\sigma\prod_{l=1}^S(c_l-\sigma)$. Then $r_jP(\sigma_j)$ has the same sign for all $j$ with $r_j\neq0$ (recall $\sigma_j>0$), and it is not identically zero because $\sum_jr_j=1$. Thus $\sum_jr_jP(\sigma_j)\ne0$. But $P$ is a linear combination of $\sigma,\dots,\sigma^{S+1}$ with $S+1\le D-1$, so $\sum_jr_jP(\sigma_j)=0$ by \eqref{eq:vanish}; contradiction. The remaining assertions are immediate ($S$ sign changes need $S+1$ non-zero terms; with $r_0>0$ and at least $D-1$ sign changes, at least $\lceil(D-1)/2\rceil$ terms are negative).
\end{proof}

\begin{corollary}[one dimension: complete alternation]\label{cor:1d}
For $d=1$, target $a=N$ and start $x=1$: $p=N-1$, all residues are non-zero and $\operatorname{sign}r_j=(-1)^j$.
\end{corollary}
\begin{proof}
The $N$ eigenvalues $1-\cos\theta_j$ are distinct and visible from $a=N$ by \eqref{eq:cornerweights}, so $p=N-1$; $D=N-1$, so the $p$ residues have at least $p-1$ sign changes, which forces all to be non-zero and alternating; $r_0>0$ by Theorem~\ref{thm:residues}(c).
\end{proof}

\begin{corollary}[corner to corner, $d\ge2$]\label{cor:ccsigns}
For the box with $a=(N,\dots,N)$, $x_0=(1,\dots,1)$: $D=d(N-1)$, the continuous-time density satisfies $f(t)=c\,t^{D-1}(1+o(1))$ as $t\downarrow0$ with $c>0$, and the residue sequence has at least $d(N-1)-1$ sign changes.
\end{corollary}

\begin{proposition}[stationary start; sizes of the residues]\label{prop:uniform}
Assume \eqref{eq:A} and put
\begin{equation}\label{eq:varrho}
\varrho_j=\frac1n\sum_{x\in\Omega'}r_j(x)=\frac{1}{n\,\sigma_j^2\,G'(\sigma_j)}\qquad(0\le j\le p-1).
\end{equation}
\begin{enumerate}
\item[(a)] $\varrho_j>0$ for every $j$ and $\sum_j\varrho_j=1-\frac1n$. If the start is uniformly distributed on $\Omega$ (the stationary law), then $\PP(T>t)=\sum_j\varrho_j\,e^{-q\sigma_jt}$ in continuous time: the first-passage law is a mixture of exponentials (plus an atom $\frac1n$ at $0$), its density is completely monotone, and it has no interior mode. In discrete time $\PP(T>t)=\sum_j\varrho_j\nu_j^{\,t}$.
\item[(b)] For every start $x\in\Omega'$,
\[
\sum_{j=0}^{p-1}\frac{r_j(x)^2}{n\,\varrho_j}\le1,\qquad\text{hence}\qquad|r_j(x)|\le\sqrt{n\,\varrho_j}\quad\text{and}\quad\sum_{j}|r_j(x)|\le\sqrt{n-1}.
\]
\item[(c)] (factorisation) For $s$ not a zero of $G$ and $s\ne0$: $\dfrac{G_x(s)}{G(s)}=\big[-ns\,G_x(s)\big]\cdot\Big[\dfrac{-1}{ns\,G(s)}\Big]$, where $-1/(nsG(s))=\frac1n+\sum_j\varrho_j\sigma_j/(\sigma_j-s)$ is the transform of the stationary-start law and $-nsG_x(-u/q)=u\int_0^\infty e^{-ut}\,n\,p_t(x,a)\,dt$ is the transform of the derivative of the normalised arrival profile $n\,p_t(x,a)$ of the walk without target.
\end{enumerate}
\end{proposition}
\begin{proof}
(a),(b): By \eqref{eq:res} and \eqref{eq:psi}, $\sum_{x}r_j(x)=\ip{\psi_j}{\one}^2/\|\psi_j\|^2=1/(\sigma_j^2G'(\sigma_j))>0$, and $\sum_j\ip{\psi_j}{\one}^2/\|\psi_j\|^2=\|\one_{\Omega'}\|^2=n-1$ because the $\psi_j$ are an orthogonal basis of $K_0\ni\one_{\Omega'}$. Averaging \eqref{eq:cont}, \eqref{eq:disc} over $x\in\Omega$ (with $T=0$ for $x=a$) gives the mixture representation. Next, $r_j(x)^2=\pi_j(x)\,\ip{\psi_j}\one^2/\|\psi_j\|^2=\pi_j(x)\,n\varrho_j$ with $\pi_j(x):=\psi_j(x)^2/\|\psi_j\|^2$, and $\sum_j\pi_j(x)$ is the squared norm of the orthogonal projection of $e_x$ onto $K_0$, hence $\le1$. The last inequality is Cauchy--Schwarz: $\sum_j|r_j|\le(\sum_j\pi_j(x))^{1/2}(\sum_jn\varrho_j)^{1/2}$.
(c) $\frac1n\sum_{x\in\Omega}G_x(s)=\frac1n\ip{\one}{(L-s)^{-1}e_a}=-\frac1{ns}$, so the average over $x\in\Omega$ of $G_x/G$ (which is $1$ for $x=a$) equals $-1/(nsG(s))$; by \eqref{eq:renewal} and \eqref{eq:varrho} it also equals $\frac1n+\sum_j\varrho_j\sigma_j/(\sigma_j-s)$. Finally $G_x(-u/q)=q\int_0^\infty e^{-ut}p_t(x,a)\,dt$ for the walk with generator $-qL$.
\end{proof}

\begin{proposition}[quasi-stationary averages]\label{prop:qsd}
Assume \eqref{eq:A} and that the graph of $L_a$ on $\Omega'$ is connected, and let $\nu(x)=\psi_0(x)/\sum_z\psi_0(z)$, $x\in\Omega'$, be the normalised principal eigenvector of $L_a$ (the quasi-stationary distribution). Then
\begin{equation}\label{eq:qsd}
\sum_{x\in\Omega'}\nu(x)\,r_j(x)=\delta_{j0}\quad(0\le j\le p-1),\qquad \sum_{x\in\Omega'}\nu(x)\,m(x)=\frac1{\sigma_0},
\end{equation}
where $m(x)=\E_xT$ for the unit-rate walk. Equivalently, started from $\nu$ the hitting time is exactly exponential: $\PP_\nu(T>t)=e^{-q\sigma_0t}$ (continuous time), $\PP_\nu(T>t)=\nu_0^{\,t}$ (discrete time). Consequently
\[
\min_{x\in\Omega'}r_0(x)\le1\le\max_{x\in\Omega'}r_0(x),\qquad \min_{x\in\Omega'}m(x)\le\frac1{\sigma_0}\le\max_{x\in\Omega'}m(x),
\]
and each of these four inequalities is an equality if and only if $\one_{\Omega'}$ is an eigenvector of $L_a$, i.e.\ iff $L_{ya}$ takes the same value for all $y\in\Omega'$.
\end{proposition}
\begin{proof}
By \eqref{eq:res} and the orthogonality of the $\psi_j$, $\sum_x\psi_0(x)r_j(x)=\ip{\psi_0}{\psi_j}\ip{\psi_j}{\one}/\|\psi_j\|^2=\delta_{j0}\ip{\psi_0}{\one}$; divide by $\ip{\psi_0}\one>0$ (Theorem~\ref{thm:residues}(c)). Inserting this into \eqref{eq:cont}, \eqref{eq:disc} gives the exponential (geometric) law, and with Theorem~\ref{thm:residues}(b) the mean $\sum_x\nu(x)m(x)=\sum_jr_j^{\nu}/\sigma_j=1/\sigma_0$. The two-sided inequalities hold because $\nu$ is a probability vector with strictly positive entries, and equality in any of them forces $r_0(\cdot)$, respectively $m(\cdot)$, to be constant on $\Omega'$. By \eqref{eq:res}, $r_0(\cdot)$ is constant iff $\psi_0$ is, i.e.\ iff $\one_{\Omega'}$ is an eigenvector of $L_a$; and $m=L_a^{-1}\one_{\Omega'}$ (Lemma~\ref{lem:taum}) is constant iff $L_a\one_{\Omega'}$ is a multiple of $\one_{\Omega'}$. Finally $(L_a\one_{\Omega'})(y)=-L_{ya}$ because the rows of $L$ sum to zero.
\end{proof}

\subsection{Reflection-symmetric pairs: the sign of $r_1$}

\begin{theorem}[sign rule for a symmetric start--target pair]\label{thm:r1}
Assume \eqref{eq:A} and let $R$ be a permutation of $\Omega$ with $RLR^{-1}=L$, $R^2=\mathrm{id}$ and $Ra=x_0\neq a$. Put $v_\pm=\frac12(e_a\pm e_{x_0})$ and
\[
G^{\pm}(s)=\ip{v_\pm}{(L-s)^{-1}v_\pm}=\sum_{i}\frac{w_i^{\pm}}{\lambda^{(i)}-s},\qquad w_i^\pm=\|E_iv_\pm\|^2 .
\]
\begin{enumerate}
\item[(a)] $G=G^++G^-$, $G_{x_0}=G^+-G^-$, and for every $j$
\begin{equation}\label{eq:rsign}
r_j(x_0)=\frac{2\,G^-(\sigma_j)}{\sigma_j\,G'(\sigma_j)}=-\frac{2\,G^+(\sigma_j)}{\sigma_j\,G'(\sigma_j)} .
\end{equation}
Thus $r_j(x_0)>0$ iff $G^-(\sigma_j)>0$ iff $G^+(\sigma_j)<0$.
\item[(b)] Suppose in addition that (H1) the first visible eigenvalue $\Lambda_1$ is \emph{purely odd}, $E_{\Lambda_1}v_+=0$, and (H2) $L$ has at least two eigenvalues, counted with multiplicity, in $(0,\Lambda_1]$. Then $r_1(x_0)<0$.
\end{enumerate}
\end{theorem}
\begin{proof}
(a) $R$ commutes with $(L-s)^{-1}$ and maps $v_\pm$ to $\pm v_\pm$, so $\ip{v_+}{(L-s)^{-1}v_-}=0$. With $e_a=v_++v_-$, $e_{x_0}=v_+-v_-$ this gives the two decompositions, and \eqref{eq:rsign} follows from \eqref{eq:res} and $G^+(\sigma_j)=-G^-(\sigma_j)$.

(b) Let $L_{a,x_0}$ be the principal submatrix of $L$ on $\Omega''=\Omega\setminus\{a,x_0\}$ and $\eta_1>0$ its smallest eigenvalue ($L_a$ is positive definite by Theorem~\ref{thm:structure}(iii), hence so is its principal submatrix $L_{a,x_0}$). Applying Theorem~\ref{thm:interlace}(a) twice ($L\to L_a\to L_{a,x_0}$) gives $\eta_1\le\mu_2\le\lambda_3$, the third smallest eigenvalue of $L$ counted with multiplicity; by (H2), $\lambda_3\le\Lambda_1$. So
\begin{equation}\label{eq:eta}
0<\eta_1\le\Lambda_1 .
\end{equation}
By (H2), $n\ge3$, so $\Omega''\ne\emptyset$. With $c=\max_xL_{xx}$ the matrix $cI-L_{a,x_0}$ is non-negative and symmetric, hence (Horn--Johnson, Theorem~8.3.1, valid for every order $\ge1$: the spectral radius of a non-negative matrix is an eigenvalue with a non-negative eigenvector) has a non-negative eigenvector $\chi\ne0$ for its spectral radius, which for a symmetric matrix is its largest eigenvalue $c-\eta_1$. $R$ preserves $\Omega''$ and commutes with $L_{a,x_0}$, so $\psi=\chi+\chi\circ R$ is a non-negative, non-zero, $R$-invariant eigenvector of $L_{a,x_0}$ for $\eta_1$. Extend $\psi$ by zero to $\Omega$. Then $L\psi=\eta_1\psi+\alpha(e_a+e_{x_0})$ with $\alpha=(L\psi)(a)=\sum_yL_{ay}\psi(y)\le0$ (the coefficients at $a$ and $x_0$ agree by $R$-invariance). If $\alpha=0$, $\psi$ would be an eigenvector of $L$ with eigenvalue $\eta_1>0$, hence orthogonal to $\one$, contradicting $\psi\ge0$, $\psi\ne0$. So $\alpha<0$ and, applying $E_i$,
\[
(\lambda^{(i)}-\eta_1)\,E_i\psi=2\alpha\,E_iv_+\qquad\text{for all }i .
\]
If $\lambda^{(i)}=\eta_1$ this gives $E_iv_+=0$; otherwise $w_i^+/(\lambda^{(i)}-\eta_1)=\ip{E_iv_+}{E_i\psi}/(2\alpha)$. Summing,
\[
G^+(\eta_1)=\frac1{2\alpha}\sum_{i:\lambda^{(i)}\ne\eta_1}\ip{E_iv_+}{\psi}=\frac1{2\alpha}\ip{v_+}{\psi}=\frac{\psi(a)+\psi(x_0)}{4\alpha}=0 .
\]
By (H1) the poles of $G^+$ are $0$ (with $w_0^+=1/n$) and visible eigenvalues $\ge\Lambda_2$; let $\Lambda^+\ge\Lambda_2$ be the smallest positive one (there is one, otherwise $G^+(s)=-1/(ns)$ would have no zero). On $(0,\Lambda^+)$, $G^+$ increases strictly from $-\infty$ to $+\infty$, and by \eqref{eq:eta} its zero there is $\eta_1\le\Lambda_1$. Since $\Lambda_1<\sigma_1<\Lambda_2\le\Lambda^+$ we get $G^+(\sigma_1)>G^+(\eta_1)=0$, and \eqref{eq:rsign} gives $r_1(x_0)<0$.
\end{proof}

\begin{corollary}[corner to corner]\label{cor:r1cc}
For the reflecting box with $d\ge2$, every $N\ge2$, $a=(N,\dots,N)$ and $x_0=(1,\dots,1)$:
\[
r_0>0,\qquad r_1<0 .
\]
\end{corollary}
\begin{proof}
$\Omega'$ is connected, so $r_0>0$ by Theorem~\ref{thm:residues}(c). Take $R:x_i\mapsto N+1-x_i$ for all $i$. Then $\phi_k\circ R=(-1)^{|k|}\phi_k$; the eigenspace of $\Lambda_1$ is spanned by the $\phi_{e_i}$, which are odd, so (H1) holds; it has dimension $d\ge2$ and $\Lambda_1$ is the smallest non-zero eigenvalue, so (H2) holds.
\end{proof}

\begin{remark}[which residues are positive]\label{rem:signrule}
In the situation of Theorem~\ref{thm:r1}(a), let $0=\Lambda^+_0<\Lambda^+_1<\dots$ be the poles of $G^+$ and $\rho^+_i\in(\Lambda^+_i,\Lambda^+_{i+1})$ its zeros. The $\rho_i^+$ are eigenvalues of the walk with \emph{both} $a$ and $x_0$ absorbing in the $R$-even sector (same argument as Theorem~\ref{thm:structure}, applied to $v_+$), and $G^+>0$ exactly on $\bigcup_i(\rho^+_i,\Lambda^+_{i+1})$. Hence
\[
r_j(x_0)<0\iff\sigma_j\in\bigcup_i\big(\rho^+_i,\Lambda^+_{i+1}\big),\qquad r_j(x_0)=0\iff\sigma_j\in\{\rho^+_i\}.
\]
For the box, $G^{+}$ ($G^-$) is the part of \eqref{eq:G} with $|k|$ even (odd), and $\E_{x_0}T=\frac{2n}{q}\,G^{-}(0)$ (Lemma~\ref{lem:taum} below). In one dimension the rule is consistent with the alternation of Corollary~\ref{cor:1d}. In $d\ge2$ with $N\ge3$ the residues do \emph{not} alternate (Theorem~\ref{thm:noalt}, Proposition~\ref{prop:signs}); formula \eqref{eq:rsign} and the rule were checked against brute-force residues for all poles (script \texttt{12}).
\end{remark}

\subsection{Alternation of the residues holds only for $d=1$ or $N=2$}\label{sec:noalt}

In one dimension the residues alternate in sign (Corollary~\ref{cor:1d}). The natural conjecture that they also alternate for the corner-to-corner geometry in $d\ge2$ is \emph{false} (already for $d=2$, $N=3$ the signs are \texttt{+--+-}, Proposition~\ref{prop:signs}). The following two statements are the true replacement: a monotonicity rule across every eigenvalue of pure parity, and, as a consequence, a complete answer for the box.

\begin{proposition}[monotonicity of the signs across an eigenvalue of pure parity]\label{prop:pure}
In the situation of Theorem~\ref{thm:r1}(a) call a visible eigenvalue $\Lambda_i$ \emph{purely even} if $E_{\Lambda_i}v_-=0$ and \emph{purely odd} if $E_{\Lambda_i}v_+=0$ ($E_{\Lambda_i}$ the spectral projector of $L$ for $\Lambda_i$). Let $1\le i\le p-1$ and $r_j=r_j(x_0)$.
\begin{enumerate}
\item[(a)] If $\Lambda_i$ is purely even, then $G^-(\sigma_{i-1})<G^-(\sigma_i)$; hence $r_{i-1}\ge0$ implies $r_i>0$, and $r_i\le0$ implies $r_{i-1}<0$.
\item[(b)] If $\Lambda_i$ is purely odd, then $G^+(\sigma_{i-1})<G^+(\sigma_i)$; hence $r_{i-1}\le0$ implies $r_i<0$, and $r_i\ge0$ implies $r_{i-1}>0$.
\end{enumerate}
In particular the pair $(r_{i-1},r_i)$ cannot have the signs $(+,-)$ in case (a) nor $(-,+)$ in case (b); and if $\Lambda_i,\Lambda_{i+1},\dots,\Lambda_{i+l-1}$ ($i+l-1\le p-1$) are all purely even (all purely odd), then the sequence $\operatorname{sign}r_{i-1},\dots,\operatorname{sign}r_{i+l-1}\in\{-1,0,1\}$ is non-decreasing (non-increasing), with at most one zero.
\end{proposition}
\begin{proof}
As in the proof of Theorem~\ref{thm:r1}(a), $R$ commutes with every spectral projector and $Rv_\pm=\pm v_\pm$, so $\ip{E_iv_+}{E_iv_-}=0$ and $w(\lambda^{(i)})=\|E_ie_a\|^2=w_i^++w_i^-$ for every eigenvalue. Hence every pole of $G^+$ or of $G^-$ is a visible eigenvalue. (a) If $w^-(\Lambda_i)=0$, the only visible eigenvalue in the open interval $I=(\Lambda_{i-1},\Lambda_{i+1})$ is not a pole of $G^-$, so $G^-$ is real-analytic on $I$ with $(G^-)'(s)=\sum_lw_l^-(\lambda^{(l)}-s)^{-2}>0$ (some $w_l^-$ is positive because $\sum_lw_l^-=\|v_-\|^2=\frac12$). By \eqref{eq:strict}, $\sigma_{i-1}<\Lambda_i<\sigma_i$ both lie in $I$, so $G^-(\sigma_{i-1})<G^-(\sigma_i)$, and by \eqref{eq:rsign} $\operatorname{sign}r_j=\operatorname{sign}G^-(\sigma_j)$ because $\sigma_jG'(\sigma_j)>0$. (b) is the same argument for $G^+$, whose poles in $I$ could only be $\Lambda_i$ (note $0=\Lambda_0\notin I$), with $\operatorname{sign}r_j=-\operatorname{sign}G^+(\sigma_j)$. The last statement follows by applying (a) (resp.\ (b)) to consecutive indices.
\end{proof}

\begin{theorem}[the residues alternate if and only if $d=1$ or $N=2$]\label{thm:noalt}
Consider the reflecting box $\{1,\dots,N\}^d$, $d\ge1$, $N\ge2$, with $a=(N,\dots,N)$, $x_0=(1,\dots,1)$, $D=d(N-1)$, and say that the residues \emph{alternate} if $\operatorname{sign}r_j=(-1)^j$ for all $0\le j\le p-1$.
\begin{enumerate}
\item[(a)] If $d=1$ or $N=2$, then $p=D$ and the residues alternate.
\item[(b)] Let $d\ge2$ and $N\ge3$, and put $i=p-3$ if $N\ne4$, $i=p-4$ if $N=4$. Then $i\ge1$, and the two consecutive eigenvalues $\Lambda_i$, $\Lambda_{i+1}$ are both purely even or both purely odd in the sense of Proposition~\ref{prop:pure}: purely even if $D$ is even and $N\ne4$, or if $D$ is odd and $N=4$; purely odd otherwise. Consequently the sequence $(\operatorname{sign}r_{i-1},\operatorname{sign}r_i,\operatorname{sign}r_{i+1})$ is non-decreasing in the first case and non-increasing in the second; in particular the residues do not alternate.
\item[(c)] (failures at small indices) Moreover, writing ``$(r_l,r_{l+1})\ne(+,-)$'' for ``$r_l>0>r_{l+1}$ does not hold'': for $d=2$, $N\ge3$, $(r_2,r_3)\ne(+,-)$; for $d\ge4$, $N\ge4$, $(r_4,r_5)\ne(+,-)$; for $d=3$, $N\ge7$, $(r_6,r_7)\ne(+,-)$. In each of these cases alternation would require exactly the excluded signs.
\end{enumerate}
The proof uses no computer assistance.
\end{theorem}
\begin{proof}
Let $u=\frac{\pi}{2N}$, $x=\sin^2u$ and $\varsigma_j=\dfrac{1-\cos\theta_j}{1-\cos\theta_1}=\dfrac{\sin^2(ju)}{\sin^2u}$ for $0\le j\le N-1$. Then $\varsigma_0=0$, $\varsigma_1=1$, $\varsigma_2=4\cos^2u=4(1-x)=:t$, $\varsigma_3=(3-4x)^2$ (from $\sin3u=\sin u\,(3-4\sin^2u)$), and $\varsigma_j$ is strictly increasing in $j$ because $ju\in[0,\frac\pi2)$. By \eqref{eq:boxeig} and Corollary~\ref{cor:box}(c), $\lambda_k=\Lambda_1V(k)$ with $V(k)=\sum_i\varsigma_{k_i}$; every eigenvalue is visible from the corner, so $\Lambda_0<\Lambda_1<\dots<\Lambda_p$ is the increasing list of the distinct values of $V$ (times $\Lambda_1$). With $R:x_i\mapsto N+1-x_i$ we have $\phi_k\circ R=(-1)^{|k|}\phi_k$; if all $k$ with $\lambda_k=\Lambda_i$ have $|k|$ even (odd), then $R$ acts as $+1$ ($-1$) on the range of $E_{\Lambda_i}$ and $E_{\Lambda_i}v_\mp=\frac12(E_{\Lambda_i}e_a\mp RE_{\Lambda_i}e_a)=0$, i.e.\ $\Lambda_i$ is purely even (purely odd) in the sense of Proposition~\ref{prop:pure}.

(a) For $d=1$ this is Corollary~\ref{cor:1d}. For $N=2$: $k\in\{0,1\}^d$ and $V(k)=|k|\in\{0,1,\dots,d\}$, so $p=d=D$. By Theorem~\ref{thm:residues}(d) the $p$ residues have at least $D-1=p-1$ sign changes, which forces all of them to be non-zero and of alternating sign, and $r_0>0$ by Theorem~\ref{thm:residues}(c) ($\Omega'$ is connected).

(b) We count from the top of the spectrum. For $0\le m\le N-1$ let $\Delta_m=\varsigma_{N-1}-\varsigma_{N-1-m}$, so $\Delta_0=0$. Since $\sin^2A-\sin^2B=\sin(A+B)\sin(A-B)$ and $(2N-2m-1)u=\pi-(2m+1)u$,
\begin{equation}\label{eq:drops}
\Delta_m-\Delta_{m-1}=\frac{\sin^2((N-m)u)-\sin^2((N-m-1)u)}{\sin^2u}=\frac{\sin((2m+1)u)}{\sin u}>0\qquad(1\le m\le N-1),
\end{equation}
because $0<(2m+1)u<2Nu=\pi$; so $\Delta_m$ is strictly increasing in $m$. For $k\in\{0,\dots,N-1\}^d$ put $m_l=N-1-k_l$ and $\Delta(m)=\sum_l\Delta_{m_l}$; then $V(k)=d\varsigma_{N-1}-\Delta(m)$ and $|k|=D-|m|$. Hence, if $0=\Delta^{(0)}<\Delta^{(1)}<\dots$ are the distinct values of $\Delta$ on $\{0,\dots,N-1\}^d$, then $\Lambda_{p-j}=\Lambda_1\big(d\varsigma_{N-1}-\Delta^{(j)}\big)$, and if all $m$ with $\Delta(m)=\Delta^{(j)}$ have $|m|$ of one parity, then $\Lambda_{p-j}$ is purely even if $D-|m|$ is even and purely odd if it is odd. Below a vector $m$ is described by the multiset of its non-zero entries, and $n_l$ denotes the number of entries equal to $l$. In this proof let
\[
\alpha=\Delta_1=\frac{\sin3u}{\sin u}=3-4x,\qquad\beta=\Delta_2-\Delta_1=\frac{\sin5u}{\sin u}=5-20x+16x^2
\]
(multiple-angle formulas). For $N\ge3$ we have $0<x\le\frac14$, hence $2\alpha-\beta=1+12x-16x^2>0$; and $\beta-\alpha=2(1-8x+8x^2)$ is positive for $x<\sin^2\frac\pi8=\frac{2-\sqrt2}{4}$, i.e.\ for $N\ge5$, and zero for $N=4$.

\emph{Case $N\ge5$.} Here $0<\alpha<\beta<2\alpha$. Let $\Delta(m)\le\alpha+\beta=\Delta_2$. As $\Delta_l$ increases with $l$, all entries of $m$ are $\le2$; if $m$ has an entry $2$, then $\Delta(m)\ge\Delta_2$ with equality only for $m=\{2\}$; otherwise $\Delta(m)=n_1\alpha$, and $n_1\alpha\le\alpha+\beta<3\alpha$ forces $n_1\le2$. So the four smallest values of $\Delta$ are $0<\alpha<2\alpha<\alpha+\beta$, attained only by $\emptyset$, $\{1\}$, $\{1,1\}$ (which exists because $d\ge2$) and $\{2\}$. In particular $\Lambda_{p-2}$ ($\Delta=2\alpha$) and $\Lambda_{p-3}$ ($\Delta=\alpha+\beta$) are pure, both with $|m|=2$.

\emph{Case $N=3$.} Here $x=\frac14$, $\Delta_1=2$, $\Delta_2=3$, so $\Delta(m)=2n_1+3n_2$; the four smallest values are $0<2<3<4$, attained only by $\emptyset$, $\{1\}$, $\{2\}$, $\{1,1\}$. So $\Lambda_{p-2}$ ($\Delta=3$) and $\Lambda_{p-3}$ ($\Delta=4$) are pure, both with $|m|=2$.

\emph{Case $N=4$.} Here $x=\sin^2\frac\pi8$, $\alpha=\beta=1+\sqrt2$ and $\Delta_3-\Delta_2=\sin\frac{7\pi}8/\sin\frac\pi8=1$, so $\Delta_1=\alpha$, $\Delta_2=2\alpha$, $\Delta_3=2\alpha+1$. Thus $\Delta(m)=c\,\alpha+n_3$ with $c=n_1+2n_2+2n_3$, while $|m|=n_1+2n_2+3n_3=c+n_3$. Since $\alpha$ is irrational, $\Delta(m)$ determines $(c,n_3)$ and hence $|m|$; so every eigenvalue is pure when $N=4$. As $2n_3\le c$ and $\alpha>1$, the values of $\Delta$ below $3\alpha$ are those with $c\le2$, namely $0<\alpha<2\alpha<2\alpha+1$, and the value $3\alpha$ is attained (by $\{1,2\}$, as $d\ge2$). So $\Lambda_{p-3}$ ($\Delta=2\alpha+1$, $(c,n_3)=(2,1)$) and $\Lambda_{p-4}$ ($\Delta=3\alpha$, $(c,n_3)=(3,0)$) are pure, both with $|m|=3$.

\emph{Conclusion.} The eigenvalue $\Lambda_0=0$ belongs to $m=(N-1,\dots,N-1)$, with $\Delta(m)=d\Delta_{N-1}\ge2\Delta_2$, which exceeds all the values found above ($2\alpha+2\beta>\alpha+\beta$ for $N\ge5$, $6>4$ for $N=3$, $4\alpha>3\alpha$ for $N=4$). Hence $p\ge4$ ($N\ne4$), resp.\ $p\ge5$ ($N=4$), so $1\le i$ and $i+1\le p-2$. The eigenvalues $\Lambda_i,\Lambda_{i+1}$ are pure with $|m|=2$ ($N\ne4$), resp.\ $|m|=3$ ($N=4$), so both are purely even iff $D-2$, resp.\ $D-3$, is even, and both are purely odd otherwise. By the last statement of Proposition~\ref{prop:pure} (with $l=2$), $(\operatorname{sign}r_{i-1},\operatorname{sign}r_i,\operatorname{sign}r_{i+1})$ is non-decreasing, resp.\ non-increasing, whereas alternation would require $(+,-,+)$ or $(-,+,-)$.

(c) $d=2$, $N\ge3$. Here $x\le\sin^2\frac\pi6=\frac14$, so $t\ge3$. If $V(k)\le t$, then no component of $k$ is $\ge3$ (otherwise $V(k)\ge\varsigma_3>\varsigma_2=t$), so $V(k)\in\{0,1,2,t,t+1,2t\}$. Hence the four smallest distinct values are $0<1<2<t$, i.e.\ $\Lambda_3=t\Lambda_1$ with eigenspace spanned by $\phi_{(2,0)},\phi_{(0,2)}$: $\Lambda_3$ is purely even, and $3\le p-1$ because $t+1$ and $2t$ are two further values. By Proposition~\ref{prop:pure}(a), $r_2\ge0$ implies $r_3>0$. Alternation would require $r_2>0>r_3$.

$d\ge4$, $N\ge4$. Here $0<x\le\sin^2\frac\pi8<0.1465$, so $3<t<4$ and $\varsigma_3=(3-4x)^2>5.8$. Let $V(k)\le4$. Then $k$ has no component $\ge3$ (as $\varsigma_3>4$) and at most one component equal to $2$ (as $2t>6$); if it has one, the other components vanish (as $t+1>4$) and $V(k)=t$; otherwise $V(k)=|k|\in\{0,1,2,3,4\}$. So the six smallest distinct values are $0<1<2<3<t<4$, $\Lambda_5=4\Lambda_1$, and its eigenspace is spanned by the $\phi_k$ with $k$ a permutation of $(1,1,1,1,0,\dots,0)$, $|k|=4$: $\Lambda_5$ is purely even. Further values exist ($t+1$, $2t$), so $5\le p-1$. By Proposition~\ref{prop:pure}(a), $r_4\ge0$ implies $r_5>0$; alternation would require $r_4>0>r_5$.

$d=3$, $N\ge7$. Here $0<x\le\sin^2\frac\pi{14}<0.0496$, so $3.8<t<4$ and $\varsigma_3-2t=1-16x+16x^2>0$. If $V(k)\le2t$, then $k$ has no component $\ge3$, so with $l$ components equal to $2$ and $j$ equal to $1$ ($j+l\le3$), $V(k)=lt+j\in\{0,1,2,3,t,t+1,t+2,2t,2t+1,3t\}$. As $3<t$ and $t+2<2t$, the eight smallest distinct values are $0<1<2<3<t<t+1<t+2<2t$, so $\Lambda_7=2t\Lambda_1$ with eigenspace spanned by the $\phi_k$, $k$ a permutation of $(2,2,0)$, $|k|=4$: purely even; and $7\le p-1$ ($2t+1$ and $3t$ are further values). By Proposition~\ref{prop:pure}(a), $r_6\ge0$ implies $r_7>0$; alternation would require $r_6>0>r_7$.
\end{proof}

\begin{remark}[the sign pattern as the sign of one polynomial; general chain]\label{rem:signpoly}
For an arbitrary chain satisfying \eqref{eq:A} and a start $x\ne a$ with $D=\dist(x,a)$, write $G_x(s)=\sum_{i=0}^p\omega_i/(\Lambda_i-s)$, $\omega_i=\ip{e_x}{E_{\Lambda_i}e_a}$. Then
\[
\mathcal N_x(s):=G_x(s)\prod_{l=0}^p(\Lambda_l-s)=\sum_{i=0}^p\omega_i\prod_{l\ne i}(\Lambda_l-s)
\]
is a real polynomial of degree \emph{exactly} $p-D$, with leading coefficient $(-1)^p(L^D)_{xa}$ of sign $(-1)^{p+D}$, and
\begin{equation}\label{eq:signpoly}
\operatorname{sign}r_j(x)=(-1)^j\operatorname{sign}\mathcal N_x(\sigma_j)\qquad(0\le j\le p-1),\qquad \operatorname{sign}\mathcal N_x(\Lambda_i)=(-1)^i\operatorname{sign}\omega_i\quad(0\le i\le p).
\end{equation}
Indeed, for $|s|>\|L\|$ the Neumann series gives $G_x(s)=-\sum_{m\ge0}(L^m)_{xa}s^{-m-1}$, where $(L^m)_{xa}=0$ for $m<D$ and $(L^D)_{xa}\neq0$ is a sum over geodesics of products of $D$ negative entries (as in the proof of Theorem~\ref{thm:residues}(d)); multiplying by $\prod_l(\Lambda_l-s)=(-1)^{p+1}s^{p+1}(1+O(s^{-1}))$ gives the degree and the leading coefficient. The first relation in \eqref{eq:signpoly} follows from \eqref{eq:res}, $\sigma_jG'(\sigma_j)>0$ and the fact that exactly $j+1$ factors of $\prod_l(\Lambda_l-\sigma_j)$ are negative; the second from the definition. Consequences: (1) $\operatorname{sign}r_j=(-1)^j$ for all $j$ iff $\mathcal N_x>0$ at every $\sigma_j$; in particular this holds whenever $p=D$ (then $\mathcal N_x$ is the constant $(-1)^p(L^D)_{xa}>0$), which is case (a) of Theorem~\ref{thm:noalt}. (2) If $r_jr_{j+1}>0$ then $\mathcal N_x$ has a zero in $(\sigma_j,\sigma_{j+1})$, so at most $p-D$ of the $p-1$ consecutive pairs fail to alternate; this is again the bound ``at least $D-1$ sign changes'' of Theorem~\ref{thm:residues}(d) when no residue vanishes. (3) The sequence $\big((-1)^i\operatorname{sign}\omega_i\big)_{1\le i\le p}$ (zeros deleted) has at most $p-D$ sign changes. (4) If the residues alternate, then $\mathcal N_x>0$ at $\sigma_{i-1}$ and $\sigma_i$, so every $i\in\{1,\dots,p-1\}$ with $(-1)^i\omega_i\le0$ forces two zeros of $\mathcal N_x$ (with multiplicity) in $(\sigma_{i-1},\sigma_i)$; hence alternation implies $2\,\#\{1\le i\le p-1:(-1)^i\omega_i\le0\}\le p-D$. For the corner-to-corner box $\omega_i=\frac1n\sum_{k:\lambda_k=\Lambda_i}(-1)^{|k|}\prod_l\varepsilon_{k_l}\cos^2\frac{\theta_{k_l}}2$, and this counting criterion (which needs only the ordered spectrum, no pole) excludes alternation in $95$ of the $104$ pairs $(d,N)$ with $d\ge2$, $N\ge3$ examined in script~\texttt{18}, but not in all of them (not for $(d,N)=(3,4),(3,20),(3,21),(3,23)$ and $(d,4)$, $4\le d\le8$); it is therefore not a substitute for Theorem~\ref{thm:noalt}, whose proof uses the finer information of Proposition~\ref{prop:pure}.
\end{remark}

\part*{Part II. Explicit bounds (any finite symmetric chain)}
\section{Explicit two-sided bounds on $\sigma_0$, $\sigma_1$, $r_0$, $r_1$ (general chain)}\label{sec:bounds}

This section is still for a general $L$ satisfying \eqref{eq:A}; all hypotheses are stated in terms of a few explicit functionals of the Green's function at the target. Their values for the box are computed in Sections~\ref{sec:corner}--\ref{sec:d3}; the resulting explicit statements are collected in Section~\ref{sec:explicit}.

\subsection{The Green's function functionals}
With the visible spectrum $0=\Lambda_0<\Lambda_1<\dots<\Lambda_p$ and weights $w_i$ of Definition~\ref{def:visible}, write $G(s)=-\frac1{ns}+g(s)$ and, for a start $x\in\Omega'$, $G_x(s)=-\frac1{ns}+g_x(s)$, where
\begin{equation}\label{eq:g}
g(s)=\sum_{i=1}^p\frac{w_i}{\Lambda_i-s},\qquad g_x(s)=\sum_{i=1}^p\frac{\sqrt{w_i}\,u_i(x)}{\Lambda_i-s}
\end{equation}
are the regular parts ($g(0)=\ip{e_a}{L^+e_a}$ and $g_x(0)=\ip{e_x}{L^+e_a}$ with $L^+$ the Moore--Penrose inverse). Define
\begin{equation}\label{eq:functionals}
\tau=n\,g(0)=n\sum_{i\ge1}\frac{w_i}{\Lambda_i},\qquad K_2=n\,g'(0)=n\sum_{i\ge1}\frac{w_i}{\Lambda_i^2},\qquad m(x)=n\big(g(0)-g_x(0)\big),
\end{equation}
\begin{equation}\label{eq:dimless}
\Xb=\Lambda_1\tau,\qquad \kappa=\frac{K_2}{\tau^2},\qquad W=n\,w_1,\qquad \mu_x=\frac{m(x)-\tau}{\tau}.
\end{equation}

\begin{lemma}[meaning of $\tau$ and $m$]\label{lem:taum}
For the unit-rate walk ($q=1$), $m(x)=\E_xT_a$ for every $x$ (with $m(a)=0$), and $\tau=\frac1n\sum_{x\in\Omega}\E_xT_a$ is the mean hitting time from the uniform (stationary) distribution. For activity $q$ both are divided by $q$, in continuous and in discrete time. Moreover $\kappa\le1/\Xb$.
\end{lemma}
\begin{proof}
Let $h(x)=n(g(0)-g_x(0))=n\ip{e_a-e_x}{L^+e_a}$. Since $LL^+=I-\frac1n\one\one^{\mathsf T}$, $(Lh)(x)=-n\,(LL^+e_a)(x)=-n(\delta_{xa}-\frac1n)=1$ for $x\ne a$, and $h(a)=0$. Hence $L_a(h|_{\Omega'})=\one_{\Omega'}$, i.e.\ $h|_{\Omega'}=L_a^{-1}\one_{\Omega'}$, which is $\int_0^\infty e^{-tL_a}\one\,dt=\E_\cdot T_a$ by \eqref{eq:surv} (and $\sum_{t\ge0}Q^t\one=(qL_a)^{-1}\one$ in discrete time). Averaging over $x$: $\frac1n\sum_xh(x)=n g(0)-\ip{\one}{L^+e_a}=ng(0)$. Finally $K_2=n\sum_{i\ge1}w_i\Lambda_i^{-2}\le\Lambda_1^{-1}\tau$, so $\kappa\le1/(\Lambda_1\tau)$.
\end{proof}

We shall use repeatedly the elementary inequalities, valid for $0\le s<\Lambda\le\Lambda'$:
\begin{equation}\label{eq:elem}
\frac1{\Lambda'}+\frac s{\Lambda'^2}\ \le\ \frac1{\Lambda'-s}=\frac1{\Lambda'}+\frac{s}{\Lambda'(\Lambda'-s)}\ \le\ \frac1{\Lambda'}+\frac{s}{\Lambda'^2}\cdot\frac1{1-s/\Lambda},\qquad
\frac1{\Lambda'^2}\le\frac1{(\Lambda'-s)^2}\le\frac1{\Lambda'^2(1-s/\Lambda)^2}.
\end{equation}

\subsection{The principal eigenvalue}

\begin{theorem}[principal killed eigenvalue]\label{thm:sigma0}
Assume \eqref{eq:A}. Then
\begin{equation}\label{eq:s0basic}
\tau\ \le\ \frac1{\sigma_0}\ \le\ \tau+\frac1{\Lambda_1},\qquad\text{i.e.}\qquad \frac{\Xb}{1+\Xb}\le\sigma_0\tau\le1 .
\end{equation}
If $\Xb>1$, put $\beta=\Xb/(\Xb-1)$. Then
\begin{equation}\label{eq:s0sharp}
1-\beta\kappa\ \le\ \frac{2}{1+\sqrt{1+4\beta\kappa}}\ \le\ \sigma_0\tau\ \le\ \frac{2}{1+\sqrt{1+4\kappa}}\ \le\ 1-\kappa+2\kappa^2 .
\end{equation}
\end{theorem}
\begin{proof}
$\sigma_0$ is the zero of $G$ in $(0,\Lambda_1)$, i.e.\ $1=\sigma_0\,n\,g(\sigma_0)$. By \eqref{eq:elem} with $\Lambda=\Lambda_1$, for $0\le s<\Lambda_1$,
\[
\tau+sK_2\ \le\ n\,g(s)\ \le\ \tau+\frac{sK_2}{1-s/\Lambda_1},\qquad n\,g(s)\le\frac{\tau}{1-s/\Lambda_1}.
\]
The first inequality at $s=\sigma_0$ gives $1\ge\sigma_0\tau$ (hence $\sigma_0\le1/\tau$) and, with $t=\sigma_0\tau$, $1\ge t+\kappa t^2$, i.e.\ $t\le 2/(1+\sqrt{1+4\kappa})$. The third gives $1-\sigma_0/\Lambda_1\le\sigma_0\tau$, i.e.\ $1/\sigma_0\le\tau+1/\Lambda_1$. If $\Xb>1$ then $\sigma_0/\Lambda_1\le1/\Xb<1$, and the second inequality gives $1\le t+\beta\kappa t^2$, i.e.\ $t\ge2/(1+\sqrt{1+4\beta\kappa})$. The outer inequalities in \eqref{eq:s0sharp} are $\frac2{1+\sqrt{1+4y}}\ge\frac1{1+y}\ge1-y$ and $\frac{2}{1+\sqrt{1+4y}}=\frac{\sqrt{1+4y}-1}{2y}\le1-y+2y^2$ (from $\sqrt{1+z}\le1+\frac z2-\frac{z^2}8+\frac{z^3}{16}$, $z\ge0$).
\end{proof}

\begin{remark}
\eqref{eq:s0basic} says that the quasi-stationary mean exit time $1/\sigma_0$ lies between the stationary mean hitting time and that time plus the relaxation time $1/\Lambda_1$ of the visible sector. For reversible chains inequalities of this type go back to D.~Aldous and M.~Brown, \emph{Inequalities for rare events in time-reversible Markov chains I}, in: Stochastic Inequalities, IMS Lecture Notes 22 (1992), 1--16; the proof above is self-contained and gives the refinement \eqref{eq:s0sharp}, whose two sides agree to first order in $\kappa$.
\end{remark}

\begin{corollary}[$\sigma_0\cdot\mathrm{MFPT}$]\label{cor:s0m}
For every start $x\in\Omega'$: $\sigma_0\,m(x)=\sigma_0\tau\,(1+\mu_x)$, hence, if $\Xb>1$,
\[
(1-\beta\kappa)(1+\mu_x)\ \le\ \sigma_0\,m(x)\ \le\ (1-\kappa+2\kappa^2)(1+\mu_x)\qquad(\text{for }1+\mu_x\ge0),
\]
and always $\frac{\Xb}{1+\Xb}(1+\mu_x)\le\sigma_0m(x)\le1+\mu_x$.
\end{corollary}

\subsection{The second pole and the gap}
Assume $p\ge2$ and put
\begin{equation}\label{eq:h}
h(s)=\sum_{i=2}^p\frac{w_i}{\Lambda_i-s},\qquad Y_0=n\Lambda_1h(0)=\Xb-W,\qquad \vartheta=n\Lambda_1^2h'(0)=\Lambda_1^2K_2-W,\qquad\varrho=\frac{\Lambda_1}{\Lambda_2}<1 .
\end{equation}

\begin{theorem}[second pole]\label{thm:sigma1}
Assume \eqref{eq:A}, $p\ge2$, and write $\sigma_1=\Lambda_1(1+\delta)$, $\delta>0$. Then
\begin{equation}\label{eq:deltaexact}
\delta=\frac{W}{Y(\sigma_1)-\frac1{1+\delta}},\qquad Y(s):=n\Lambda_1h(s).
\end{equation}
\begin{enumerate}
\item[(i)] If $Y_0+\vartheta>1$, then $\delta\le\delta_+:=\dfrac{W}{Y_0+\vartheta-1}$.
\item[(ii)] If moreover $(1+\delta_+)\varrho<1$, put $\omega=\dfrac{1+\delta_+}{1-(1+\delta_+)\varrho}$. Then $\delta\ge\delta_-:=\dfrac{W}{Y_0+\omega\vartheta-\frac1{1+\delta_+}}$, and
\[
Y_0+\vartheta\le Y(\sigma_1)\le Y_0+\omega\vartheta,\qquad \vartheta\le n\Lambda_1^2h'(\sigma_1)\le\frac{\omega^2\vartheta}{(1+\delta_+)^2}.
\]
\end{enumerate}
\end{theorem}
\begin{proof}
$G(\sigma_1)=0$ reads $-\frac1{n\sigma_1}-\frac{w_1}{\sigma_1-\Lambda_1}+h(\sigma_1)=0$; multiplying by $n\Lambda_1$ gives $W/\delta=Y(\sigma_1)-\Lambda_1/\sigma_1$, which is \eqref{eq:deltaexact}.
(i) By \eqref{eq:elem} (first inequality, with $\Lambda'=\Lambda_i$, $i\ge2$, and $s=\sigma_1<\Lambda_2$), $h(\sigma_1)\ge h(0)+\sigma_1h'(0)\ge h(0)+\Lambda_1h'(0)$, so $Y(\sigma_1)\ge Y_0+\vartheta$; and $\Lambda_1/\sigma_1<1$. Hence $W/\delta>Y_0+\vartheta-1>0$.
(ii) By (i), $\sigma_1\le(1+\delta_+)\Lambda_1<\Lambda_2$, and \eqref{eq:elem} (second inequality, with $\Lambda=\Lambda_2$) gives $h(\sigma_1)\le h(0)+\sigma_1h'(0)/(1-\sigma_1/\Lambda_2)$, i.e.\ $Y(\sigma_1)\le Y_0+\frac{(1+\delta)\vartheta}{1-(1+\delta)\varrho}\le Y_0+\omega\vartheta$, the last step because $y\mapsto(1+y)/(1-(1+y)\varrho)$ is increasing. Also $\Lambda_1/\sigma_1\ge1/(1+\delta_+)$. Hence $W/\delta\le Y_0+\omega\vartheta-\frac1{1+\delta_+}$ (the right-hand side is positive, being $\ge W/\delta$). The bounds on $h'(\sigma_1)$ follow from the last inequality in \eqref{eq:elem}.
\end{proof}

\begin{corollary}[spectral gap]\label{cor:gap}
Always $\sigma_1-\sigma_0>\Lambda_1-\frac1\tau=\Lambda_1(1-\Xb^{-1})$ and $\sigma_1-\sigma_0<\Lambda_2$. Under the hypotheses of Theorem~\ref{thm:sigma1}(ii) and $\Xb>1$,
\[
\Lambda_1\Big(1+\delta_--\frac{1}{\Xb}\Big)\ \le\ \sigma_1-\sigma_0\ \le\ \Lambda_1\Big(1+\delta_+-\frac{1}{1+\Xb}\Big),\qquad
\Xb(1+\delta_-)\le\frac{\sigma_1}{\sigma_0}\le(1+\Xb)(1+\delta_+).
\]
\end{corollary}
\begin{proof}
Interlacing \eqref{eq:strict}, \eqref{eq:s0basic} and Theorem~\ref{thm:sigma1}.
\end{proof}

\subsection{The residues $r_0$ and $r_1$}
For the start $x$ let $K_2^{(x)}=n\sum_{i\ge1}\|E_ie_x\|^2/(\lambda^{(i)})^2$ (the functional $K_2$ with $x$ as target) and
\[
\kappa_x=\frac{\sqrt{K_2K_2^{(x)}}}{\tau^2},\qquad A_x=-n\sqrt{w_1}\,u_1(x)=-n\,\ip{e_x}{E_{\Lambda_1}e_a},\qquad
\vartheta_x=\Lambda_1^2\sqrt{\big(K_2-\tfrac{W}{\Lambda_1^2}\big)K_2^{(x)}} .
\]
(In a symmetric situation such as corner to corner, $K_2^{(x_0)}=K_2$, so $\kappa_{x_0}=\kappa$, and $\vartheta_{x_0}\le\sqrt{\vartheta(\vartheta+W)}$; for CC one may take $\vartheta_{x_0}=\vartheta$, see the proof.)

\begin{theorem}[the residue $r_0$]\label{thm:r0}
Assume \eqref{eq:A} and $\Xb>1$, $\beta=\Xb/(\Xb-1)$, $t=\sigma_0\tau$. Then
\begin{equation}\label{eq:r0}
r_0(x)=\frac{1+t\,\mu_x+\varepsilon_1}{1+\varepsilon_2},\qquad|\varepsilon_1|\le\beta\kappa_xt^2,\qquad\kappa t^2\le\varepsilon_2\le\beta^2\kappa t^2 .
\end{equation}
In particular, with $t_-\le t\le t_+$ the bounds of \eqref{eq:s0sharp}, if $\mu_x\ge0$ and $\beta\kappa_xt_+^2\le1$,
\[
\frac{1+t_-\mu_x-\beta\kappa_xt_+^2}{1+\beta^2\kappa t_+^2}\ \le\ r_0(x)\ \le\ \frac{1+t_+\mu_x+\beta\kappa_xt_+^2}{1+\kappa t_-^2} .
\]
Moreover the stationary-start weight \eqref{eq:varrho} of the principal pole is $\varrho_0=1/(1+\varepsilon_2)$, so that
\begin{equation}\label{eq:rho0}
1-\beta^2\kappa\ \le\ \varrho_0\ \le 1,\qquad \sum_{j\ge1}\varrho_j\ \le\ \beta^2\kappa .
\end{equation}
\end{theorem}
\begin{proof}
By \eqref{eq:res}, $r_0=-G_x(\sigma_0)/(\sigma_0G'(\sigma_0))$ with $G_x(s)=-\frac1{ns}+g_x(s)$, $G'(s)=\frac1{ns^2}+g'(s)$; multiplying numerator and denominator by $n\sigma_0$,
\[
r_0=\frac{1-n\sigma_0\,g_x(\sigma_0)}{1+n\sigma_0^2\,g'(\sigma_0)} .
\]
Now $n\,g_x(\sigma_0)=n\,g_x(0)+\sigma_0\,n\sum_{i\ge1}\frac{\sqrt{w_i}u_i(x)}{\Lambda_i(\Lambda_i-\sigma_0)}$ and $-n\sigma_0g_x(0)=\sigma_0(m(x)-\tau)=t\mu_x$. By Cauchy--Schwarz and $\sqrt{w_i}|u_i(x)|=|\ip{E_ie_x}{E_ie_a}|\le\sqrt{w_i}\,\|E_ie_x\|$,
\[
n\sum_{i\ge1}\frac{\sqrt{w_i}|u_i(x)|}{\Lambda_i(\Lambda_i-\sigma_0)}\le\frac1{1-\sigma_0/\Lambda_1}\,n\sum_{i\ge1}\frac{\sqrt{w_i}\,\|E_ie_x\|}{\Lambda_i^2}\le\beta\sqrt{K_2K_2^{(x)}},
\]
which gives $|\varepsilon_1|\le\sigma_0^2\beta\sqrt{K_2K_2^{(x)}}=\beta\kappa_xt^2$. Finally $\varepsilon_2=n\sigma_0^2g'(\sigma_0)\in[\sigma_0^2K_2,\beta^2\sigma_0^2K_2]$ by \eqref{eq:elem}. The two-sided bound follows because the numerator lies in $[1+t_-\mu_x-\beta\kappa_xt_+^2,\,1+t_+\mu_x+\beta\kappa_xt_+^2]$, the lower end being $\ge0$ by assumption, and the denominator in $[1+\kappa t_-^2,1+\beta^2\kappa t_+^2]$. For \eqref{eq:rho0}: $\varrho_0=1/(n\sigma_0^2G'(\sigma_0))=1/(1+n\sigma_0^2g'(\sigma_0))=1/(1+\varepsilon_2)\ge1-\beta^2\kappa t^2\ge1-\beta^2\kappa$, and $\sum_{j\ge1}\varrho_j=1-\frac1n-\varrho_0$ by Proposition~\ref{prop:uniform}(a).
\end{proof}

\begin{theorem}[the residue $r_1$]\label{thm:r1size}
Assume \eqref{eq:A}, $p\ge2$, and let $\delta=\sigma_1/\Lambda_1-1$, $\varsigma=1/(1+\delta)$. Then
\begin{equation}\label{eq:r1exact}
r_1(x)=-\,\frac{\delta\,\big[A_x+\delta\,(Y_x(\sigma_1)-\varsigma)\big]}{(1+\delta)W+\delta^2\big[\varsigma+(1+\delta)\,n\Lambda_1^2h'(\sigma_1)\big]},
\qquad Y_x(s):=n\Lambda_1\sum_{i\ge2}\frac{\sqrt{w_i}\,u_i(x)}{\Lambda_i-s}.
\end{equation}
Moreover $Y_x(0)=A_x-\Lambda_1(m(x)-\tau)$, and under the hypotheses of Theorem~\ref{thm:sigma1}(ii), $|Y_x(\sigma_1)-Y_x(0)|\le\omega\vartheta_x$. Consequently, with
$y_-=Y_x(0)-\omega\vartheta_x-1$, $y_+=Y_x(0)+\omega\vartheta_x-\frac1{1+\delta_+}$,
\[
\mathcal N_-=A_x+\min(\delta_-y_-,\delta_+y_-),\qquad \mathcal N_+=A_x+\max(\delta_-y_+,\delta_+y_+),
\]
\[
\mathcal D_-=(1+\delta_-)W+\delta_-^2\Big[\frac1{1+\delta_+}+(1+\delta_-)\vartheta\Big],\qquad
\mathcal D_+=(1+\delta_+)W+\delta_+^2\Big[1+\frac{\omega^2\vartheta}{1+\delta_+}\Big],
\]
one has, provided $\mathcal N_->0$,
\begin{equation}\label{eq:r1bounds}
\frac{\delta_-\,\mathcal N_-}{\mathcal D_+}\ \le\ -r_1(x)\ \le\ \frac{\delta_+\,\mathcal N_+}{\mathcal D_-} .
\end{equation}
\end{theorem}
\begin{proof}
Since $\sqrt{w_1}u_1(x)=-A_x/n$,
\[
n\Lambda_1G_x(\sigma_1)=-\frac{\Lambda_1}{\sigma_1}+\frac{-A_x}{\Lambda_1-\sigma_1}\Lambda_1+Y_x(\sigma_1)=-\varsigma+\frac{A_x}{\delta}+Y_x(\sigma_1),
\]
\[
n\Lambda_1\sigma_1G'(\sigma_1)=\frac{\Lambda_1}{\sigma_1}+\frac{\sigma_1\Lambda_1\,nw_1}{(\sigma_1-\Lambda_1)^2}+\sigma_1n\Lambda_1h'(\sigma_1)=\varsigma+\frac{(1+\delta)W}{\delta^2}+(1+\delta)\,n\Lambda_1^2h'(\sigma_1).
\]
Dividing and multiplying through by $\delta^2$ gives \eqref{eq:r1exact}. From $n g_x(0)=-(m(x)-\tau)$ and $g_x(0)=\sqrt{w_1}u_1(x)/\Lambda_1+\sum_{i\ge2}\sqrt{w_i}u_i(x)/\Lambda_i$ we get $Y_x(0)=A_x-\Lambda_1(m(x)-\tau)$. As in the proof of Theorem~\ref{thm:r0}, using \eqref{eq:elem} with $\Lambda=\Lambda_2$,
\[
|Y_x(\sigma_1)-Y_x(0)|\le n\Lambda_1\sigma_1\sum_{i\ge2}\frac{\sqrt{w_i}|u_i(x)|}{\Lambda_i(\Lambda_i-\sigma_1)}\le\frac{(1+\delta)}{1-(1+\delta)\varrho}\,\Lambda_1^2\,n\sum_{i\ge2}\frac{\sqrt{w_i}\,\|E_ie_x\|}{\Lambda_i^2}\le\omega\vartheta_x .
\]
(The last step is Cauchy--Schwarz. In a symmetric situation $x_0=Ra$, $\|E_ie_{x_0}\|=\|E_ie_a\|=\sqrt{w_i}$, the middle expression equals $\frac{1+\delta}{1-(1+\delta)\varrho}\,\Lambda_1^2\,n\sum_{i\ge2}w_i\Lambda_i^{-2}\le\omega\vartheta$, so one may take $\vartheta_{x_0}=\vartheta$.) The bracket in the numerator of \eqref{eq:r1exact} therefore lies in $[\mathcal N_-,\mathcal N_+]$ because $\delta\in[\delta_-,\delta_+]$, $\varsigma\in[\frac1{1+\delta_+},1]$ and $Y_x(\sigma_1)-\varsigma\in[y_-,y_+]$; the denominator lies in $[\mathcal D_-,\mathcal D_+]$ by Theorem~\ref{thm:sigma1}(ii). This gives \eqref{eq:r1bounds}.
\end{proof}

\begin{remark}[using bounds on the functionals instead of their exact values]\label{rem:monotone}
In applications $K_2$ (hence $\kappa$ and $\vartheta$) is only known between two bounds. The proofs above use the functionals monotonically, so the statements remain true in the following form. Let $\kappa_-\le\kappa\le\kappa_+$ and $\vartheta_-\le\vartheta\le\vartheta_+$ (and $\vartheta_x\le\vartheta_+$).
(a) In \eqref{eq:s0sharp}: $2/(1+\sqrt{1+4\beta\kappa_+})\le\sigma_0\tau\le2/(1+\sqrt{1+4\kappa_-})$; in \eqref{eq:r0} and \eqref{eq:rho0} $\kappa$ may be replaced by $\kappa_+$ in the upper bounds for $|\varepsilon_1|$, $\varepsilon_2$, $\sum_{j\ge1}\varrho_j$.
(b) In Theorem~\ref{thm:sigma1}: if $Y_0+\vartheta_->1$ then $\delta\le\hat\delta_+:=W/(Y_0+\vartheta_--1)$ [because $Y(\sigma_1)\ge Y_0+\vartheta\ge Y_0+\vartheta_-$]; if moreover $(1+\hat\delta_+)\varrho<1$ and $\hat\omega:=(1+\hat\delta_+)/(1-(1+\hat\delta_+)\varrho)$, then $\delta\ge\hat\delta_-:=W/(Y_0+\hat\omega\vartheta_+-\frac1{1+\hat\delta_+})$, $Y(\sigma_1)\le Y_0+\hat\omega\vartheta_+$ and $\vartheta_-\le n\Lambda_1^2h'(\sigma_1)\le\hat\omega^2\vartheta_+/(1+\hat\delta_+)^2$ [because $y\mapsto(1+y)/(1-(1+y)\varrho)$ is increasing and $\delta\le\hat\delta_+$].
(c) In Theorem~\ref{thm:r1size}: \eqref{eq:r1bounds} holds with $\delta_\pm,\omega$ replaced by $\hat\delta_\pm,\hat\omega$, with $\omega\vartheta_x$ replaced by $\hat\omega\vartheta_+$ in $y_\pm$, with $\vartheta$ replaced by $\vartheta_-$ in $\mathcal D_-$ and by $\vartheta_+$ in $\mathcal D_+$, and with any lower (upper) bound for $Y_x(0)$ in $y_-$ ($y_+$) [the denominator of \eqref{eq:r1exact} is increasing in $\delta$ and in $h'(\sigma_1)$, and $\delta^2\varsigma=\delta^2/(1+\delta)$ is increasing in $\delta$].
\end{remark}

\begin{remark}[first-order content]\label{rem:firstorder}
When $\Xb\to\infty$ with $W$, $\vartheta$, $A_x$, $\Lambda_1(m(x)-\tau)$ bounded and $\varrho$ bounded away from $1$ (this is the situation of the box, Section~\ref{sec:explicit}), the theorems of this section give, with explicit error terms,
\[
\sigma_0\tau=1-\kappa+O(\Xb^{-3}),\quad
\frac{\sigma_1}{\Lambda_1}=1+\frac{W}{\Xb}\big(1+O(\Xb^{-1})\big),\quad
r_0=1+\mu_x+O(\Xb^{-2}),\quad r_1=-\frac{A_x}{\Xb}\big(1+O(\Xb^{-1})\big),
\]
which are the relations ``$\nu_0=\mathrm{MFPT}^{-1}(1+O(1/X))$, $\nu_1=\mu(1+W/X+\dots)$, $a_0\approx\nu_0$, $a_1\approx-A\nu_0$'' obtained by formal perturbation theory in Section~6 of the article ($a_j=r_j\sigma_j$), here with proofs and constants.
\end{remark}

\part*{Part III. The reflecting box: corner to corner}
\section{Corner to corner in the reflecting box: general dimension}\label{sec:corner}

From now on $L$ is the generator \eqref{eq:boxL} of the reflecting box, $d\ge1$, $N\ge2$, $n=N^d$, the target is the corner $a=(N,\dots,N)$ and the start is the opposite corner $x_0=(1,\dots,1)$ (``CC''). We write
\[
u=\frac{\pi}{2N},\qquad c_k=\cos\theta_k=\cos\frac{\pi k}{N},\qquad w_k=n\,\phi_k(a)^2=\prod_{i=1}^d\varepsilon_{k_i}\cos^2(k_iu)\qquad(k\in\{0,\dots,N-1\}^d),
\]
so that, by Lemma~\ref{lem:boxspec}, \eqref{eq:cornerweights}, \eqref{eq:functionals} and Corollary~\ref{cor:box},
\begin{equation}\label{eq:ccsums}
\begin{gathered}\tau=\sum_{k\ne0}\frac{w_k}{\lambda_k},\qquad m:=m(x_0)=\sum_{k\ne0}\frac{\big(1-(-1)^{|k|}\big)w_k}{\lambda_k},\qquad K_2=\sum_{k\ne0}\frac{w_k}{\lambda_k^2},\\
\Lambda_1=\frac2d\sin^2u,\qquad W=2d\cos^2u .\end{gathered}
\end{equation}
By Lemma~\ref{lem:taum}, $\tau/q$ is the mean hitting time of the corner from a uniformly distributed start and $m/q$ is the corner-to-corner MFPT of the article. The point reflection $R:x_i\mapsto N+1-x_i$ exchanges $a$ and $x_0$ and commutes with $L$; consequently $\|E_ie_{x_0}\|=\|E_ie_a\|$ for all $i$, and in the notation of Section~\ref{sec:bounds}
\begin{equation}\label{eq:ccdata}
K_2^{(x_0)}=K_2,\qquad\kappa_{x_0}=\kappa,\qquad A_{x_0}=W,\qquad \vartheta_{x_0}\ \text{may be replaced by}\ \vartheta,\qquad\varrho=\tfrac12\quad(d\ge2),
\end{equation}
where $A_{x_0}=-n\sum_{i}\phi_{e_i}(a)\phi_{e_i}(x_0)=d\cdot2\cos^2u$ by \eqref{eq:cornerweights}.

\begin{lemma}[elementary inequalities]\label{lem:elem}
\begin{enumerate}
\item[(a)] $\sin v\ge\frac{2}{\pi}v$ for $0\le v\le\frac\pi2$, and $v^2(1-\frac{v^2}3)\le\sin^2v\le v^2$ for $0\le v\le1$.
\item[(b)] $\big(\frac{\sin v}{v}\big)^2\ge\cos v$ for $0<v\le\frac\pi2$.
\item[(c)] $\operatorname{arsinh}$ is concave on $[0,\infty)$; hence $\operatorname{arsinh}\rho\ge\rho\,\operatorname{arsinh}(\rho_0)/\rho_0$ for $0\le\rho\le\rho_0$.
\item[(d)] $0\le x-\operatorname{arsinh}(\sin x)\le\frac{x^3}{3}$ and $\operatorname{arsinh}(\sin x)\ge\frac{2\operatorname{arsinh}(1)}{\pi}x$ for $0\le x\le\frac\pi2$.
\item[(e)] $(1-y)^4(1+2y)\le1$ for $0\le y\le1$.
\item[(f)] $1-\frac y2-\frac{y^2}2\le\sqrt{1-y}\le1-\frac y2$ for $0\le y\le1$.
\end{enumerate}
\end{lemma}
\begin{proof}
(a) Jordan's inequality (concavity of $\sin$ on $[0,\frac\pi2]$). $\sin v\le v$, and $\sin v\ge v-\frac{v^3}6\ge0$ gives $\sin^2v\ge v^2(1-\frac{v^2}6)^2\ge v^2(1-\frac{v^2}3)$.
(b) Let $h(v)=\sin^2v-v^2\cos v$. Then $h(0)=0$ and $h'(v)=2\cos v(\sin v-v)+v^2\sin v\ge-\frac{v^3}3\cos v+v^2\sin v=v^2\big(\sin v-\frac v3\cos v\big)\ge0$, because $\sin v-v\ge-\frac{v^3}{6}$ and $\tan v\ge v$.
(c) $\operatorname{arsinh}''(\rho)=-\rho(1+\rho^2)^{-3/2}\le0$.
(d) Let $A(x)=\operatorname{arsinh}(\sin x)$; $A'(x)=\cos x/\sqrt{1+\sin^2x}\in[0,1]$ is decreasing on $[0,\frac\pi2]$, so $A$ is concave with $A(0)=0$, $A\le x$, and $A(x)\ge xA(\frac\pi2)/\frac\pi2$. With $s=\sin x$: $A'(x)=\sqrt{(1-s^2)/(1+s^2)}\ge1-s^2$ (square both sides: $1\ge(1-s^2)(1+s^2)$), so $\frac{d}{dx}\big(\frac{x^3}3-x+A(x)\big)\ge x^2-\sin^2x\ge0$.
(e) $(1-y)^2(1+2y)=1-3y^2+2y^3\le1$ and $(1-y)^2\le1$.
(f) The upper bound is concavity. For the lower bound, $F(y)=\sqrt{1-y}-1+\frac y2+\frac{y^2}2$ has $F(0)=F(1)=0$ and $F'(y)=\frac12+y-\frac1{2\sqrt{1-y}}$, which vanishes on $(0,1)$ only where $(1+2y)^2(1-y)=1$, i.e.\ $y(3-4y^2)=0$, i.e.\ at $y=\sqrt3/2$; since $F''(0)=\frac34>0$, $F$ increases on $[0,\frac{\sqrt3}2]$ and decreases afterwards, hence $F\ge0$.
\end{proof}

\begin{lemma}[arrival profile; $m-\tau$]\label{lem:arrival}
Let $u_1(s)=1+2\sum_{k=1}^{N-1}(-1)^k\cos^2(ku)\,e^{-(1-c_k)s}$. For the unit-rate walk, $n\,p_t(x_0,a)=u_1(t/d)^d$,
\begin{equation}\label{eq:profile}
0\le1-n\,p_t(x_0,a)\le\min\big\{1,\ W e^{-\Lambda_1t}\big\}\qquad(t\ge0),
\end{equation}
and therefore
\begin{equation}\label{eq:mtau}
\frac{d\,(N^2-1)}{3}\ \le\ m-\tau=\int_0^\infty\big(1-n\,p_t(x_0,a)\big)\,dt\ \le\ \frac{1+\ln W}{\Lambda_1},
\end{equation}
with equality on the left if and only if $d=1$. Equivalently, with $\mu:=\mu_{x_0}=(m-\tau)/\tau$,
\begin{equation}\label{eq:muminus}
\frac{\mu_-(N)}{\Xb}\ \le\ \mu\ \le\ \frac{1+\ln W}{\Xb},\qquad
\mu_-(N):=\frac23\,(N^2-1)\sin^2\frac{\pi}{2N}\ \ge\ \max\Big\{1,\ \frac{\pi^2}{6}\Big(1-\frac1{N^2}\Big)\Big(1-\frac{\pi^2}{12N^2}\Big)\Big\},
\end{equation}
and $\mu_-(N)\to\pi^2/6$. In particular $m-\tau\ge1/\Lambda_1$: the corner-to-corner mean first-passage time exceeds the stationary-start mean hitting time by at least one relaxation time.
\end{lemma}
\begin{proof}
For $d=1$ the spectral theorem and Lemma~\ref{lem:boxspec} give
\[N\,p^{(1)}_s(1,N)=\sum_{k}N\phi_k(1)\phi_k(N)e^{-(1-c_k)s}=u_1(s),\]
using $N\phi_k(1)\phi_k(N)=\varepsilon_k(-1)^k\cos^2(ku)$. Hence $u_1\ge0$. Moreover $1-u_1(s)=2\sum_{k=1}^{N-1}(-1)^{k+1}b_k$ with $b_k=\cos^2(ku)e^{-(1-c_k)s}$, and $b_1>b_2>\dots>b_{N-1}>0$ since both factors decrease in $k$. An alternating sum of decreasing positive terms lies between $0$ and its first term: $0\le1-u_1(s)\le2\cos^2(u)\,e^{-(1-c_1)s}$. In dimension $d$, $L=\frac1d\sum_iL^{(i)}$ with commuting one-dimensional operators, so $e^{-tL}=\bigotimes_ie^{-(t/d)L^{(i)}}$ and $n\,p_t(x_0,a)=u_1(t/d)^d$. With $v=u_1(t/d)\in[0,1]$, $0\le1-v^d\le\min\{1,d(1-v)\}\le\min\{1,2d\cos^2(u)e^{-(1-c_1)t/d}\}$, which is \eqref{eq:profile}.
Next, $L^+=\int_0^\infty(e^{-tL}-E_0)\,dt$, so by \eqref{eq:functionals} $m-\tau=-n\,g_{x_0}(0)=-n\ip{e_{x_0}}{L^+e_a}=\int_0^\infty(1-n\,p_t(x_0,a))\,dt$. The integrand is non-negative, continuous and equal to $1$ at $t=0$, so $m>\tau$. Since $W=2d\cos^2u\ge2d\cos^2\frac\pi4=d\ge1$, with $t_*=\Lambda_1^{-1}\ln W\ge0$ we get $m-\tau\le t_*+\int_{t_*}^\infty We^{-\Lambda_1t}dt=(1+\ln W)/\Lambda_1$.

For the lower bound, $v=u_1(t/d)\in[0,1]$ gives $1-v^d\ge1-v$, with strict inequality when $d\ge2$ and $0<v<1$ (which holds for all $t>0$: $u_1(s)<1$ by the first part since $1-u_1(s)\ge2(b_1-b_2)>0$ for $N\ge3$ and $=2b_1>0$ for $N=2$, and $u_1(s)>0$ for $s>0$ because the chain is irreducible). Hence
\begin{align*}
m-\tau\ &\ge\ \int_0^\infty\big(1-u_1(t/d)\big)\,dt=d\int_0^\infty\big(1-u_1(s)\big)\,ds=d\sum_{k=1}^{N-1}(-1)^{k+1}\frac{2\cos^2(ku)}{1-c_k}\\
&=d\sum_{k=1}^{N-1}(-1)^{k+1}\cot^2(ku)=\frac{d\,(N^2-1)}3 ,
\end{align*}
the last step by Lemma~\ref{lem:cot} below (whose proof does not use the present lemma); the integration of the finite sum is termwise. Multiplying by $\Lambda_1=\frac2d\sin^2u$ gives $\Lambda_1(m-\tau)\ge\mu_-(N)$. Finally $\mu_-(2)=\frac23\cdot3\cdot\frac12=1$, and for $N\ge3$, by Lemma~\ref{lem:elem}(a), $\mu_-(N)\ge\frac23(N^2-1)u^2(1-\frac{u^2}3)=\frac{\pi^2}{6}(1-N^{-2})(1-\frac{\pi^2}{12N^2})$, which is increasing in $N$ and equals $1.3285\ldots>1$ at $N=3$.
\end{proof}
\begin{remark}
For $d=1$ the lower bound in \eqref{eq:mtau} is the exact value: $m_1=N(N-1)$, $\tau_1=\frac13(N-1)(2N-1)$, $m_1-\tau_1=\frac13(N^2-1)$. For $d=2$ and $d=3$ the lower bound is the leading term of the exact reductions \eqref{eq:tau2exact}, \eqref{eq:m3exact}; the true limits of $\Lambda_1(m-\tau)$ are $\pi\ln2=2.1776$ and $2.5194$ (Theorem~\ref{thm:asym}), between $\pi^2/6=1.6449$ and $1+\ln2d$.
\end{remark}

\begin{lemma}[second moment functional]\label{lem:K2}
For $d\ge1$, $N\ge2$,
\begin{equation}\label{eq:K2}
W+\binom d2\cos^4u\ \le\ \Lambda_1^2K_2\ \le\ Z_d(N):=\sum_{\substack{k\in\Z^d\setminus\{0\}\\ |k|_\infty\le N-1}}\frac{1}{|k|^4},
\end{equation}
where $|k|$ is the Euclidean norm. $Z_d(N)$ increases to $Z_d:=\sum_{k\in\Z^d\setminus0}|k|^{-4}<\infty$ for $d\le3$; $Z_1=2\zeta(4)=\pi^4/45$, and (certified enclosures, Lemma~\ref{lem:epstein} and Section~\ref{sec:cert})
\begin{equation}\label{eq:Zvalues}
6.02681<Z_2<6.02682,\qquad 16.53231<Z_3<16.53232 .
\end{equation}
For every $d$, $Z_d(N)\le2d\,3^{d-1}\sum_{j=1}^{N-1}j^{d-5}$.
\end{lemma}
\begin{proof}
By \eqref{eq:ccsums}, $\Lambda_1^2K_2=\sum_{k\ne0}\prod_i\varepsilon_{k_i}\cdot\prod_i\cos^2(k_iu)\,(\Lambda_1/\lambda_k)^2$ and $\lambda_k/\Lambda_1=\sum_i\sin^2(k_iu)/\sin^2u$. For an integer $1\le j\le N-1$ we have $ju\in(0,\frac\pi2)$ and, by Lemma~\ref{lem:elem}(a),(b), $\sin^2(ju)/\sin^2u\ge\sin^2(ju)/u^2\ge j^2\cos(ju)$. Hence
\[\lambda_k/\Lambda_1\ge\sum_ik_i^2\cos(k_iu)\ge|k|^2\prod_i\cos(k_iu),\quad\text{i.e.}\quad\prod_i\cos^2(k_iu)(\Lambda_1/\lambda_k)^2\le|k|^{-4}.\] Since $\prod_i\varepsilon_{k_i}$ is the number of points of $\Z^d$ obtained from $k$ by sign changes, the upper bound follows. For the lower bound keep the terms $k=e_i$ ($\lambda_k=\Lambda_1$, total weight $W$) and $k=e_i+e_j$, $i<j$ ($\lambda_k=2\Lambda_1$, weight $4\cos^4u$ each).
The number of $k\in\Z^d$ with $|k|_\infty=j$ is $(2j+1)^d-(2j-1)^d\le2d(2j+1)^{d-1}\le2d\,3^{d-1}j^{d-1}$, and $|k|\ge j$ for them; this gives the last bound and the convergence for $d\le3$. An elementary enclosure, $6.02679<Z_2<6.02682$ and $16.4629<Z_3<16.5484$ (script \texttt{06}), is obtained from the partial sums over $|k|_\infty\le R$ and the tail bound $\sum_{|k|_\infty>R}|k|^{-4}\le(1+\frac{\sqrt d}{2(R+1)})^4\int_{|x|\ge R+\frac12}|x|^{-4}dx$ (compare each term with the integral over the unit cube centred at $k$: for $x$ in that cube $|x|\le|k|+\frac{\sqrt d}2\le|k|(1+\frac{\sqrt d}{2(R+1)})$, and the cubes are disjoint and lie outside the ball of radius $R+\frac12$); the integral equals $\pi/(R+\frac12)^2$ for $d=2$ and $4\pi/(R+\frac12)$ for $d=3$. The sharp enclosures \eqref{eq:Zvalues} are proved in Lemma~\ref{lem:epstein}.
\end{proof}
\begin{remark}
$Z_2=4\zeta(2)G=6.0268120\ldots$ ($G$ Catalan's constant) by the Lorenz--Hardy formula $\sum_{(m,n)\ne0}(m^2+n^2)^{-s}=4\zeta(s)\beta(s)$ (see e.g.\ I.~J.~Zucker and M.~M.~Robertson, J.~Phys.~A 8 (1975) 874, or J.~M.~Borwein et al., \emph{Lattice Sums Then and Now}, Cambridge Univ.\ Press 2013, Section~1.3); the closed form is not used, only \eqref{eq:Zvalues} (the certified value of Lemma~\ref{lem:epstein} agrees with it to $39$ digits). Numerically $\Lambda_1^2K_2\to Z_d$ for $d=2,3$ (the ratio is $0.998$ at $d=2$, $N=60$), so \eqref{eq:K2} is asymptotically sharp.
\end{remark}

\begin{lemma}[Epstein zeta values by theta splitting]\label{lem:epstein}
For $d\ge1$ and real $s>d/2$ let $Z_d(s)=\sum_{k\in\Z^d\setminus0}|k|^{-2s}$ (so $Z_d=Z_d(2)$), and for $c>0$, $\alpha\in\R$ let $E_\alpha(c)=\int_1^\infty t^{\alpha-1}e^{-ct}\,dt=c^{-\alpha}\Gamma(\alpha,c)$ (incomplete gamma function). Then
\begin{equation}\label{eq:epstein}
\pi^{-s}\Gamma(s)\,Z_d(s)=\frac{1}{s-\frac d2}-\frac1s+\sum_{k\in\Z^d\setminus0}\Big[E_s\big(\pi|k|^2\big)+E_{\frac d2-s}\big(\pi|k|^2\big)\Big],
\end{equation}
where all terms of the series are positive and, for $s\ge1$ and every integer $R\ge1$ with $\pi R^2>s-1$,
\begin{equation}\label{eq:epsteintail}
\sum_{|k|_\infty>R}\Big[E_s\big(\pi|k|^2\big)+E_{\frac d2-s}\big(\pi|k|^2\big)\Big]\le\frac{4d\,(1.1)^{d-1}\cdot1.001\;e^{-\pi(R+1)^2}}{\pi R^2-s+1}.
\end{equation}
Moreover, for $Z_d'=\sum_{|k|^2\ge2}\big(|k|^2(|k|^2-1)\big)^{-1}$ ($d\le3$) and every integer $J\ge2$,
\begin{equation}\label{eq:Zprime}
Z_d'=\sum_{j=2}^{J}\big(Z_d(j)-2d\big)+\sum_{|k|^2\ge2}\frac{|k|^{-2J}}{|k|^2-1},\qquad \sum_{|k|_\infty>R}\frac{|k|^{-2J}}{|k|^2-1}\le\frac{4d\,3^{d-1}}{2J+2-d}\,R^{\,d-2-2J}.
\end{equation}
Evaluating \eqref{eq:epstein}--\eqref{eq:Zprime} in ball arithmetic ($R=7$; $J=8$, $R=40$; script \texttt{14\_certify\_epstein.py}) gives
\begin{equation}\label{eq:Zcert}
\begin{aligned}Z_2&=6.02681203969194012\ldots,&Z_3&=16.5323159597616696\ldots,\\ Z_2'&=3.22591076832683889\ldots,&Z_3'&=14.7022972380072310\ldots\end{aligned}
\end{equation}
Here the decimal expansions are \emph{truncated}, not rounded, and every digit shown is certified: each value lies strictly between the printed number and the printed number plus one unit in its last place (the script asserts exactly these inequalities).
\end{lemma}
\begin{proof}
For $c>0$, $\Gamma(s)c^{-s}=\int_0^\infty t^{s-1}e^{-ct}dt$. With $c=\pi|k|^2$, summing over $k\ne0$ (Tonelli),
$\pi^{-s}\Gamma(s)Z_d(s)=\int_0^\infty t^{s-1}(\Theta(t)-1)\,dt$, where $\Theta(t)=\theta(t)^d$ and $\theta(t)=\sum_{j\in\Z}e^{-\pi j^2t}$. Jacobi's transformation $\theta(1/t)=\sqrt t\,\theta(t)$ ($t>0$; DLMF 20.7.32 with $z=0$, $\tau=it$, or Poisson summation) gives $\Theta(1/t)=t^{d/2}\Theta(t)$. Substituting $t\mapsto1/t$ in the part $\int_0^1$,
\[
\int_0^1t^{s-1}(\Theta(t)-1)\,dt=\int_1^\infty t^{-s-1}\big(t^{d/2}\Theta(t)-1\big)dt=\int_1^\infty t^{\frac d2-s-1}(\Theta(t)-1)\,dt+\int_1^\infty\big(t^{\frac d2-s-1}-t^{-s-1}\big)dt,
\]
and the last integral equals $\frac{1}{s-d/2}-\frac1s$ because $s>\frac d2$. Since $\Theta(t)-1=\sum_{k\ne0}e^{-\pi|k|^2t}$, Tonelli's theorem gives \eqref{eq:epstein}.
For the tail: $\frac d2-s-1\le s-1$, so each bracket is at most $2E_s(c)$, and for $c>s-1\ge0$, $E_s(c)=e^{-c}\int_0^\infty(1+x)^{s-1}e^{-cx}dx\le e^{-c}\int_0^\infty e^{-(c-s+1)x}dx=e^{-c}/(c-s+1)$. For $|k|_\infty>R$ we have $c=\pi|k|^2>\pi R^2$, hence the left side of \eqref{eq:epsteintail} is at most $\frac{2}{\pi R^2-s+1}\sum_{|k|_\infty>R}e^{-\pi|k|^2}$, and
$\sum_{|k|_\infty>R}e^{-\pi|k|^2}\le d\,\theta(1)^{d-1}\cdot2\sum_{j\ge R+1}e^{-\pi j^2}\le d\,(1.1)^{d-1}\cdot\frac{2e^{-\pi(R+1)^2}}{1-e^{-2\pi(R+1)}}$,
using $\theta(1)=1+2\sum_{j\ge1}e^{-\pi j^2}<1.0865$, $j^2-(R+1)^2\ge2(R+1)(j-R-1)$ and $(1-e^{-4\pi})^{-1}<1.001$.
For \eqref{eq:Zprime}: $\frac1{n(n-1)}=\sum_{j=2}^{J}n^{-j}+\frac{n^{-J}}{n-1}$ for $n\ge2$, and exactly $2d$ points of $\Z^d$ have $|k|^2=1$. For $|k|_\infty=j>R$: $|k|^2-1\ge\frac12|k|^2$ and $|k|\ge j$, and there are at most $2d\,3^{d-1}j^{d-1}$ such points (proof of Lemma~\ref{lem:K2}); so the tail is at most $4d\,3^{d-1}\sum_{j>R}j^{d-3-2J}\le4d\,3^{d-1}\int_R^\infty x^{d-3-2J}dx$.
The script evaluates the finitely many remaining terms with $300$-bit balls ($\Gamma(\alpha,c)$ by Arb), adds the tail bounds, and checks the result against (i) the closed form $Z_2=4\zeta(2)G$, (ii) direct summation of $Z_d(4)$, $Z_d(6)$, (iii) a second choice of $R$ and $J$.
\end{proof}

\begin{lemma}[size of $\tau$ in general dimension]\label{lem:taugen}
For every $d\ge1$, $N\ge2$,
\begin{equation}\label{eq:taulow}
\tau\ \ge\ 2n\Big(1-\frac1n\Big)^2,\qquad\text{hence}\qquad \Xb=\Lambda_1\tau\ \ge\ \frac4d\,N^{d-2}\big(1-N^{-d}\big)^2 .
\end{equation}
For $d\ge3$ also $\tau\le\dfrac{d^2\,3^{d-1}}{d-2}\,N^d$, so that $\Xb$ is of exact order $N^{d-2}$.
\end{lemma}
\begin{proof}
Let $v=e_a-\frac1n\one\perp\one$. Then $\ip v{L^+v}=\ip{e_a}{L^+e_a}=\tau/n$, $\ip vv=1-\frac1n$ and $\ip v{Lv}=L_{aa}=\frac{d}{2d}=\frac12$ (a corner has $d$ neighbours). On $\one^\perp$, $L$ is positive definite with inverse $L^+$, so by Cauchy--Schwarz $\ip vv^2=\ip{(L^+)^{1/2}v}{L^{1/2}v}^2\le\ip v{L^+v}\ip v{Lv}$, which is the first claim; the second uses $\Lambda_1=\frac2d\sin^2u\ge\frac{2}{dN^2}$ (Lemma~\ref{lem:elem}(a)). For the upper bound, $\lambda_k=\frac2d\sum_i\sin^2(k_iu)\ge\frac{2|k|^2}{dN^2}$ and $w_k\le\prod_i\varepsilon_{k_i}$ give $\tau\le\frac{dN^2}2\sum_{0<|k|_\infty<N}|k|^{-2}\le\frac{dN^2}{2}\,2d\,3^{d-1}\sum_{j=1}^{N-1}j^{d-3}\le\frac{d^23^{d-1}}{d-2}N^d$.
\end{proof}

\begin{lemma}[monotone coupling]\label{lem:coupling}
Let $d\ge1$, $N\ge2$, target $a=(N,\dots,N)$, and order $\Omega$ coordinatewise: $x\preceq y$ iff $x_i\le y_i$ for all $i$. If $x\preceq y$ then, in continuous and in discrete time and for every activity $q\in(0,1]$,
\[
\PP_x(T>t)\ \ge\ \PP_y(T>t)\quad\text{for all }t\ge0 .
\]
Consequently $m(x)\ge m(y)$ and $r_0(x)\ge r_0(y)$ (with $m(a)=r_0(a)=0$): the mean first-passage time and the principal residue are coordinatewise non-increasing, and both are maximal at the opposite corner $x_0=(1,\dots,1)$.
\end{lemma}
\begin{proof}
Run two copies $(X_s)$, $(Y_s)$ of the reflecting walk \emph{without} target from $x$ and $y$, driven by the same jump times and the same proposed directions $\pm e_i$. A proposed move $+e_i$ maps the $i$-th coordinates to $\min(X^i+1,N)$, $\min(Y^i+1,N)$, a proposed move $-e_i$ to $\max(X^i-1,1)$, $\max(Y^i-1,1)$ (cancelled moves are exactly the cases where $\min$/$\max$ is active), and the other coordinates are unchanged. Both maps are non-decreasing, so $X_s\preceq Y_s$ for all $s$. At the hitting time $T^X$ of $a$ by $X$ we have $a=X_{T^X}\preceq Y_{T^X}$, hence $Y_{T^X}=a$ because $a$ is the maximum of $\Omega$; thus $T^Y\le T^X$ pathwise, which proves the inequality. Integrating over $t$ gives $m(x)\ge m(y)$ (Lemma~\ref{lem:taum}). By \eqref{eq:cont} and $\sigma_0<\sigma_j$ ($j\ge1$), $r_0(x)=\lim_{t\to\infty}e^{\sigma_0t}\PP_x(T>t)$ for the unit-rate walk, so $r_0(x)\ge r_0(y)$.
\end{proof}

\begin{theorem}[$r_0>1$ and $\sigma_0\cdot\mathrm{MFPT}>1$, all $d$ and $N$]\label{thm:r0gt1}
For the reflecting box with corner target and start at the opposite corner, for every $d\ge1$ and $N\ge2$ with $(d,N)\ne(1,2)$:
\[
r_0=\max_{x\in\Omega'}r_0(x)>1,\qquad \sigma_0\,m=\max_{x\in\Omega'}\sigma_0\,m(x)>1 ,
\]
and quantitatively $\ \sigma_0m\ \ge\ \dfrac{\Xb+\mu_-}{\Xb+1}=1+\dfrac{\mu_--1}{\Xb+1}$ with $\mu_-=\mu_-(N)\ge1$ from \eqref{eq:muminus}. (For $(d,N)=(1,2)$: $r_0=\sigma_0m=1$.)
\end{theorem}
\begin{proof}
By Proposition~\ref{prop:qsd} the $\psi_0$-weighted averages over $\Omega'$ of $r_0(\cdot)$ and of $\sigma_0m(\cdot)$ are equal to $1$, with $\psi_0>0$ (Theorem~\ref{thm:residues}(c); $\Omega'$ is connected). By Lemma~\ref{lem:coupling} both functions attain their maximum at $x_0$, so $r_0(x_0)\ge1$ and $\sigma_0m(x_0)\ge1$, with equality only if the function is constant on $\Omega'$. By Proposition~\ref{prop:qsd} this happens iff $\one_{\Omega'}$ is an eigenvector of $L_a$, i.e.\ iff $-L_{ya}$ is the same for all $y\in\Omega'$, i.e.\ iff every site of $\Omega'$ is a neighbour of $a$. The corner $a$ has $d$ neighbours and $|\Omega'|=N^d-1$, and $N^d-1=d$ only for $(d,N)=(1,2)$. The quantitative bound follows from $\sigma_0m=(1+\mu)\,\sigma_0\tau$, $\sigma_0\tau\ge\Xb/(1+\Xb)$ \eqref{eq:s0basic} and $\mu\ge\mu_-/\Xb$ \eqref{eq:muminus}, all valid for every $d\ge1$; it gives a second proof of $\sigma_0m\ge1$ that does not use the coupling.
\end{proof}

\begin{theorem}[corner to corner, explicit in $\Xb$]\label{thm:cc}
Let $d\ge2$, $N\ge2$, and for the corner-to-corner geometry let $\Xb=\Lambda_1\tau$, $\mu=(m-\tau)/\tau$,
\[
W=2d\cos^2\frac{\pi}{2N},\quad \ell_W=1+\ln W,\quad\vartheta_0=\binom d2\cos^4\frac\pi{2N},\quad\mu_-=\mu_-(N)=\frac23(N^2-1)\sin^2\frac{\pi}{2N}\ (\ge1),
\]
and let $Z$ be any number with $\Lambda_1^2K_2\le Z$ (for instance $Z=Z_d(N)$, or $Z=Z_d$ if $d\le3$). Assume $\Xb>1$ and put $\beta=\Xb/(\Xb-1)$. Then:
\begin{enumerate}
\item[(i)] \emph{(principal eigenvalue)}
\[
1-\frac{\beta Z}{\Xb^2}\ \le\ \frac{2}{1+\sqrt{1+4\beta Z/\Xb^2}}\ \le\ \sigma_0\tau\ \le\ \frac{2}{1+\sqrt{1+4W/\Xb^2}}\ \le\ 1-\frac{W}{\Xb^2}+\frac{2W^2}{\Xb^4}.
\]
\item[(ii)] \emph{(principal eigenvalue times MFPT)} $\mu_-/\Xb\le\mu\le\ell_W/\Xb$ and $\sigma_0m=(1+\mu)\,\sigma_0\tau$; in particular
\begin{gather*}
\Big(1+\frac{\mu_-}{\Xb}\Big)\Big(1-\frac{\beta Z}{\Xb^2}\Big)\ \le\ \sigma_0\,m\ <\ 1+\frac{\ell_W}{\Xb};\\
\text{without any assumption on }\Xb:\qquad \frac{\Xb+\mu_-}{1+\Xb}\le\sigma_0\,m<1+\frac{\ell_W}{\Xb}.
\end{gather*}
Since $\mu_-\ge1$, the last lower bound is $\ge1$: \emph{always} $\ \tau\le\dfrac1{\sigma_0}\le\tau+\dfrac1{\Lambda_1}\le m$ (see also Theorem~\ref{thm:r0gt1}).
\item[(iii)] \emph{(second pole)} $\sigma_1=\Lambda_1(1+\delta)$ with $0<\delta<1$. If $\Xb>W+1-\vartheta_0$ then
$\delta\le\delta_+:=\dfrac{W}{\Xb-W-1+\vartheta_0}$. If moreover $\delta_+<1$, put $\omega=\dfrac{2(1+\delta_+)}{1-\delta_+}$; then
$\delta\ge\delta_-:=\dfrac{W}{\Xb-W+\omega(Z-W)-\frac1{1+\delta_+}}$.
\item[(iv)] \emph{(gap)} $\Lambda_1\big(1-\frac1{\Xb}\big)<\sigma_1-\sigma_0<2\Lambda_1$, and under the hypotheses of (iii)
\[
\Lambda_1\Big(1+\delta_--\frac1{\Xb}\Big)\le\sigma_1-\sigma_0\le\Lambda_1\Big(1+\delta_+-\frac{1}{1+\Xb}\Big).
\]
\item[(v)] \emph{(residue of the principal pole)} $r_0>0$ and
\begin{gather*}
\Big(1+\frac{\mu_-}{\Xb}\Big)\Big(1-\frac{\beta(1+\beta)Z}{\Xb^2}\Big)\ \le\ (1+\mu)\Big(1-\frac{\beta(1+\beta)Z}{\Xb^2}\Big)\ \le\ r_0,\\
r_0\ \le\ 1+\mu+\frac{\beta Z}{\Xb^2}\ \le\ 1+\frac{\ell_W}{\Xb}+\frac{\beta Z}{\Xb^2}
\end{gather*}
(the first inequality provided $\beta(1+\beta)Z\le\Xb^2$). Moreover $r_0>1$ without any hypothesis on $\Xb$ (Theorem~\ref{thm:r0gt1}).
\item[(vi)] \emph{(residue of the second pole)} $r_1<0$. Under the hypotheses of (iii) let
\[
y_-=W-\ell_W-\omega(Z-W)-1,\qquad y_+=W-\mu_-+\omega(Z-W)-\tfrac{1}{1+\delta_+},
\]
\[
\mathcal N_-=W+\min(\delta_-y_-,\delta_+y_-),\qquad\mathcal N_+=W+\max(\delta_-y_+,\delta_+y_+),
\]
\[
\mathcal D_-=(1+\delta_-)W+\delta_-^2\Big[\tfrac1{1+\delta_+}+(1+\delta_-)\vartheta_0\Big],\qquad
\mathcal D_+=(1+\delta_+)W+\delta_+^2\Big[1+\tfrac{\omega^2(Z-W)}{1+\delta_+}\Big].
\]
Then $-r_1\le\delta_+\mathcal N_+/\mathcal D_-$, and if $\mathcal N_->0$ also $-r_1\ge\delta_-\mathcal N_-/\mathcal D_+$.
\end{enumerate}
\end{theorem}
\begin{proof}
By Lemma~\ref{lem:K2}, $W/\Xb^2\le\kappa=\Lambda_1^2K_2/\Xb^2\le Z/\Xb^2$ and $\vartheta_0\le\vartheta=\Lambda_1^2K_2-W\le Z-W$; by Lemma~\ref{lem:arrival}, $\mu_-/\Xb\le\mu\le\ell_W/\Xb$ with $\mu_-\ge1$; and \eqref{eq:ccdata} holds.
(i) is Theorem~\ref{thm:sigma0} with Remark~\ref{rem:monotone}(a).
(ii) $\sigma_0m=\sigma_0\tau(1+\mu)$ (Corollary~\ref{cor:s0m}); combine (i), $\sigma_0\tau\le1$ and \eqref{eq:s0basic}: $\sigma_0m\ge(1+\frac{\mu_-}{\Xb})\,\sigma_0\tau$ with $\sigma_0\tau\ge\max\{0,1-\beta Z/\Xb^2\}$, respectively $\sigma_0\tau\ge\Xb/(1+\Xb)$; the upper inequality is strict because $\sigma_0\tau<1$ (the first inequality in the proof of Theorem~\ref{thm:sigma0} is strict since $K_2>0$). Finally $(\Xb+\mu_-)/(1+\Xb)\ge1$ because $\mu_-\ge1$, i.e.\ $\frac1{\sigma_0}\le\tau+\frac1{\Lambda_1}\le\tau+\frac{\mu_-}{\Lambda_1}\le m$ by \eqref{eq:s0basic} and \eqref{eq:muminus}.
(iii) $\delta<1$ is \eqref{eq:W}. The rest is Theorem~\ref{thm:sigma1} in the form of Remark~\ref{rem:monotone}(b) with $Y_0=\Xb-W$, $\vartheta_-=\vartheta_0$, $\vartheta_+=Z-W$, $\varrho=\frac12$, for which $\hat\omega=(1+\delta_+)/(1-\frac{1+\delta_+}2)=\omega$.
(iv) Corollary~\ref{cor:gap}.
(v) $r_0>0$ is Corollary~\ref{cor:r1cc}. By \eqref{eq:r0} with $\kappa_{x_0}=\kappa$, $t=\sigma_0\tau\le1$: $r_0\le1+t\mu+\varepsilon_1\le1+\mu+\beta\kappa\le1+\mu+\beta Z/\Xb^2$. For the lower bound let $\kappa'=Z/\Xb^2$ with $\beta\kappa'\le1$. By (i), $1-t\le\beta\kappa'$, so the numerator of \eqref{eq:r0} is at least $1+\mu-\mu\beta\kappa'-\beta\kappa'=(1+\mu)(1-\beta\kappa')\ge0$, and $1/(1+\varepsilon_2)\ge1-\beta^2\kappa'$; the product of $(1+\mu)(1-\beta\kappa')$ with a number in $[1-\beta^2\kappa',1]\cap(0,1]$ is at least $(1+\mu)(1-\beta\kappa'-\beta^2\kappa')$. (If $\beta\kappa'>1$ this lower bound is negative and holds trivially because $r_0>0$; so the middle inequality needs no hypothesis.) If $1-\beta(1+\beta)\kappa'\ge0$ we may further replace $1+\mu$ by its lower bound $1+\mu_-/\Xb$. The last upper bound uses $\mu\le\ell_W/\Xb$.
(vi) $r_1<0$ is Corollary~\ref{cor:r1cc}. The bounds are Theorem~\ref{thm:r1size} in the form of Remark~\ref{rem:monotone}(c), with $A_{x_0}=W$ and $Y_{x_0}(0)=W-\Lambda_1(m-\tau)\in[W-\ell_W,W-\mu_-]$ by \eqref{eq:mtau}, \eqref{eq:muminus}.
\end{proof}

\begin{corollary}[all dimensions $d\ge3$]\label{cor:alld}
For $d\ge3$ and $N\ge2$ with $\Xb>1$: $\ 0<1-\sigma_0\tau\le\beta Z_d(N)/\Xb^2$ and $0<\sigma_0m-1<(1+\ln 2d)/\Xb$, where $\Xb\ge\frac4dN^{d-2}(1-N^{-d})^2$ and $Z_d(N)\le Z_3$ ($d=3$), $Z_4(N)\le216\,(1+\ln N)$, $Z_d(N)\le\frac{2d\,3^{d-1}}{d-4}N^{d-4}$ ($d\ge5$). In particular $\sigma_0\tau=1-O(N^{-2})$ for $d=3$, $1-O(N^{-4}\ln N)$ for $d=4$ and $1-O(N^{-d})$ for $d\ge5$, with explicit constants.
\end{corollary}
\begin{proof}
Theorem~\ref{thm:cc}(i),(ii), Theorem~\ref{thm:r0gt1}, Lemma~\ref{lem:taugen} and Lemma~\ref{lem:K2} ($\sum_{j=1}^{N-1}j^{-1}\le1+\ln N$, $\sum_{j=1}^{N-1}j^{d-5}\le N^{d-4}/(d-4)$ for $d\ge5$).
\end{proof}

\section{Exact reductions of the lattice Green's function at the corner}\label{sec:reductions}

The functionals $\tau$ and $m$ are $d$-fold sums \eqref{eq:ccsums}. One of the $d$ summations can always be carried out in closed form. Throughout, $\varepsilon'_0=\frac12$, $\varepsilon'_k=1$ ($k\ge1$), so that $w_k=\prod_i\varepsilon'_{k_i}(1+c_{k_i})$.

\begin{lemma}[one summation in closed form]\label{lem:cycle}
Let $N\ge1$ and $\varphi>0$. Then
\begin{align}
\sum_{k=0}^{N-1}\varepsilon'_k\,\frac{1+c_k}{\cosh\varphi-c_k}&=N\Big[\coth\frac\varphi2\,\coth(N\varphi)-1\Big],\label{eq:sum1}\\
\sum_{k=0}^{N-1}\varepsilon'_k\,(-1)^k\frac{1+c_k}{\cosh\varphi-c_k}&=\frac{N\coth\frac\varphi2}{\sinh(N\varphi)},\label{eq:sum2}\\
\frac1\pi\int_0^\pi\frac{1+\cos\theta}{\cosh\varphi-\cos\theta}\,d\theta&=\coth\frac\varphi2-1 .\label{eq:int1}
\end{align}
\end{lemma}
\begin{proof}
For $r=e^{-\varphi}\in(0,1)$ the Poisson kernel identity $\sum_{m\in\Z}r^{|m|}e^{im\theta}=\frac{1-r^2}{1-2r\cos\theta+r^2}$ reads
\begin{equation}\label{eq:poisson}
\frac{1}{\cosh\varphi-\cos\theta}=\frac{1}{\sinh\varphi}\sum_{m\in\Z}e^{-|m|\varphi}e^{im\theta},
\end{equation}
the series converging absolutely and uniformly in $\theta$. Let $\theta_l=\pi l/N$, $l=0,\dots,2N-1$, and let $j$ be an integer with $0\le j\le2N$. Since $\sum_{l=0}^{2N-1}e^{im\theta_l}$ equals $2N$ if $2N\mid m$ and $0$ otherwise,
\[
\begin{aligned}\frac{1}{2N}\sum_{l=0}^{2N-1}\frac{\cos(j\theta_l)}{\cosh\varphi-\cos\theta_l}&=\frac{1}{\sinh\varphi}\sum_{i\in\Z}e^{-|2Ni+j|\varphi}
=\frac{1}{\sinh\varphi}\Big[e^{-j\varphi}+\frac{(e^{-j\varphi}+e^{j\varphi})e^{-2N\varphi}}{1-e^{-2N\varphi}}\Big]\\&=\frac{\cosh((N-j)\varphi)}{\sinh\varphi\,\sinh(N\varphi)}=:\gamma_j\end{aligned}
\]
(we used $|2Ni+j|=2Ni+j$ for $i\ge0$ and $=2N|i|-j$ for $i\le-1$, as $0\le j\le2N$). The summand is invariant under $l\mapsto2N-l$ and $1+\cos\theta_N=0$, so for any function $F$,
$\sum_{k=0}^{N-1}\varepsilon'_k(1+c_k)F(c_k)=\frac12\sum_{l=0}^{2N-1}(1+\cos\theta_l)F(\cos\theta_l)$. With $F(c)=1/(\cosh\varphi-c)$ this gives for the left side of \eqref{eq:sum1} the value $N(\gamma_0+\gamma_1)=N\frac{\cosh(N\varphi)+\cosh((N-1)\varphi)}{\sinh\varphi\sinh N\varphi}=N\frac{\cosh((N-\frac12)\varphi)}{\sinh\frac\varphi2\,\sinh N\varphi}$, which is the right side because $\cosh((N-\frac12)\varphi)=\cosh N\varphi\cosh\frac\varphi2-\sinh N\varphi\sinh\frac\varphi2$. For \eqref{eq:sum2} use $(-1)^l=\cos(N\theta_l)$ and $\cos(N\theta)\cos\theta=\frac12[\cos((N+1)\theta)+\cos((N-1)\theta)]$: the left side equals $N[\gamma_N+\frac12(\gamma_{N-1}+\gamma_{N+1})]=N\frac{1+\cosh\varphi}{\sinh\varphi\sinh N\varphi}=N\coth\frac\varphi2/\sinh(N\varphi)$. Finally, the integrand of \eqref{eq:int1} is even, so the left side is $\frac1{2\pi}\int_{-\pi}^{\pi}$; inserting \eqref{eq:poisson} and integrating term by term, only $m=0$ and $m=\pm1$ contribute and give $(1+e^{-\varphi})/\sinh\varphi=2/(e^\varphi-1)=\coth\frac\varphi2-1$.
\end{proof}

\begin{lemma}[cotangent sums]\label{lem:cot}
For $N\ge2$: $\ \displaystyle\sum_{k=1}^{N-1}\cot^2\frac{\pi k}{2N}=\frac{(N-1)(2N-1)}{3}$, $\ \displaystyle\sum_{k=1}^{N-1}(-1)^k\cot^2\frac{\pi k}{2N}=-\frac{N^2-1}{3}$, $\ \displaystyle\sum_{k=1}^{N-1}\cos^2\frac{\pi k}{2N}=\frac{N-1}2$, $\ \displaystyle\prod_{k=1}^{N-1}\sin\frac{\pi k}{2N}=\frac{\sqrt N}{2^{N-1}}$.
\end{lemma}
\begin{proof}
$(1+c_k)/(1-c_k)=\cot^2(\pi k/2N)$. Subtract the $k=0$ term $1/(\cosh\varphi-1)$ from \eqref{eq:sum1} and from \eqref{eq:sum2} and let $\varphi\downarrow0$. With $\coth x=\frac1x+\frac x3+O(x^3)$, $1/\sinh x=\frac1x-\frac x6+O(x^3)$ and $1/(\cosh\varphi-1)=\frac{2}{\varphi^2}-\frac16+O(\varphi^2)$:
$N[\coth\frac\varphi2\coth N\varphi-1]=\frac2{\varphi^2}+\frac{2N^2}3+\frac16-N+O(\varphi^2)$ and $N\coth\frac\varphi2/\sinh N\varphi=\frac{2}{\varphi^2}-\frac{N^2}{3}+\frac16+O(\varphi^2)$, which gives the two cotangent sums. The third identity follows from $\cos^2\frac{\pi k}{2N}+\cos^2\frac{\pi(N-k)}{2N}=1$. The last one follows from the classical $\prod_{k=1}^{M-1}\sin\frac{\pi k}{M}=M/2^{M-1}$ with $M=2N$ and the symmetry $k\mapsto2N-k$.
\end{proof}

\begin{proposition}[$d=2$: single sums]\label{prop:d2exact}
Let $d=2$, $x_k=\frac{\pi k}{2N}$, $y_k=2N\operatorname{arsinh}(\sin x_k)$ and
\begin{equation}\label{eq:Phi}
\Phi(x)=\frac{\cos^2x\,\sqrt{1+\sin^2x}}{\sin x}\qquad(0<x\le\tfrac\pi2).
\end{equation}
Then for every $N\ge2$
\begin{equation}\label{eq:tau2exact}
\tau_2=4N\sum_{k=1}^{N-1}\Phi(x_k)\coth y_k-\frac23(N^2-1),\qquad
m_2=4N\sum_{k=1}^{N-1}\Phi(x_k)\Big[\coth y_k-\frac{(-1)^k}{\sinh y_k}\Big].
\end{equation}
\end{proposition}
\begin{proof}
Here $\lambda_k=\frac12(2-c_{k_1}-c_{k_2})$ and $\tau_2=2\sum_{k\ne0}\varepsilon'_{k_1}\varepsilon'_{k_2}\frac{(1+c_{k_1})(1+c_{k_2})}{2-c_{k_1}-c_{k_2}}$. For $k_1\ge1$ define $\varphi>0$ by $\cosh\varphi=2-c_{k_1}$; then $\sinh^2\frac\varphi2=\frac{1-c_{k_1}}{2}=\sin^2x_{k_1}$, so $\varphi=2\operatorname{arsinh}(\sin x_{k_1})$, $N\varphi=y_{k_1}$, $\coth\frac\varphi2=\sqrt{1+\sin^2x_{k_1}}/\sin x_{k_1}$, and $1+c_{k_1}=2\cos^2x_{k_1}$. By \eqref{eq:sum1} the sum over $k_2$ is $N[\coth\frac\varphi2\coth y_{k_1}-1]$. For $k_1=0$ the sum over $k_2\ge1$ is $\frac12\cdot2\sum_{k_2\ge1}\cot^2x_{k_2}=\frac{(N-1)(2N-1)}3$. Hence
$\tau_2=\frac23(N-1)(2N-1)+4N\sum_{k\ge1}\Phi(x_k)\coth y_k-4N\sum_{k\ge1}\cos^2x_k$, and Lemma~\ref{lem:cot} gives the first formula. In the same way, by \eqref{eq:sum2}, $\tau_2-m_2=\sum_{k\ne0}(-1)^{|k|}w_k/\lambda_k=-\frac23(N^2-1)+4N\sum_{k\ge1}(-1)^k\Phi(x_k)/\sinh y_k$.
\end{proof}
\begin{remark}
The formula for $m_2$ is the single-sum MFPT identity proved in Theorem~4.2 of the article; it is equivalent, for even $N$, to the corner-to-corner resistance formula of J.~W.~Essam and F.~Y.~Wu, J.~Phys.~A 42 (2009) 025205. The formula for the stationary-start time $\tau_2$ is its analogue for the diagonal Green's function.
\end{remark}

\begin{proposition}[$d=3$: double sums]\label{prop:d3exact}
Let $d=3$. For $(\theta_1,\theta_2)\in[0,\pi]^2\setminus\{0\}$ define $\phi=\phi(\theta_1,\theta_2)>0$ by $\sinh^2\frac\phi2=\sin^2\frac{\theta_1}2+\sin^2\frac{\theta_2}2$ (equivalently $\cosh\phi=3-\cos\theta_1-\cos\theta_2$) and
\begin{equation}\label{eq:HP}
H=(1+\cos\theta_1)(1+\cos\theta_2)\Big(\coth\frac\phi2-1\Big),\qquad P_N=(1+\cos\theta_1)(1+\cos\theta_2)\coth\frac\phi2\,\big(\coth(N\phi)-1\big).
\end{equation}
Then, with $\sum'$ denoting the sum over $(k_1,k_2)\in\{0,\dots,N-1\}^2\setminus\{0\}$ and all functions evaluated at $(\theta_{k_1},\theta_{k_2})$,
\begin{align}
\tau_3&=(N-1)(2N-1)+3N\,{\sum}'\varepsilon'_{k_1}\varepsilon'_{k_2}\,H+3N\,{\sum}'\varepsilon'_{k_1}\varepsilon'_{k_2}\,P_N,\label{eq:tau3exact}\\
m_3-\tau_3&=(N^2-1)-3N\,{\sum}'\varepsilon'_{k_1}\varepsilon'_{k_2}(-1)^{k_1+k_2}\,\frac{(1+c_{k_1})(1+c_{k_2})\coth\frac\phi2}{\sinh(N\phi)}.\label{eq:m3exact}
\end{align}
\end{proposition}
\begin{proof}
$\tau_3=3\sum_{k\ne0}\prod_i\varepsilon'_{k_i}(1+c_{k_i})/(3-c_{k_1}-c_{k_2}-c_{k_3})$. For $(k_1,k_2)\ne0$ apply \eqref{eq:sum1} to the sum over $k_3$ with $\cosh\phi=3-c_{k_1}-c_{k_2}>1$: it equals $N[\coth\frac\phi2\coth N\phi-1]=N[(\coth\frac\phi2-1)+\coth\frac\phi2(\coth N\phi-1)]$. For $k_1=k_2=0$ the sum over $k_3\ge1$ is $\sum\cot^2x_{k_3}=\frac{(N-1)(2N-1)}3$. This is \eqref{eq:tau3exact}; \eqref{eq:m3exact} follows in the same way from \eqref{eq:sum2} and Lemma~\ref{lem:cot}, since $\tau_3-m_3=\sum_{k\ne0}(-1)^{|k|}w_k/\lambda_k$.
\end{proof}

\section{Two dimensions: logarithmic asymptotics with an explicit remainder}\label{sec:d2}

\subsection{Statement}
Let $\gamma$ be Euler's constant, $\kappa_2=2\operatorname{arsinh}1=\ln(3+2\sqrt2)$, and
\begin{equation}\label{eq:consts2}
I_\Psi=-\tfrac12+\tfrac32\ln2-\ln\pi,\qquad
\Lambda_*=\sum_{j\ge1}\frac{1}{j\,(e^{2\pi j}-1)},\qquad \Sigma_*=\sum_{j\ge1}\frac{(-1)^{j+1}}{j\,\sinh(\pi j)},
\end{equation}
\begin{equation}\label{eq:ctau}
c_\tau=\frac8\pi\big(\gamma+I_\Psi+2\Lambda_*\big)-\frac23=-0.727904612791689\ldots,\qquad
c_{m\tau}=\frac23+\frac8\pi\Sigma_*=0.882542400610606\ldots .
\end{equation}

\begin{theorem}[$d=2$, corner target]\label{thm:d2}
For every $N\ge2$ the unit-rate stationary-start mean hitting time $\tau_2$ of the corner and the corner-to-corner mean first-passage time $m_2$ of the $N\times N$ reflecting box satisfy
\begin{align}
\tau_2&=\frac8\pi N^2\ln N+c_\tau N^2+R_\tau(N),&-2.14&\le R_\tau(N)\le3.53,\label{eq:tau2}\\
m_2-\tau_2&=c_{m\tau}N^2+R_{m\tau}(N),&-9.06&\le R_{m\tau}(N)\le7.72 .\label{eq:m2}
\end{align}
Consequently $m_2=\frac8\pi N^2\ln N+C_2N^2+R_m(N)$ with $C_2=c_\tau+c_{m\tau}=0.154637787818917\ldots$ and $|R_m(N)|\le11.3$.
(For activity $q$ divide $\tau_2$ and $m_2$ by $q$.) The proof is analytic except for the numerical constants: the bound $\sup_{[0,\pi/2]}|\Psi''|\le2.35$ of Lemma~\ref{lem:Psi}(d) and the values of four elementary convergent series are certified by interval arithmetic (Section~\ref{sec:cert}).
\end{theorem}

\begin{lemma}[closed forms]\label{lem:closed}
$\Lambda_*=-\frac\pi{12}-\ln\Gamma(\tfrac14)+\ln2+\tfrac34\ln\pi$ and $\Sigma_*=\tfrac12\ln2-\tfrac\pi{12}$. Hence
\[
c_\tau=\frac8\pi\Big(\gamma-\frac12+\frac72\ln2+\frac12\ln\pi-2\ln\Gamma(\tfrac14)\Big)-2,\qquad c_{m\tau}=\frac{4\ln2}{\pi},
\]
and $C_2$ is the constant of Theorem~4.5 of the article.
\end{lemma}
\begin{proof}
Expanding $1/(e^{2\pi j}-1)=\sum_{l\ge1}e^{-2\pi jl}$ and $1/\sinh(\pi j)=2\sum_{l\ \mathrm{odd}}e^{-\pi jl}$ and summing over $j$ first (absolute convergence),
$\Lambda_*=-\ln\prod_{l\ge1}(1-e^{-2\pi l})$ and $\Sigma_*=2\ln\prod_{l\ \mathrm{odd}}(1+e^{-\pi l})$.
We use the following classical facts, with the notation and equation numbers of the NIST Digital Library of Mathematical Functions (DLMF, \texttt{dlmf.nist.gov}), at $\tau=i$, i.e.\ nome $q=e^{i\pi\tau}=e^{-\pi}$:
(20.5.3) $\theta_3(0,q)=\prod_{l\ge1}(1-q^{2l})(1+q^{2l-1})^2$;
(20.9.1)--(20.9.2) $k=\theta_2^2(0,q)/\theta_3^2(0,q)$, $K(k)=\frac\pi2\theta_3^2(0,q)$, $K'(k)=-i\tau K(k)$;
(23.5.5) $K(k)=K'(k)=\Gamma(\frac14)^2/(4\sqrt\pi)$ for $k^2=\frac12$;
(20.4.6) $\theta_1'(0,q)=\theta_2\theta_3\theta_4(0,q)$; (20.7.5) $\theta_3^4=\theta_2^4+\theta_4^4$;
(23.15.9), (23.17.8) $\eta(\tau)=(\frac12\theta_1'(0,q))^{1/3}=q^{1/12}\prod_{l\ge1}(1-q^{2l})$.
At $\tau=i$, $K'=K$, so $k^2=\frac12$, $\theta_3(0,e^{-\pi})^2=\frac2\pi K=\Gamma(\frac14)^2/(2\pi^{3/2})$ and $\theta_2^4=\frac12\theta_3^4=\theta_4^4$; hence $\eta(i)^3=\frac12\theta_2\theta_3\theta_4=\theta_3^3/(2\sqrt2)$, i.e.\ $\eta(i)=\theta_3(0,e^{-\pi})/\sqrt2=\Gamma(\frac14)/(2\pi^{3/4})$. Therefore
$\prod_{l}(1-e^{-2\pi l})=e^{\pi/12}\eta(i)$, which gives $\Lambda_*$, and $\prod_{l\ \mathrm{odd}}(1+e^{-\pi l})^2=\theta_3(0,e^{-\pi})/(e^{\pi/12}\eta(i))=\sqrt2\,e^{-\pi/12}$, which gives $\Sigma_*$.
\end{proof}
\begin{remark}
The inequalities \eqref{eq:tau2}, \eqref{eq:m2} do not depend on Lemma~\ref{lem:closed}: the constants are defined by the rapidly convergent series \eqref{eq:consts2}. In Section~\ref{sec:cert} the series and the closed forms are both enclosed by ball arithmetic and the enclosures overlap (to 30 digits).
\end{remark}

\subsection{Three auxiliary lemmas}

\begin{lemma}[harmonic numbers]\label{lem:harm}
For $N\ge2$: $H_{N-1}:=\sum_{k=1}^{N-1}\frac1k=\ln N+\gamma-\frac{1}{2N}-\eta_N$ with $\ \frac{1}{12(N+1)^2}<\eta_N<\frac1{12N^2}+\frac1{6N^3}\le\frac{1}{6N^2}$.
\end{lemma}
\begin{proof}
Let $a_M=H_M-\ln M-\frac1{2M}$. Then $a_{M+1}-a_M=\frac12\big(\frac1M+\frac1{M+1}\big)-\int_M^{M+1}\frac{dx}{x}$ is the error of the trapezoidal rule for the convex function $1/x$ on $[M,M+1]$, which equals $\frac1{12}f''(\xi)=\frac{1}{6\xi^3}$ for some $\xi\in(M,M+1)$. So $\frac{1}{6(M+1)^3}<a_{M+1}-a_M<\frac1{6M^3}$; $a_M$ increases to $\gamma$, and
$\gamma-a_N=\sum_{M\ge N}(a_{M+1}-a_M)\in\big(\frac16\sum_{M\ge N+1}M^{-3},\frac16\sum_{M\ge N}M^{-3}\big)\subset\big(\frac{1}{12(N+1)^2},\frac{1}{6N^3}+\frac{1}{12N^2}\big)$ by comparison with $\int x^{-3}dx$. Finally $H_{N-1}=H_N-\frac1N=\ln N+\gamma-\frac1{2N}-(\gamma-a_N)$, and $\frac{1}{12N^2}+\frac{1}{6N^3}\le\frac{1}{6N^2}$ for $N\ge2$.
\end{proof}

\begin{lemma}[the function $\Psi$]\label{lem:Psi}
Let $\Psi(x)=\Phi(x)-\frac1x$ on $(0,\frac\pi2]$, with $\Phi$ from \eqref{eq:Phi}. Then:
\begin{enumerate}
\item[(a)] $\Psi(x)=c(x)+Q(\sin x)$ with $c(x)=\dfrac{1}{\sin x}-\dfrac1x$ and $Q(s)=-\dfrac{s\,(1+s^2-s^4)}{1+(1-s^2)\sqrt{1+s^2}}$; $\Psi$ extends to a real-analytic function on a neighbourhood of $[0,\frac\pi2]$, with $\Psi(0)=0$ and $\Psi(\frac\pi2)=-\frac2\pi$.
\item[(b)] $-1.09\,x\le\Psi(x)\le0$ on $[0,\frac\pi2]$.
\item[(c)] $\int_0^{\pi/2}\Psi(x)\,dx=I_\Psi=-\frac12+\frac32\ln2-\ln\pi$.
\item[(d)] (computer-assisted) $-0.7382\le\Psi''(x)\le2.3125$ on $[0,\frac\pi2]$; in particular $|\Psi''|\le M_2:=2.35$.
\end{enumerate}
\end{lemma}
\begin{proof}
(a) With $s=\sin x$, $\Phi(x)=\frac{(1-s^2)\sqrt{1+s^2}}{s}$, and
$\Phi-\frac1s=\frac{(1-s^2)\sqrt{1+s^2}-1}{s}=\frac{(1-s^2)^2(1+s^2)-1}{s\,[(1-s^2)\sqrt{1+s^2}+1]}=Q(s)$ since $(1-s^2)^2(1+s^2)-1=-s^2(1+s^2-s^4)$. The denominator $D(s)=1+(1-s^2)\sqrt{1+s^2}$ is analytic and $\ge1$ on $[0,1]$, and $c$ is analytic for $|x|<\pi$ with $c(0)=0$.
(b) By Abramowitz and Stegun, \emph{Handbook of Mathematical Functions}, formula 4.3.68, $c(x)=\sum_{j\ge1}\frac{2(2^{2j-1}-1)|B_{2j}|}{(2j)!}x^{2j-1}=\frac x6+\frac{7x^3}{360}+\dots$ for $|x|<\pi$, with positive coefficients. Hence $c$, $c'$, $c''$ are non-negative and increasing on $[0,\pi)$ and $c(x)/x$ is increasing, so $\frac x6\le c(x)\le x\,c(\frac\pi2)/\frac\pi2=\frac2\pi(1-\frac2\pi)x$ on $[0,\frac\pi2]$. On $[0,1]$, $1\le D(s)\le2$ and $1\le1+s^2-s^4\le\frac54$, so $-\frac54s\le Q(s)\le-\frac s2$. Therefore $\Psi(x)\le\frac2\pi(1-\frac2\pi)x-\frac12\sin x\le\big[\frac2\pi(1-\frac2\pi)-\frac1\pi\big]x=\frac1\pi(1-\frac4\pi)x\le0$ (Lemma~\ref{lem:elem}(a)) and $\Psi(x)\ge\frac x6-\frac54x=-\frac{13}{12}x\ge-1.09x$.
(c) For $0<\epsilon<\frac\pi2$ substitute $s=\sin x$ and then $v=s^2$:
$\int_\epsilon^{\pi/2}\Phi\,dx=\int_{\sin\epsilon}^1\frac{\sqrt{1-s^4}}{s}ds=\frac12\int_{\sin^2\epsilon}^1\frac{\sqrt{1-v^2}}{v}dv$, and $\frac{d}{dv}\big[\sqrt{1-v^2}+\ln v-\ln(1+\sqrt{1-v^2})\big]=\frac{\sqrt{1-v^2}}{v}$. So $\int_\epsilon^{\pi/2}\Phi\,dx=\frac12\big[-\sqrt{1-\sin^4\epsilon}-2\ln\sin\epsilon+\ln(1+\sqrt{1-\sin^4\epsilon})\big]=-\frac12-\ln\epsilon+\frac12\ln2+o(1)$, while $\int_\epsilon^{\pi/2}\frac{dx}x=\ln\frac\pi2-\ln\epsilon$. Subtract and let $\epsilon\to0$.
(d) $\Psi''(x)=c''(x)+Q''(\sin x)\cos^2x-Q'(\sin x)\sin x$. The script \texttt{06\_certify\_constants.py} splits $[0,\frac\pi2]$ into $4000$ intervals $[\alpha,\beta]$; on each it bounds $c''$ by $c''(\alpha)\le c''\le c''(\beta)$ (monotonicity from (b); $c''(x)=\csc x(\csc^2x+\cot^2x)-2x^{-3}$ evaluated at the end points with $200$-bit balls, $c''(0)=0$) and evaluates the explicit rational-algebraic expressions for $Q'$, $Q''$ in ball arithmetic on the ball containing $\sin([\alpha,\beta])$. The union of the resulting enclosures is contained in $[-0.7382,2.3125]$.
\end{proof}

\begin{lemma}[trapezoidal sum]\label{lem:trap}
$\displaystyle\sum_{k=1}^{N-1}\Psi(x_k)=\frac{2N}{\pi}I_\Psi+\frac1\pi+e_1(N)$ with $|e_1(N)|\le\dfrac{\pi^2M_2}{48N}$, where $x_k=\frac{\pi k}{2N}$. Moreover
\begin{equation}\label{eq:e1limit}
\lim_{N\to\infty}N\,e_1(N)=\frac{\pi}{24}\Big(\Psi'(\tfrac\pi2)-\Psi'(0)\Big)=\frac{\pi}{24}\Big(\frac4{\pi^2}+\frac13\Big)=\frac1{6\pi}+\frac{\pi}{72}.
\end{equation}
\end{lemma}
\begin{proof}
Let $h=\frac\pi{2N}$ and $T_h=h\big[\frac{\Psi(0)}2+\sum_{k=1}^{N-1}\Psi(x_k)+\frac{\Psi(\pi/2)}2\big]$. On each cell $[x_{k-1},x_k]$ two integrations by parts give $\frac h2(\Psi(x_{k-1})+\Psi(x_k))-\int_{x_{k-1}}^{x_k}\Psi=\int_{x_{k-1}}^{x_k}q_h\Psi''$ with $q_h(x)=\frac12(x-x_{k-1})(x_k-x)\ge0$ and $\int_{x_{k-1}}^{x_k}q_h=\frac{h^3}{12}$. Summing over the $N$ cells, $T_h-\int_0^{\pi/2}\Psi=\int_0^{\pi/2}q_h\Psi''$, which is at most $\frac\pi2\cdot\frac{h^2}{12}\sup|\Psi''|$ in modulus. Since $\Psi(0)=0$ and $\Psi(\frac\pi2)=-\frac2\pi$, $T_h=h\big[\sum_k\Psi(x_k)-\frac1\pi\big]$, so $e_1(N)=h^{-1}\big(T_h-I_\Psi\big)$ and the bound follows. For the limit, $\Psi''$ is continuous on $[0,\frac\pi2]$ (Lemma~\ref{lem:Psi}(a)) and $q_h\ge0$, so by the mean value theorem for integrals $\int_{x_{k-1}}^{x_k}q_h\Psi''=\frac{h^3}{12}\Psi''(\xi_k)$ with some $\xi_k\in[x_{k-1},x_k]$; hence $h^{-2}(T_h-I_\Psi)=\frac1{12}\sum_{k=1}^{N}h\,\Psi''(\xi_k)$ is a Riemann sum and tends to $\frac1{12}\int_0^{\pi/2}\Psi''=\frac1{12}(\Psi'(\frac\pi2)-\Psi'(0))$. Thus $N e_1(N)=\frac\pi2\,h^{-2}(T_h-I_\Psi)\to\frac\pi{24}(\Psi'(\frac\pi2)-\Psi'(0))$. Finally, by Lemma~\ref{lem:Psi}(a), $\Psi'(x)=c'(x)+Q'(\sin x)\cos x$ with $c'(x)=-\frac{\cos x}{\sin^2x}+\frac1{x^2}$, $c'(0)=\frac16$, $c'(\frac\pi2)=\frac4{\pi^2}$ and $Q'(0)=-\frac12$ (as $Q(s)=-\frac s2+O(s^3)$), so $\Psi'(0)=-\frac13$ and $\Psi'(\frac\pi2)=\frac4{\pi^2}$.
\end{proof}

\subsection{Proof of Theorem \ref{thm:d2}}
Write $F(y)=\coth y-1=\frac{2}{e^{2y}-1}$ and recall $y_k=2N\operatorname{arsinh}(\sin x_k)$. By Lemma~\ref{lem:elem}(d),
\begin{equation}\label{eq:yk}
\kappa_2k\ \le\ y_k\ \le\ \pi k,\qquad 0\le\pi k-y_k\le\frac{2N}{3}x_k^3\qquad(1\le k\le N-1).
\end{equation}
Introduce the series (certified values in Section~\ref{sec:cert})
\[
\begin{aligned}S_1&=\sum_{j\ge1}\frac{j^2}{\sinh^2(\kappa_2j)}=0.1398576\ldots,& S_2&=\sum_{j\ge1}\frac{2j}{e^{2\kappa_2j}-1}=0.0642886\ldots,\\
S_3&=\sum_{j\ge1}\frac{j^2\cosh(\kappa_2j)}{\sinh^2(\kappa_2j)}=0.7395811\ldots,& S_4&=\sum_{j\ge1}\frac{j}{\sinh(\kappa_2j)}=0.5105105\ldots\end{aligned}
\]

\emph{Step 1 (the sum of $\Phi$).} By Lemmas~\ref{lem:harm} and \ref{lem:trap}, using $\sum_{k=1}^{N-1}\frac1{x_k}=\frac{2N}\pi H_{N-1}$,
\begin{equation}\label{eq:sumPhi}
\sum_{k=1}^{N-1}\Phi(x_k)=\frac{2N}{\pi}\big(\ln N+\gamma+I_\Psi\big)-\frac{2N}{\pi}\eta_N+e_1(N)
\end{equation}
(the terms $-\frac{2N}{\pi}\cdot\frac{1}{2N}$ and $+\frac1\pi$ cancel).

\emph{Step 2 (the exponentially small part).} Since $\frac{1}{x_k}=\frac{2N}{\pi k}$, we have $\frac{2N}\pi\,2\Lambda_*=\sum_{k\ge1}\frac1{x_k}F(\pi k)$, and
\[
E_c(N):=\sum_{k=1}^{N-1}\Phi(x_k)F(y_k)-\frac{2N}{\pi}\,2\Lambda_*=\underbrace{\sum_{k=1}^{N-1}\Psi(x_k)F(y_k)}_{(1)}+\underbrace{\sum_{k=1}^{N-1}\frac{F(y_k)-F(\pi k)}{x_k}}_{(2)}-\underbrace{\sum_{k\ge N}\frac{2N}{\pi k}F(\pi k)}_{(3)}.
\]
$F$ is positive and decreasing with $|F'(y)|=\sinh^{-2}y$ decreasing. By Lemma~\ref{lem:Psi}(b) and \eqref{eq:yk}: $0\ge(1)\ge-1.09\sum_kx_kF(\kappa_2k)\ge-1.09\frac{\pi}{2N}S_2$. By the mean value theorem and \eqref{eq:yk}: $0\le(2)\le\sum_k\frac1{x_k}\frac{2Nx_k^3}{3}\sinh^{-2}(\kappa_2k)=\frac{\pi^2}{6N}S_1$. Finally $0\le(3)\le\frac{2N}{\pi}\cdot\frac1N\sum_{k\ge N}\frac{2e^{-2\pi k}}{1-e^{-2\pi}}=\frac{4e^{-2\pi N}}{\pi(1-e^{-2\pi})^2}$, and $Ne^{-2\pi N}\le2e^{-4\pi}$ for $N\ge2$. Hence
\begin{equation}\label{eq:Ec}
-c_E^-\le N\,E_c(N)\le c_E^+,\qquad c_E^-=1.09\,\frac\pi2S_2+\frac{8e^{-4\pi}}{\pi(1-e^{-2\pi})^2}<0.11009,\qquad c_E^+=\frac{\pi^2}6S_1<0.23006 .
\end{equation}

\emph{Step 3 ($\tau_2$).} By \eqref{eq:tau2exact}, \eqref{eq:sumPhi} and the definition of $E_c$ and $c_\tau$,
\[
\begin{gathered}\tau_2=4N\Big[\sum_k\Phi(x_k)+\sum_k\Phi(x_k)F(y_k)\Big]-\frac23(N^2-1)=\frac8\pi N^2\ln N+c_\tau N^2+R_\tau,\\
R_\tau=\frac23-\frac{8N^2}{\pi}\eta_N+4Ne_1(N)+4NE_c(N).\end{gathered}
\]
By Lemma~\ref{lem:harm}, $-\frac{4}{3\pi}<-\frac{8N^2}\pi\eta_N<0$; by Lemma~\ref{lem:trap}, $|4Ne_1|\le\frac{\pi^2M_2}{12}$; with \eqref{eq:Ec},
\[
-2.1309<\frac23-\frac4{3\pi}-\frac{\pi^2M_2}{12}-4c_E^-\ \le R_\tau\le\ \frac23+\frac{\pi^2M_2}{12}+4c_E^+<3.5197 .
\]

\emph{Step 4 ($m_2-\tau_2$).} By \eqref{eq:tau2exact}, $m_2-\tau_2=\frac23(N^2-1)+4N\,A_N$ with $A_N=\sum_{k=1}^{N-1}(-1)^{k+1}\Phi(x_k)/\sinh y_k$, and $\frac{2N}{\pi}\Sigma_*=\sum_{k\ge1}(-1)^{k+1}\frac{1}{x_k\sinh(\pi k)}$. As in Step 2, with $y\mapsto1/\sinh y$ positive, decreasing, and $|\frac{d}{dy}\frac1{\sinh y}|=\frac{\cosh y}{\sinh^2y}$ decreasing,
\[
\begin{aligned}\Big|A_N-\frac{2N}\pi\Sigma_*\Big|&\le\sum_{k=1}^{N-1}\frac{|\Psi(x_k)|}{\sinh y_k}+\sum_{k=1}^{N-1}\frac1{x_k}\Big(\frac{1}{\sinh y_k}-\frac1{\sinh\pi k}\Big)+\sum_{k\ge N}\frac{2N}{\pi k\sinh(\pi k)}\\
&\le\frac{1.09\,\pi}{2N}S_4+\frac{\pi^2}{6N}S_3+\frac{4e^{-\pi N}}{\pi(1-e^{-\pi})(1-e^{-2\pi N})} .\end{aligned}
\]
For $N\ge2$ the right side is at most $c_A/N$ with $c_A=1.09\frac\pi2S_4+\frac{\pi^2}6S_3+\frac{8e^{-2\pi}}{\pi(1-e^{-\pi})(1-e^{-4\pi})}<2.0957$. Therefore
$m_2-\tau_2=\big(\frac23+\frac8\pi\Sigma_*\big)N^2-\frac23+4N\big(A_N-\frac{2N}\pi\Sigma_*\big)$, i.e.\ \eqref{eq:m2} with $|R_{m\tau}+\frac23|\le4c_A<8.383$. \qed

\begin{observation}\label{obs:d2}
In ball arithmetic (script \texttt{19}; certified comparisons) $R_\tau(N)<R_\tau(N+1)$ and $R_{m\tau}(N)<R_{m\tau}(N+1)$ for every $2\le N<300$, and both sequences keep increasing along $N=200,\,500,\,1000,\,2000,\,5000,\,10^4,\,3\cdot10^4,\,10^5,\,3\cdot10^5,\,10^6$: $R_\tau(N)$ \emph{increases} from $R_\tau(2)=0.851279\ldots$ towards $0.883035\ldots$ and $R_{m\tau}(N)$ \emph{increases} from $R_{m\tau}(2)=-0.530169\ldots$ towards $-0.350894\ldots$ (a check of this note obtained the same picture with separately written $40$-digit floating-point code up to $N=10^6$, within the same project). The two limits are proved in Theorem~\ref{thm:d2limits} below; the monotonicity for \emph{all} $N$ (hence the two-sided bounds $0.85127<R_\tau(N)<0.88304$ and $-0.53017<R_{m\tau}(N)<-0.35089$ for all $N\ge2$) is only observed. Also observed (certified values of $R_\tau(N)$, $R_{m\tau}(N)$ for the listed $N\le10^6$): $(R_\tau(N)-c^{(0)}_\tau)N^2$ settles at $-0.0670577$ and $(R_{m\tau}(N)-c^{(0)}_{m\tau})N^2$ at $-1.0025013$ (rounded to seven decimals, the same values for every listed $N$ from $3\cdot10^4$ to $10^6$), and the sum $-1.069559$ of these two numbers agrees, after rounding to six decimals, with the coefficient $C_{-2}$ of Theorem~4.5 of the article (Eq.~(25)); no rate of convergence is proved here. Thus the true remainders are about four times (resp.\ twenty times) smaller than the bounds proved in Theorem~\ref{thm:d2}.
\end{observation}

\subsection{The constant terms: limits of the remainders}
Besides $S_1,\dots,S_4$ we need the four series
\begin{equation}\label{eq:S58}
\begin{aligned}S_5&=\sum_{k\ge1}\frac{2k}{e^{2\pi k}-1}=0.0037558\ldots,&S_6&=\sum_{k\ge1}\frac{k^2}{\sinh^2(\pi k)}=0.0075537\ldots,\\
S_7&=\sum_{k\ge1}\frac{(-1)^{k+1}k}{\sinh(\pi k)}=0.0795774\ldots,&S_8&=\sum_{k\ge1}\frac{(-1)^{k+1}k^2\cosh(\pi k)}{\sinh^2(\pi k)}=0.0733219\ldots\end{aligned}
\end{equation}
(certified values, script \texttt{19\_certify\_d2\_limits.py}).

\begin{theorem}[$d=2$: the constant terms]\label{thm:d2limits}
For the remainders of Theorem~\ref{thm:d2}, as $N\to\infty$,
\begin{align}
R_\tau(N)&\to c^{(0)}_{\tau}:=\frac23+\frac\pi{18}-\frac{2\pi}{3}S_5+\frac{2\pi^2}{3}S_6=0.8830352\ldots,\label{eq:ctau0}\\
R_{m\tau}(N)&\to c^{(0)}_{m\tau}:=-\frac23-\frac{2\pi}{3}S_7+\frac{2\pi^2}{3}S_8=-0.3508943\ldots,\label{eq:cmtau0}
\end{align}
so that $m_2=\frac8\pi N^2\ln N+C_2N^2+C_2^{(0)}+o(1)$ with $C_2^{(0)}=c^{(0)}_\tau+c^{(0)}_{m\tau}=0.5321408\ldots$ . The limits are proved analytically (no rate is claimed); only the decimal values are computer-assisted.
\end{theorem}
\begin{proof}
\emph{$R_\tau$.} By Step~3 of the proof of Theorem~\ref{thm:d2}, $R_\tau=\frac23-\frac{8N^2}\pi\eta_N+4Ne_1(N)+4NE_c(N)$. By Lemma~\ref{lem:harm}, $\frac{N^2}{12(N+1)^2}<N^2\eta_N<\frac1{12}+\frac1{6N}$, so $\frac{8N^2}{\pi}\eta_N\to\frac{2}{3\pi}$; by \eqref{eq:e1limit}, $4Ne_1(N)\to\frac2{3\pi}+\frac\pi{18}$. It remains to show $NE_c(N)\to-\frac\pi6S_5+\frac{\pi^2}6S_6$. In the notation of Step~2, $NE_c(N)=\sum_{k=1}^{N-1}a_k(N)+\sum_{k=1}^{N-1}b_k(N)-N\cdot(3)$ with
\[
a_k(N)=N\,\Psi(x_k)\,F(y_k),\qquad b_k(N)=\frac{N}{x_k}\big(F(y_k)-F(\pi k)\big),
\]
and the estimates of Step~2 are the $N$-independent, summable majorants
\[
|a_k(N)|\le1.09\,\frac{\pi k}{2}\,F(\kappa_2k),\qquad0\le b_k(N)\le\frac{N}{x_k}\cdot\frac{2Nx_k^3}{3}\cdot\frac{1}{\sinh^2(\kappa_2k)}=\frac{\pi^2k^2}{6\sinh^2(\kappa_2k)}
\]
(their sums are $1.09\frac\pi2S_2$ and $\frac{\pi^2}{6}S_1$), while $0\le N\cdot(3)\le\frac{4Ne^{-2\pi N}}{\pi(1-e^{-2\pi})^2}\to0$. Fix $k$ and let $N\to\infty$. Then $x_k=\frac{\pi k}{2N}\to0$; since $\arsinh(\sin x)=x-\frac{x^3}{3}+O(x^5)$ as $x\to0$ (from $\sin x=x-\frac{x^3}6+O(x^5)$ and $\arsinh s=s-\frac{s^3}{6}+O(s^5)$),
\[
y_k=2N\arsinh(\sin x_k)\to\pi k,\qquad \frac{N}{x_k}\,(\pi k-y_k)=2N^2x_k^2\cdot\frac{x_k-\arsinh(\sin x_k)}{x_k^3}\to\frac{\pi^2k^2}{2}\cdot\frac13=\frac{\pi^2k^2}{6}.
\]
As $\Psi$ is differentiable at $0$ with $\Psi(0)=0$, $\Psi'(0)=-\frac13$ (proof of Lemma~\ref{lem:trap}), $N\Psi(x_k)=\frac{\pi k}{2}\cdot\frac{\Psi(x_k)}{x_k}\to-\frac{\pi k}{6}$, and $F$ is continuous, so $a_k(N)\to-\frac{\pi k}6F(\pi k)$. By the mean value theorem $F(y_k)-F(\pi k)=(\pi k-y_k)\sinh^{-2}(\xi)$ with $\xi\in[y_k,\pi k]$, so $b_k(N)\to\frac{\pi^2k^2}{6}\sinh^{-2}(\pi k)$. By dominated convergence for series (Tannery's theorem; the terms with $k\ge N$ are set to zero),
\[
\sum_{k=1}^{N-1}a_k(N)\to-\frac\pi6\sum_{k\ge1}kF(\pi k)=-\frac\pi6S_5,\qquad\sum_{k=1}^{N-1}b_k(N)\to\frac{\pi^2}6S_6 .
\]
Hence $R_\tau\to\frac23-\frac{2}{3\pi}+\frac{2}{3\pi}+\frac{\pi}{18}+4\big(-\frac\pi6S_5+\frac{\pi^2}{6}S_6\big)$, which is \eqref{eq:ctau0}.

\emph{$R_{m\tau}$.} By Step~4, $R_{m\tau}=-\frac23+4N\big(A_N-\frac{2N}{\pi}\Sigma_*\big)$ and, since $\Phi(x)=\Psi(x)+\frac1x$ and $\frac1{x_k}=\frac{2N}{\pi k}$,
\[
N\Big(A_N-\frac{2N}\pi\Sigma_*\Big)=\sum_{k=1}^{N-1}(-1)^{k+1}\big(a'_k(N)+b_k'(N)\big)-N\sum_{k\ge N}\frac{(-1)^{k+1}2N}{\pi k\sinh(\pi k)},
\]
\[
a_k'(N)=\frac{N\Psi(x_k)}{\sinh y_k},\qquad b_k'(N)=\frac N{x_k}\Big(\frac1{\sinh y_k}-\frac1{\sinh\pi k}\Big).
\]
By the estimates of Step~4, $|a_k'(N)|\le1.09\frac{\pi k}{2}/\sinh(\kappa_2k)$ and $0\le b_k'(N)\le\frac{\pi^2k^2\cosh(\kappa_2k)}{6\sinh^2(\kappa_2k)}$ (summable, independent of $N$), and the last sum is at most $\frac{4Ne^{-\pi N}}{\pi(1-e^{-\pi})(1-e^{-2\pi N})}\to0$ in modulus. For fixed $k$, exactly as above, $a_k'(N)\to-\frac{\pi k}{6\sinh(\pi k)}$ and, with $\frac{d}{dy}\frac{1}{\sinh y}=-\frac{\cosh y}{\sinh^2y}$, $b_k'(N)\to\frac{\pi^2k^2\cosh(\pi k)}{6\sinh^2(\pi k)}$. By dominated convergence $N(A_N-\frac{2N}{\pi}\Sigma_*)\to-\frac\pi6S_7+\frac{\pi^2}6S_8$, which gives \eqref{eq:cmtau0}.

\emph{Values.} The script \texttt{19\_certify\_d2\_limits.py} sums the series \eqref{eq:S58} over $k\le60$ in ball arithmetic and adds a tail bound (for $k\ge1$ each summand is at most $8k^2e^{-\pi k}$ in modulus, the ratio of consecutive bounds is below $\frac1{20}$, and so $\sum_{k>60}8k^2e^{-\pi k}<10^{-75}$); it certifies $0.8830352<c^{(0)}_\tau<0.8830353$, $-0.3508944<c^{(0)}_{m\tau}<-0.3508943$ and $0.5321408<C_2^{(0)}<0.5321409$.
\end{proof}

\begin{remark}[closed forms; not used]\label{rem:d2closed}
The classical evaluations of the Eisenstein series at the lemniscatic point, $E_2(i)=\frac3\pi$ and $E_4(i)=\frac{3\,\Gamma(\frac14)^8}{(2\pi)^6}$, together with Ramanujan's identity $q\frac{dE_2}{dq}=\frac{E_2^2-E_4}{12}$, give
$S_5=\frac1{12}-\frac1{4\pi}$, $S_7=\frac1{4\pi}$, $S_6=\frac{1}{72}\big(E_4(i)-\frac9{\pi^2}\big)$, $S_8=\frac1{8\pi^2}+\frac{\Gamma(\frac14)^8}{512\pi^6}$, and therefore
\[
c^{(0)}_\tau=\frac34+\frac{\Gamma(\frac14)^8}{2304\,\pi^4},\qquad c^{(0)}_{m\tau}=-\frac34+\frac{\Gamma(\frac14)^8}{768\,\pi^4},\qquad C_2^{(0)}=\frac{\Gamma(\frac14)^8}{576\,\pi^4}=0.5321408832\ldots,
\]
the constant $C_0$ of the asymptotic expansion of Theorem~4.5 of the article, which is thereby proved in the series form \eqref{eq:ctau0}--\eqref{eq:cmtau0}. These closed forms are \emph{not} proved in the present note and are not used anywhere; script \texttt{19} verifies each of them against the certified series values to more than $50$ digits (status: classical identities, quoted; numerically confirmed).
\end{remark}

\section{Three dimensions: the Watson-type constant with explicit error}\label{sec:d3}

\subsection{Statement}
With $H$, $P_N$, $\phi$ from \eqref{eq:HP} define, for $\theta\in(0,\pi]$ and $s\in(0,1]$,
\begin{equation}\label{eq:JK}
J(\theta)=\int_0^\pi H(\theta,\theta')\,d\theta',\qquad K(s)=4\int_0^1\sqrt{\frac{(1-t^2)(1+s^2+t^2)}{s^2+t^2}}\;dt,\qquad C_3=\frac{3}{\pi^2}\int_0^\pi J(\theta)\,d\theta ,
\end{equation}
and the constants ($F_*(r)=1/(r(e^{2\pi r}-1))$, $|k|$ the Euclidean norm on $\Z^2$)
\begin{align}
E_\infty&=\frac4\pi\sum_{k\in\Z^2\setminus0}F_*(|k|)+\frac{16}{\pi}\sum_{j\ge1}\int_0^\infty F_*\big(\sqrt{j^2+v^2}\big)\,dv=0.014720007338\ldots,\notag\\
c_{\tau3}&=2+\frac{12}\pi\big(\gamma-\ln\pi-2\ln2\big)+3E_\infty=-5.418839026578\ldots,\label{eq:c3consts}\\
c_{m\tau3}&=1-\frac6\pi\sum_{k\in\Z^2\setminus0}\frac{(-1)^{k_1+k_2}}{|k|\sinh(\pi|k|)}=1.531584884077\ldots .\notag
\end{align}

For $\kappa=(\kappa_1,\kappa_2)\in\R^2\setminus\{0\}$ put $r=|\kappa|$, $\ \chi_+(\kappa)=\dfrac{\kappa_1^4+\kappa_2^4}{6r^2}+\dfrac{r^2}{6}$, $\ \chi_-(\kappa)=\dfrac{\kappa_1^4+\kappa_2^4}{6r^2}-\dfrac{r^2}{2}$, and
\begin{equation}\label{eq:ppa}
\begin{aligned}
p(\kappa)&=\frac{2\pi}{r}\,\chi_-(\kappa)\,\big(\coth(\pi r)-1\big)+\frac{2\pi^2\chi_+(\kappa)}{\sinh^2(\pi r)},&\qquad
p_a(\kappa)&=\frac{2\pi}{r}\,\frac{\chi_-(\kappa)}{\sinh(\pi r)}+\frac{2\pi^2\chi_+(\kappa)\cosh(\pi r)}{\sinh^2(\pi r)},\\
\delta_\infty&=\frac14\sum_{k\in\Z^2\setminus0}p(k)+\sum_{j\ge1}\int_0^\infty p(j,v)\,dv,&\qquad
\alpha_\infty&=\frac14\sum_{k\in\Z^2\setminus0}(-1)^{k_1+k_2}p_a(k)
\end{aligned}
\end{equation}
(absolutely convergent; certified values $\delta_\infty=0.0695430676\ldots$, $\alpha_\infty=-0.2336942394\ldots$).

\begin{theorem}[$d=3$, corner target]\label{thm:d3}
\begin{enumerate}
\item[(a)] \emph{(the constant)} $C_3=\dfrac{3}{\pi^3}\displaystyle\int_{[0,\pi]^3}\frac{\prod_{i=1}^3(1+\cos\theta_i)}{3-\sum_i\cos\theta_i}\,d\theta=3\int_0^\infty\big[e^{-s}(I_0(s)+I_1(s))\big]^3ds=\sum_{y\in\{0,1\}^3}u(y)$, where $I_\nu$ are modified Bessel functions and $u(y)=\sum_{t\ge0}\PP_0(S_t=y)$ is the Green's function of the simple random walk on $\Z^3$ (so $u(0)$ is Watson's integral). Moreover (computer-assisted enclosure)
\[
4.320460169373<C_3<4.320460169375 .
\]
\item[(b)] \emph{(elementary first-order bounds)} For every $N\ge2$,
\begin{equation}\label{eq:tau3}
\begin{gathered}C_3N^3-8.51\,N^2-2.54\ \le\ \tau_3\ \le\ C_3N^3-0.53\,N^2+3.90,\\ N^2-1\le m_3-\tau_3\le\frac{1+\ln6}{\Lambda_1}\le\tfrac32(1+\ln 6)\,N^2 .\end{gathered}
\end{equation}
\item[(c)] \emph{(second order with explicit $O(1)$ remainders; computer-assisted constants)} For every $N\ge2$,
\begin{equation}\label{eq:tau3sharp}
\begin{aligned}
\tau_3&=C_3N^3+c_{\tau3}N^2+R_3(N),&-17.4&\le R_3(N)\le19.4,\\
m_3-\tau_3&=c_{m\tau3}N^2+R_3'(N),&-82.4&\le R_3'(N)\le80.4 .
\end{aligned}
\end{equation}
\item[(d)] \emph{(the constant terms)} As $N\to\infty$,
\begin{gather*}
R_3(N)\to c^{(0)}_{\tau3}:=1+\frac{\pi}{12}+\frac{3\zeta(3)}{4\pi}+3\delta_\infty=1.7573985\ldots,\\
R_3'(N)\to c^{(0)}_{m\tau3}:=-1-3\alpha_\infty=-0.2989172\ldots,
\end{gather*}
so that $m_3=C_3N^3+C_3'N^2+C_3''+o(1)$ with $C_3'=c_{\tau3}+c_{m\tau3}=-3.887254142500\ldots$ and $C_3''=c^{(0)}_{\tau3}+c^{(0)}_{m\tau3}=1.4584812\ldots$ .
\end{enumerate}
\end{theorem}
\begin{remark}
$u(0)=\frac{\sqrt6}{32\pi^3}\Gamma(\frac1{24})\Gamma(\frac5{24})\Gamma(\frac7{24})\Gamma(\frac{11}{24})=1.5163860591\ldots$ (G.~N.~Watson, \emph{Three triple integrals}, Quart.\ J.\ Math.\ Oxford 10 (1939) 266--276; closed form: M.~L.~Glasser and I.~J.~Zucker, \emph{Extended Watson integrals for the cubic lattices}, Proc.\ Natl.\ Acad.\ Sci.\ USA 74 (1977) 1800--1801), and $u(e_1)=u(0)-1$. The constants $C_3=4.32046$ and $c_{\tau3}+c_{m\tau3}=-3.887$ are the constants $K_3$, $K_3'$ of Theorem~4.8 of the article (numerical estimates in its Supplementary Section~S5.3); (c) and (d) prove their values, an explicit $O(1)$ remainder for every $N$, and the next term $C_3''$. The proof of (b) is elementary (Lemmas~\ref{lem:KJ}--\ref{lem:small3}); (c) and (d) need in addition the closed form of $K$ by complete elliptic integrals and a trapezoidal estimate for a function whose second derivative has a logarithmic singularity (Lemmas~\ref{lem:Kclosed}--\ref{lem:quant}).
\end{remark}

\subsection{The functions $K$ and $J$}

\begin{lemma}\label{lem:KJ}
Let $s=\sin\frac\theta2\in(0,1]$.
\begin{enumerate}
\item[(a)] $J(\theta)=2(1-s^2)\,(K(s)-\pi)$, and $K(s)>\pi$.
\item[(b)] $K$ is decreasing and $M(s):=K(s)+4\ln s$ is increasing on $(0,1]$, with $M(0+)=6\ln2-2$.
\item[(c)] $g_3(s):=2s^2(K(s)-\pi)$ is increasing on $(0,1]$, with $g_3(0+)=0$.
\item[(d)] $J(\theta)=-8\ln\sin\frac\theta2+J_b(\theta)$ with $J_b=g_1-g_3$, $g_1(s):=2(M(s)-\pi)$. As functions of $\theta\in[0,\pi]$, $g_1$ and $g_3$ are increasing and absolutely continuous, with
\[
\begin{gathered}g_1(0)=12\ln2-4-2\pi=-1.965419\ldots,\qquad g_3(0)=0,\\ g_1(\pi)=g_3(\pi)=2(K(1)-\pi)=2.206121\ldots\end{gathered}
\]
(the value $K(1)=4.2446532915\ldots$ is certified in Section~\ref{sec:cert}).
\end{enumerate}
\end{lemma}
\begin{proof}
(a) In $\int_0^\pi(1+\cos\theta')\coth\frac\phi2\,d\theta'$ substitute $t=\sin\frac{\theta'}2$: $1+\cos\theta'=2(1-t^2)$, $d\theta'=2\,dt/\sqrt{1-t^2}$ and $\coth\frac\phi2=\sqrt{1+\sinh^2\frac\phi2}/\sinh\frac\phi2=\sqrt{(1+s^2+t^2)/(s^2+t^2)}$. So this integral is $K(s)$; since $\coth>1$ and $\int_0^\pi(1+\cos\theta')d\theta'=\pi$, $K(s)>\pi$ and $J(\theta)=(1+\cos\theta)(K(s)-\pi)$.

(b) For $s>0$ we may differentiate under the integral sign. With $v=s^2+t^2$, $\partial_s\sqrt{(1+v)/v}=-s\,v^{-3/2}(1+v)^{-1/2}$, so
$K'(s)=-4s\int_0^1\sqrt{1-t^2}\,v^{-3/2}(1+v)^{-1/2}dt<0$. Since $\int_0^\infty(s^2+t^2)^{-3/2}dt=s^{-2}$,
\[
M'(s)=K'(s)+\frac4s=4s\Big[\int_0^1\Big(1-\sqrt{\tfrac{1-t^2}{1+s^2+t^2}}\Big)\frac{dt}{(s^2+t^2)^{3/2}}+\int_1^\infty\frac{dt}{(s^2+t^2)^{3/2}}\Big]>0 .
\]
Write $K(s)=4\int_0^1\frac{Q_s(t)-1}{\sqrt{s^2+t^2}}dt+4\operatorname{arsinh}\frac1s$ with $Q_s(t)=\sqrt{(1-t^2)(1+s^2+t^2)}$. Since $\operatorname{arsinh}\frac1s+\ln s=\ln(1+\sqrt{1+s^2})\to\ln2$ and $|Q_s-1|\le|Q_s^2-1|=|s^2(1-t^2)-t^4|\le s^2+t^4$, the integrand is bounded by $s+t^3\le1+t^3$ and dominated convergence gives
$M(0+)=4\int_0^1\frac{\sqrt{1-t^4}-1}{t}dt+4\ln2=2\int_0^1\frac{\sqrt{1-v^2}-1}{v}dv+4\ln 2=2\big[\sqrt{1-v^2}-\ln(1+\sqrt{1-v^2})\big]_0^1+4\ln2=6\ln2-2$.

(c) $g_3'(s)=4s\big[K(s)+\frac s2K'(s)-\pi\big]$ and, by (b),
\[
K+\frac s2K'=4\int_0^1\sqrt{1-t^2}\;\frac{v(1+v)-\frac{s^2}{2}}{v^{3/2}(1+v)^{1/2}}\,dt\ \ge\ 4\int_0^1\sqrt{1-t^2}\;\frac{v+\frac12}{\sqrt{v(1+v)}}\,dt>4\int_0^1\sqrt{1-t^2}\,dt=\pi,
\]
because $s^2\le v$ and $(v+\frac12)^2=v(1+v)+\frac14$. Hence $g_3'>0$. Also $g_3(s)=2s^2(M(s)-4\ln s-\pi)\to0$ as $s\downarrow0$.

(d) $J=2(K-\pi)-2s^2(K-\pi)=2(M-4\ln s-\pi)-g_3$. The map $\theta\mapsto s$ is increasing and smooth; $g_1,g_3$ are continuous on $[0,\pi]$ (with the stated limits at $0$), increasing, and continuously differentiable on $(0,\pi]$. A continuous increasing function $g$ on $[0,\pi]$ that is $C^1$ on $(0,\pi]$ is absolutely continuous: $g(\theta)-g(\epsilon)=\int_\epsilon^\theta g'$ and monotone convergence as $\epsilon\downarrow0$ give $g(\theta)-g(0)=\int_0^\theta g'$ with $g'\ge0$ integrable.
\end{proof}

\begin{lemma}[Riemann sums]\label{lem:riemann}
Let $\theta_k=\pi k/N$ and $\mathcal E_N[f]=\sum_{k=1}^{N-1}f(\theta_k)-\frac N\pi\int_0^\pi f$.
\begin{enumerate}
\item[(a)] If $f$ is increasing on $[0,\pi]$, then $-f(\pi)\le\mathcal E_N[f]\le-f(0)$.
\item[(b)] If $f$ is absolutely continuous on $[0,\pi]$, then $\mathcal E_N[f]\to-\frac12(f(0)+f(\pi))$ as $N\to\infty$.
\item[(c)] $\mathcal E_N\big[\ln\sin\frac{\theta}2\big]=\frac12\ln N+\ln2$.
\end{enumerate}
Consequently
\begin{equation}\label{eq:sumJ}
\sum_{k=1}^{N-1}J(\theta_k)=\frac N\pi\int_0^\pi J-4\ln N-8\ln2+\varepsilon_b(N),\qquad\varepsilon_b(N):=\mathcal E_N[J_b],
\end{equation}
\begin{equation}\label{eq:epsb}
\begin{gathered}-2.2062<-2(K(1)-\pi)\le\varepsilon_b(N)\le2\pi+4-12\ln2+2(K(1)-\pi)<4.1716,\\ \lim_{N\to\infty}\varepsilon_b(N)=\pi+2-6\ln2 .\end{gathered}
\end{equation}
\end{lemma}
\begin{proof}
(a) With $h=\pi/N$, $h\sum_{k=0}^{N-1}f(\theta_k)\le\int_0^\pi f\le h\sum_{k=1}^{N}f(\theta_k)$.
(b) Let $T_N=h[\frac12f(0)+\sum_{k=1}^{N-1}f(\theta_k)+\frac12f(\pi)]$, so that $\mathcal E_N[f]=-\frac12(f(0)+f(\pi))+h^{-1}(T_N-\int_0^\pi f)$. Integrating by parts on each cell (valid for absolutely continuous $f$), $T_N-\int_0^\pi f=\int_0^\pi p_h(x)f'(x)\,dx$ with $p_h(x)=x-\frac12(\theta_{k-1}+\theta_k)$ on $[\theta_{k-1},\theta_k]$; $|p_h|\le\frac h2$ and $\int p_h=0$ over each cell. Given $\epsilon>0$ choose a continuous $\psi$ with $\|f'-\psi\|_{L^1}\le\epsilon$; then $|\int p_hf'|\le\frac h2\epsilon+|\int p_h\psi|$ and $|\int p_h\psi|=|\sum_k\int_{\mathrm{cell}}p_h(x)(\psi(x)-\psi(\theta_k))dx|\le\frac{h}{2}\pi\,\omega_\psi(h)$, with $\omega_\psi$ the modulus of continuity. Hence $h^{-1}|T_N-\int f|\le\frac\epsilon2+\frac\pi2\omega_\psi(h)\to\frac\epsilon2$.
(c) By Lemma~\ref{lem:cot}, $\sum_{k=1}^{N-1}\ln\sin\frac{\theta_k}{2}=\frac12\ln N-(N-1)\ln2$, and $\int_0^\pi\ln\sin\frac\theta2\,d\theta=-\pi\ln2$.
\eqref{eq:sumJ} is (c) and Lemma~\ref{lem:KJ}(d). By (a) applied to $g_1$ and $g_3$, $\varepsilon_b=\mathcal E_N[g_1]-\mathcal E_N[g_3]\in[-g_1(\pi)+g_3(0),\,-g_1(0)+g_3(\pi)]$, which is \eqref{eq:epsb}; by (b), $\varepsilon_b\to-\frac12(g_1(0)+g_1(\pi)-g_3(0)-g_3(\pi))=-\frac12g_1(0)$.
\end{proof}

\subsection{Row sums and the exponentially small parts}

\begin{lemma}[row sums]\label{lem:rows}
For $\theta\in(0,\pi]$ and $N\ge1$,
\[
\sum_{k=0}^{N-1}\varepsilon'_k\,H(\theta,\theta_k)=\frac N\pi J(\theta)+e_N(\theta),\qquad e_N(\theta):=\frac N\pi\int_0^\pi P_N(\theta,\theta')\,d\theta'>0 .
\]
\end{lemma}
\begin{proof}
Let $\mathcal K(\theta_1,\theta_2,\theta_3)=\prod_i(1+\cos\theta_i)/(3-\sum_i\cos\theta_i)$, which is symmetric and positive. By \eqref{eq:int1} (with $\cosh\varphi=3-\cos\theta_1-\cos\theta_2$), $H(\theta_1,\theta_2)=\frac1\pi\int_0^\pi\mathcal K\,d\theta_3$. Fix $\theta_1=\theta$ and let $\chi=\phi(\theta,\theta_3)$, i.e.\ $\cosh\chi=3-\cos\theta-\cos\theta_3$. Summing over $\theta_2=\theta_k$ under the integral and using \eqref{eq:sum1},
\[
\sum_{k=0}^{N-1}\varepsilon'_kH(\theta,\theta_k)=\frac1\pi\int_0^\pi(1+\cos\theta)(1+\cos\theta_3)\,N\Big[\coth\frac\chi2\coth(N\chi)-1\Big]d\theta_3 ,
\]
while integrating over $\theta_2$ first (Tonelli) and using \eqref{eq:int1},
$J(\theta)=\frac1\pi\int_0^\pi\!\!\int_0^\pi\mathcal K\,d\theta_2\,d\theta_3=\int_0^\pi(1+\cos\theta)(1+\cos\theta_3)\big[\coth\frac\chi2-1\big]d\theta_3$. Subtracting $\frac N\pi J(\theta)$ leaves $\frac N\pi\int_0^\pi(1+\cos\theta)(1+\cos\theta_3)\coth\frac\chi2\,(\coth N\chi-1)\,d\theta_3$.
\end{proof}

Let $\kappa_3=\sqrt2\operatorname{arsinh}\sqrt2=1.62099\ldots$ and (certified values in Section~\ref{sec:cert})
\[
B_D=\sum_{k\in\Z^2\setminus0}\frac{1}{|k|\,(e^{2\kappa_3|k|}-1)}=0.197812\ldots,\qquad
B_e=\sum_{j\ge1}\sqrt{\frac{\pi}{4\kappa_3j}}\;\frac{e^{-2\kappa_3j}}{1-e^{-2\kappa_3j}}=0.029091\ldots
\]

\begin{lemma}[exponentially small parts]\label{lem:small3}
Let $\rho=\sinh\frac\phi2=\sqrt{\sin^2\frac{\theta_1}2+\sin^2\frac{\theta_2}2}$. Then for $(\theta_1,\theta_2)\in[0,\pi]^2\setminus\{0\}$
\begin{equation}\label{eq:Pbound}
(1+\cos\theta_1)(1+\cos\theta_2)\coth\frac\phi2\le\frac4\rho,\qquad N\phi\ge\kappa_3N\rho,\qquad\rho\ge\frac{|\theta|}{\pi},\qquad
0<P_N\le\frac4\rho\cdot\frac{2}{e^{2\kappa_3N\rho}-1}.
\end{equation}
Consequently, for $N\ge2$, with $D_N:=\sum'\varepsilon'_{k_1}\varepsilon'_{k_2}P_N(\theta_{k_1},\theta_{k_2})$,
\begin{equation}\label{eq:DNe}
0<D_N\le2B_D\,N,\qquad0<\sum_{k=1}^{N-1}e_N(\theta_k)\le8B_e\,N,\qquad \lim_{N\to\infty}\frac1N\Big[D_N+\sum_{k=1}^{N-1}e_N(\theta_k)\Big]=E_\infty,
\end{equation}
and $\mathrm{Alt}_N:=\sum'\varepsilon'_{k_1}\varepsilon'_{k_2}(-1)^{k_1+k_2}(1+c_{k_1})(1+c_{k_2})\coth\frac\phi2/\sinh(N\phi)$ satisfies $\mathrm{Alt}_N/N\to\frac2\pi\sum_{k\in\Z^2\setminus0}\frac{(-1)^{k_1+k_2}}{|k|\sinh(\pi|k|)}$.
\end{lemma}
\begin{proof}
With $s_i=\sin\frac{\theta_i}2$: $(1+\cos\theta_1)(1+\cos\theta_2)=4(1-s_1^2)(1-s_2^2)\le4(1-\frac{\rho^2}2)^2$ by the arithmetic--geometric mean inequality, $\coth\frac\phi2=\sqrt{1+\rho^2}/\rho$, and $(1-\frac{\rho^2}{2})^2\sqrt{1+\rho^2}\le1$ by Lemma~\ref{lem:elem}(e) with $y=\rho^2/2\in[0,1]$. Since $\rho\le\sqrt2$, Lemma~\ref{lem:elem}(c) gives $\phi=2\operatorname{arsinh}\rho\ge2\rho\operatorname{arsinh}(\sqrt2)/\sqrt2=\kappa_3\rho$, and Lemma~\ref{lem:elem}(a) gives $s_i\ge\theta_i/\pi$. This proves \eqref{eq:Pbound} ($\coth x-1=2/(e^{2x}-1)$ is decreasing).

At a lattice point, $\rho\ge|k|/N$, and $\rho\mapsto\rho^{-1}(e^{2\kappa_3N\rho}-1)^{-1}$ is decreasing, so $P_N\le8N\,\tilde F(|k|)$ with $\tilde F(r)=1/(r(e^{2\kappa_3r}-1))$; as $\sum_{k\in\Z^2\setminus0}f(|k|)=4\sum'_{k\in\mathbb N_0^2\setminus0}\varepsilon'_{k_1}\varepsilon'_{k_2}f(|k|)$, $D_N\le2N\sum_{\Z^2\setminus0}\tilde F=2B_DN$.
Substituting $\theta'=\pi v/N$, $e_N(\theta_k)=\int_0^NP_N(\theta_k,\pi v/N)\,dv\le8N\int_0^\infty\tilde F(\sqrt{k^2+v^2})\,dv$, and with $v=k\sinh x$ (so $dv/\sqrt{k^2+v^2}=dx$)
\[
\int_0^\infty\tilde F(\sqrt{k^2+v^2})\,dv=\int_0^\infty\frac{dx}{e^{2\kappa_3k\cosh x}-1}\le\frac{e^{-2\kappa_3k}}{1-e^{-2\kappa_3k}}\int_0^\infty e^{-\kappa_3kx^2}dx=\sqrt{\frac{\pi}{4\kappa_3k}}\,\frac{e^{-2\kappa_3k}}{1-e^{-2\kappa_3k}},
\]
using $\cosh x\ge1+\frac{x^2}2$. Summing over $k\ge1$ gives the second bound in \eqref{eq:DNe}.

For the limits, fix $k=(k_1,k_2)\ne0$. As $N\to\infty$, $\rho=\frac{\pi|k|}{2N}(1+o(1))$, $N\phi\to\pi|k|$, $(1+c_{k_1})(1+c_{k_2})\to4$, hence $P_N/N\to\frac{16}{\pi}F_*(|k|)$, and likewise for fixed $k\ge1$, $v\ge0$: $P_N(\theta_k,\pi v/N)/N\to\frac{16}\pi F_*(\sqrt{k^2+v^2})$. The bounds just proved, $P_N/N\le8\tilde F$, are summable (integrable) and independent of $N$, so dominated convergence gives
$D_N/N\to\frac{16}\pi\sum'\varepsilon'\varepsilon'F_*(|k|)=\frac4\pi\sum_{\Z^2\setminus0}F_*(|k|)$ and $\frac1N\sum_ke_N(\theta_k)\to\frac{16}\pi\sum_{k\ge1}\int_0^\infty F_*(\sqrt{k^2+v^2})dv$. The statement on $\mathrm{Alt}_N$ follows in the same way: its terms divided by $N$ converge to $(-1)^{k_1+k_2}\frac{8}{\pi}\frac{1}{|k|\sinh\pi|k|}$ and are dominated by $\frac{4}{|k|\sinh(\kappa_3|k|)}$.
\end{proof}

\subsection{Proof of Theorem \ref{thm:d3}(a),(b)}
\emph{(a)} $C_3=\frac{3}{\pi^2}\int_0^\pi\!\!\int_0^\pi H=\frac3{\pi^3}\int_{[0,\pi]^3}\mathcal K$ by the proof of Lemma~\ref{lem:rows}. Writing $1/(3-\sum_i\cos\theta_i)=\int_0^\infty e^{-s(3-\sum_i\cos\theta_i)}ds$ and using Tonelli's theorem,
$C_3=3\int_0^\infty a(s)^3ds$ with $a(s)=\frac1\pi\int_0^\pi(1+\cos\theta)e^{-s(1-\cos\theta)}d\theta=e^{-s}(I_0(s)+I_1(s))$ (DLMF 10.32.3). Expanding $\prod_i(1+\cos\theta_i)=\sum_{y\in\{0,1\}^3}\prod_i\cos(y_i\theta_i)$ gives $C_3=\sum_yu(y)$ with $u(y)=\pi^{-3}\int_{[0,\pi]^3}\prod_i\cos(y_i\theta_i)\,(1-\frac13\sum_i\cos\theta_i)^{-1}d\theta$, the Fourier representation of the Green's function of the simple random walk on $\Z^3$. The enclosure is proved in Lemma~\ref{lem:C3} below.

\emph{(b)} Split the sum in \eqref{eq:tau3exact} into the rows $k_1\ge1$ and the row $k_1=0$. By Lemma~\ref{lem:rows} and \eqref{eq:sumJ},
\[
\sum_{k_1=1}^{N-1}\sum_{k_2=0}^{N-1}\varepsilon'_{k_2}H(\theta_{k_1},\theta_{k_2})=\frac N\pi\Big[\frac N\pi\int_0^\pi J-4\ln N-8\ln2+\varepsilon_b(N)\Big]+\sum_{k=1}^{N-1}e_N(\theta_k).
\]
For the row $k_1=0$: $\sinh\frac\phi2=\sin\frac{\theta}2$, so with $x=\theta/2$, $H(0,\theta)=4\cos^2x\,(\sqrt{1+\sin^2x}/\sin x-1)=4\Phi(x)-4\cos^2x$, and by \eqref{eq:sumPhi} and Lemma~\ref{lem:cot}
\[
\frac12\sum_{k=1}^{N-1}H(0,\theta_k)=2\sum_{k=1}^{N-1}\Phi(x_k)-(N-1)=\frac{4N}{\pi}(\ln N+\gamma+I_\Psi)-\frac{4N}\pi\eta_N+2e_1(N)-N+1 .
\]
The terms $\mp\frac{4N}{\pi}\ln N$ cancel. Inserting into \eqref{eq:tau3exact} and using $\frac{3N^3}{\pi^2}\int_0^\pi J=C_3N^3$, $(N-1)(2N-1)=2N^2-3N+1$:
\begin{equation}\label{eq:tau3identity}
\tau_3=C_3N^3+N^2\Big[A_0+\frac3\pi\varepsilon_b(N)\Big]+3N\Big[\sum_{k=1}^{N-1}e_N(\theta_k)+D_N\Big]+\rho_N,
\end{equation}
\[
\begin{gathered}A_0=-1-\frac{24\ln2}{\pi}+\frac{12}\pi(\gamma+I_\Psi)=-6.4014175\ldots,\\
\rho_N=1-\frac{12N^2}{\pi}\eta_N+6N\,e_1(N)\in\Big(1-\frac2\pi-\frac{\pi^2M_2}{8},\,1+\frac{\pi^2M_2}{8}\Big)\subset(-2.536,3.900),\end{gathered}
\]
by Lemmas~\ref{lem:harm} and \ref{lem:trap}. With \eqref{eq:epsb} and \eqref{eq:DNe},
\[
\begin{aligned}-8.5082<A_0-\frac6\pi(K(1)-\pi)\ &\le\ \frac{\tau_3-C_3N^3-\rho_N}{N^2}\\ &\le\ A_0+\tfrac3\pi\big(2\pi+4-12\ln2+2(K(1)-\pi)\big)\\ &\qquad+3(8B_e+2B_D)<-0.5328 ,\end{aligned}
\]
which is the first part of \eqref{eq:tau3}. The second part is Lemma~\ref{lem:arrival} with $d=3$, $W\le6$ and $\Lambda_1=\frac23\sin^2\frac{\pi}{2N}\ge\frac{2}{3N^2}$.

Parts (c) and (d) are proved in Section~\ref{sec:d3second}. \qed

\begin{lemma}[enclosure of $C_3$]\label{lem:C3}
Let $a(s)=e^{-s}(I_0(s)+I_1(s))$. Then $a(s)=\frac4\pi\int_0^1\sqrt{1-t^2}\,e^{-2st^2}dt$, and
\[
\sqrt{\tfrac{2}{\pi s}}\Big(1-\frac{1}{8s}-\frac{3}{32s^2}\Big)\le a(s)\le\sqrt{\tfrac2{\pi s}}\quad(s>0),\qquad a(s)\le\sqrt{\tfrac{2}{\pi s}}\Big(1-\frac1{8s}\Big)+\frac{e^{-2s}}{\pi}\quad(s\ge1).
\]
Consequently, for $S\ge200$, with $b=(2/\pi)^{3/2}$,
$\ b\big(2S^{-1/2}-\frac14S^{-3/2}-\frac9{80}S^{-5/2}\big)\le\int_S^\infty a^3\le b\big(2S^{-1/2}(1+4e^{-150})-\frac14S^{-3/2}+\frac{3}{160}S^{-5/2}\big)$, and $C_3=4.32046016937353\ldots$, in particular $4.320460169373<C_3<4.320460169375$.
\end{lemma}
\begin{proof}
In $a(s)=\frac1\pi\int_0^\pi(1+\cos\theta)e^{-s(1-\cos\theta)}d\theta$ substitute $t=\sin\frac\theta2$. By Lemma~\ref{lem:elem}(f) with $y=t^2$, $1-\frac{t^2}2-\frac{t^4}2\le\sqrt{1-t^2}\le1-\frac{t^2}2\le1$ on $[0,1]$. Using $\int_0^\infty t^{2j}e^{-2st^2}dt=\frac12\sqrt{\frac{\pi}{2s}}\,\{1,\frac1{4s},\frac{3}{16s^2}\}$ for $j=0,1,2$:
the bound $\sqrt{1-t^2}\le1$ gives $a\le\sqrt{2/(\pi s)}$; the lower bound gives $a\ge\frac4\pi\int_0^\infty(1-\frac{t^2}2-\frac{t^4}{2})e^{-2st^2}dt$ (the integrand is negative for $t>1$), which is the stated lower bound; and $a\le\frac4\pi\int_0^\infty(1-\frac{t^2}{2})e^{-2st^2}dt+\frac4\pi\int_1^\infty\frac{t^2}2e^{-2st^2}dt$, where $\int_1^\infty\frac{t^2}{2}e^{-2st^2}dt\le\int_1^\infty\frac{t^3}2e^{-2st^2}dt=\frac{e^{-2s}}4(\frac1{2s}+\frac{1}{4s^2})\le\frac{e^{-2s}}{4}$ for $s\ge1$.
For $s\ge200$ the last term is at most $e^{-2s}\sqrt s\cdot\sqrt{2/(\pi s)}\,(1-\frac1{8s})$, so $a^3\le b\,s^{-3/2}(1-\frac1{8s})^3(1+e^{-2s}\sqrt s)^3\le b\,s^{-3/2}(1-\frac{3}{8s}+\frac{3}{64s^2})(1+4e^{-150})$, and $a^3\ge b\,s^{-3/2}(1-\frac{3}{8s}-\frac{9}{32s^2})$ by Bernoulli's inequality; integrate from $S$ to $\infty$ (in the upper bound the factor $1+4e^{-150}$ is needed only on the leading term since the sum of the two others is negative). With $S=10^6$ the two tail bounds differ by $3b\cdot\frac{21}{160}S^{-5/2}<2.1\cdot10^{-16}$. The integral $\int_0^{S}a^3$ of the entire function $a^3$ is enclosed rigorously by the Petras algorithm of Arb (\texttt{acb.integral}; F.~Johansson, \emph{Numerical integration in arbitrary-precision ball arithmetic}, Proc.\ ICMS 2018, LNCS 10931, 255--263) in \texttt{06\_certify\_constants.py}; adding the tail gives $C_3=4.320460169373537\pm2.3\cdot10^{-16}$ (the same computation with $S=200$ gives the consistent enclosure $[4.3204598,4.3204603]$).
\end{proof}

\subsection{Second order: closed form of $K$, the $O(1)$ remainder and the constant term}\label{sec:d3second}

Throughout, $\mathbf K(k)=\int_0^{\pi/2}(1-k^2\sin^2\varphi)^{-1/2}d\varphi$ and $\mathbf E(k)=\int_0^{\pi/2}(1-k^2\sin^2\varphi)^{1/2}d\varphi$ are the complete elliptic integrals with modulus $k\in[0,1)$, $k'=\sqrt{1-k^2}$. We use the classical facts (DLMF 19.4.1, 19.4.2, 19.9.1, 19.12.1)
\begin{equation}\label{eq:ellfacts}
\frac{d\mathbf K}{dk}=\frac{\mathbf E-k'^2\mathbf K}{k\,k'^2},\qquad\frac{d\mathbf E}{dk}=\frac{\mathbf E-\mathbf K}{k},\qquad
\ln4\le\mathbf K(k)+\ln k'\le\frac\pi2,\qquad\lim_{k\to1}\big(\mathbf K(k)+\ln k'\big)=\ln4,
\end{equation}
and $1\le\mathbf E\le\frac\pi2$, $\mathbf E(1)=1$, $\mathbf K>\mathbf E$ for $k>0$ (obvious from the integrals). The limit is the term $m=0$ of the expansion DLMF 19.12.1, $\mathbf K(k)=\sum_{m\ge0}\big(\frac{(1/2)_m}{m!}\big)^2k'^{2m}\big(\ln\frac1{k'}+d(m)\big)$ with $d(0)=\psi(1)-\psi(\frac12)=\ln4$.

\begin{lemma}[closed form of $K$; the functions $B$, $b$, $e$, $\ell$]\label{lem:Kclosed}
For $s\in(0,1]$ let $k=\dfrac1{1+s^2}$, so that $k'=\dfrac{s\sqrt{s^2+2}}{1+s^2}$, and write $\mathbf K=\mathbf K(k)$, $\mathbf E=\mathbf E(k)$. Define
\begin{gather*}
B(s)=1-(1+s^2)\mathbf E+\frac{s^2(s^2+2)}{1+s^2}\mathbf K,\qquad b(s)=\frac{B(s)}{s^2}-(\mathbf K-\mathbf E),\\
e(s)=\frac{\mathbf E}{(1+s^2)(2+s^2)},\qquad\ell(s)=\mathbf K-\mathbf E+\ln s .
\end{gather*}
\begin{enumerate}
\item[(a)] $K(s)=(4+2s^2)\,\mathbf K-(2+2s^2)\,\mathbf E$.
\item[(b)] With $'=d/ds$:
\begin{gather*}
\mathbf E'=\frac{2s(\mathbf K-\mathbf E)}{1+s^2},\qquad(\mathbf K-\mathbf E)'=-\frac{2e}{s},\qquad M'=\frac{4B}{s},\qquad B'=2s(\mathbf K-\mathbf E),\\
M''=8(\mathbf K-\mathbf E)-\frac{4B}{s^2},\qquad
\Big(\frac{B}{s^2}\Big)'=-\frac{2b}{s},\qquad(s^2b)'=2s\,e,\\ b'=\frac{2(e-b)}{s},\qquad\ell'=\frac{1-2e}{s}.
\end{gather*}
\item[(c)] $B(0+)=0$, $B>0$, $b>0$; $M$, $B$, $\mathbf E$ and $s^2b$ are increasing in $s$ and $\mathbf K-\mathbf E$ is decreasing; $e(0+)=b(0+)=\frac12$, $\ell(0+)=\frac32\ln2-1$, $M(0+)=6\ln2-2$.
\item[(d)] For $0<s\le0.78$: $\ k'\ge s$, $\ 0\le\mathbf E-1\le k'^2\mathbf K\le2s^2(\frac\pi2+\ln\frac1s)$, $\ 0<\dfrac{B}{s^2}\le\dfrac\pi2-\dfrac12+\ln\dfrac1s$, and with $w:=1+2b-4e$: $\ |w(s)|\le s^2\big(13.11+5\ln\frac1s\big)$. For all $s\in(0,1]$: $0<b\le\frac\pi4$, $|\ell|\le0.72$, $-0.983\le M-\pi\le1.104$.
\end{enumerate}
\end{lemma}
\begin{proof}
(a) Fix $\theta\in(0,\pi]$, $s=\sin\frac\theta2$, $E=3-\cos\theta=2+2s^2$, and put
\[
g_{mn}=\frac1{\pi^{2}}\int_{[0,\pi]^2}\frac{\cos(m\theta_2)\cos(n\theta_3)}{E-\cos\theta_2-\cos\theta_3}\,d\theta_2\,d\theta_3 .
\]
By the proof of Lemma~\ref{lem:rows} and Lemma~\ref{lem:KJ}(a), $(1+\cos\theta)(K(s)-\pi)=J(\theta)=(1+\cos\theta)\,\pi\,(g_{00}+2g_{10}+g_{11})$, using $g_{10}=g_{01}$. Integrating the identity $1=(E-\cos\theta_2-\cos\theta_3)/(E-\cos\theta_2-\cos\theta_3)$ gives $Eg_{00}-2g_{10}=1$, hence
\[
K(s)=\pi\big[(1+E)\,g_{00}+g_{11}\big].
\]
Let $\Phi$ be continuous on $\R^2$, $2\pi$-periodic and even in each variable. Then $\pi^{-2}\int_{[0,\pi]^2}\Phi=(2\pi)^{-2}\int_{[-\pi,\pi]^2}\Phi(\theta_2,\theta_3)\,d\theta_2d\theta_3=(2\pi)^{-2}\int_{[-\pi,\pi]^2}\Phi(u+v,u-v)\,du\,dv$: for fixed $v$, periodicity in $u$ gives $\int_{-\pi}^{\pi}\Phi(u+v,u-v)\,du=\int_{-\pi}^{\pi}\Phi(w,w-2v)\,dw$, and for fixed $w$, as $v$ runs over $[-\pi,\pi]$ the point $w-2v$ runs twice over a period, so $\frac1{2\pi}\int_{-\pi}^{\pi}\Phi(w,w-2v)\,dv=\frac1{2\pi}\int_{-\pi}^{\pi}\Phi(w,t)\,dt$. Since $\cos\theta_2+\cos\theta_3=2\cos u\cos v$ and $\cos\theta_2\cos\theta_3=\frac12(\cos2u+\cos2v)$, and $\frac1{2\pi}\int_{-\pi}^{\pi}\frac{dv}{E-\beta\cos v}=(E^2-\beta^2)^{-1/2}$ for $|\beta|<E$ (this is \eqref{eq:poisson} integrated over a period), we get with $k=2/E=1/(1+s^2)$
\[
g_{00}=\frac1{2\pi}\int_{-\pi}^{\pi}\frac{du}{\sqrt{E^2-4\cos^2u}}=\frac{2}{\pi E}\mathbf K,\qquad
g_{11}=\frac1{2\pi}\int_{-\pi}^{\pi}\frac{\cos2u\,du}{\sqrt{E^2-4\cos^2u}}=\frac{2}{\pi E}\Big[\frac{2}{k^2}(\mathbf K-\mathbf E)-\mathbf K\Big]
\]
(in $g_{11}$ the two terms $\cos2u$ and $\cos2v$ contribute equally by the symmetry $u\leftrightarrow v$; then $\cos2u=2\cos^2u-1$ and $\int_0^{\pi/2}\cos^2u\,(1-k^2\cos^2u)^{-1/2}du=(\mathbf K-\mathbf E)/k^2$). Therefore $K=\frac2E\big[(1+E)\mathbf K+\frac{E^2}{2}(\mathbf K-\mathbf E)-\mathbf K\big]=(2+E)\mathbf K-E\,\mathbf E$.

(b) $dk/ds=-2sk^2$ and $k'^2=s^2(s^2+2)/(1+s^2)^2$, so by \eqref{eq:ellfacts}
$\mathbf K'=-\frac{2(1+s^2)}{s(s^2+2)}(\mathbf E-k'^2\mathbf K)$ and $\mathbf E'=\frac{2s(\mathbf K-\mathbf E)}{1+s^2}$; moreover $\frac{d(\mathbf K-\mathbf E)}{dk}=\frac{k\mathbf E}{k'^2}$, which gives $(\mathbf K-\mathbf E)'=-\frac{2\mathbf E}{s(1+s^2)(s^2+2)}=-\frac{2e}s$. Differentiating (a):
$K'=4s(\mathbf K-\mathbf E)+(4+2s^2)\mathbf K'-(2+2s^2)\mathbf E'=-\frac4s\big[(1+s^2)\mathbf E-\frac{s^2(s^2+2)}{1+s^2}\mathbf K\big]=-\frac4s+\frac{4B}s$, i.e.\ $M'=K'+\frac4s=\frac{4B}{s}$. With $c(s)=\frac{s^2(s^2+2)}{1+s^2}$ we have $c\,\mathbf K'=-2s(\mathbf E-k'^2\mathbf K)$, $c'=\frac{2s(2+2s^2+s^4)}{(1+s^2)^2}$ and $c'-2s+2sk'^2=2s$, hence
$B'=-2s\mathbf E-(1+s^2)\mathbf E'+c'\mathbf K+c\mathbf K'=-2s\mathbf E+(c'-2s+2sk'^2)\mathbf K=2s(\mathbf K-\mathbf E)$.
The remaining identities follow: $(B/s^2)'=\frac{2(\mathbf K-\mathbf E)}{s}-\frac{2B}{s^3}=-\frac{2b}s$; $b'=-\frac{2b}{s}+\frac{2e}{s}$; $(s^2b)'=2sb+s^2b'=2se$; $\ell'=-\frac{2e}s+\frac1s$; $M''=(4B/s)'=8(\mathbf K-\mathbf E)-4B/s^2$.

(c) As $s\downarrow0$, $k\to1$, $\mathbf E\to1$ and $c\,\mathbf K\le3s^2(\frac\pi2+\ln\frac1{k'})\to0$, so $B(0+)=0$ and $B(s)=\int_0^s2\sigma(\mathbf K-\mathbf E)\,d\sigma>0$; likewise $s^2b=B-s^2(\mathbf K-\mathbf E)\to0$, so $s^2b(s)=\int_0^s2\sigma e(\sigma)\,d\sigma>0$. The monotonicity statements are the signs of the derivatives in (b) ($M'>0$, $B'>0$, $\mathbf E'>0$, $(s^2b)'>0$, $(\mathbf K-\mathbf E)'<0$). $e$ is continuous on $[0,1]$ with $e(0)=\frac12\mathbf E(1)=\frac12$, hence $b(s)=s^{-2}\int_0^s2\sigma e\to\frac12$. By \eqref{eq:ellfacts}, $\mathbf K+\ln s=(\mathbf K+\ln k')+\ln\frac s{k'}\to\ln4-\frac12\ln2=\frac32\ln2$, so $\ell(0+)=\frac32\ln2-1$ and, by (a), $M(0+)=\lim\,[4(\mathbf K+\ln s)+2s^2\mathbf K-(2+2s^2)\mathbf E]=6\ln2-2$ (in agreement with Lemma~\ref{lem:KJ}(b)).

(d) $k'\ge s\iff\sqrt{s^2+2}\ge1+s^2\iff s^4+s^2\le1$, true for $s\le0.78$. Next $\mathbf E-1=\int_0^{\pi/2}\big[\sqrt{\cos^2\varphi+k'^2\sin^2\varphi}-\cos\varphi\big]d\varphi=\int_0^{\pi/2}\frac{k'^2\sin^2\varphi\,d\varphi}{\sqrt{\cos^2\varphi+k'^2\sin^2\varphi}+\cos\varphi}\le k'^2\mathbf K$, and $k'^2\le2s^2$, $\mathbf K\le\frac\pi2+\ln\frac1{k'}\le\frac\pi2+\ln\frac1s$. Hence $B=\int_0^s2\sigma(\mathbf K-\mathbf E)\le\int_0^s2\sigma(\frac\pi2-1+\ln\frac1\sigma)d\sigma=s^2(\frac\pi2-\frac12+\ln\frac1s)$. Write $e=\frac12+\eta$; then $\eta(\sigma)=\frac{\mathbf E-1}{(1+\sigma^2)(2+\sigma^2)}-\frac{3\sigma^2+\sigma^4}{2(1+\sigma^2)(2+\sigma^2)}$, where the first term lies in $[0,\sigma^2(\frac\pi2+\ln\frac1\sigma)]$ and the second in $[-\sigma^2,0]$, so $|\eta(\sigma)|\le\sigma^2(\frac\pi2+1+\ln\frac1\sigma)$. Since $b-\frac12=s^{-2}\int_0^s2\sigma\eta(\sigma)d\sigma$, $|b-\frac12|\le\frac2{s^2}\int_0^s\sigma^3(\frac\pi2+1+\ln\frac1\sigma)d\sigma=s^2\big(\frac{\pi/2+1}{2}+\frac12\ln\frac1s+\frac18\big)$. Therefore $|w|=|2(b-\frac12)-4\eta|\le s^2\big(5(\frac\pi2+1)+\frac14+5\ln\frac1s\big)\le s^2(13.11+5\ln\frac1s)$.
For all $s$: $\frac16\le e\le\frac\pi4$ (as $1\le\mathbf E\le\frac\pi2$ and $2\le(1+s^2)(2+s^2)\le6$), so $0<b=s^{-2}\int_0^s2\sigma e\le\frac\pi4$; $\ell=(\mathbf K+\ln k')-\mathbf E+\ln\frac{s}{k'}$ with $\ln\frac s{k'}=\ln\frac{1+s^2}{\sqrt{s^2+2}}\in[-\frac12\ln2,\ln\frac2{\sqrt3}]$, so $\ell\in[\ln4-\frac\pi2-\frac12\ln2,\ \frac\pi2-1+\ln\frac2{\sqrt3}]\subset[-0.54,0.72]$; and $M$ increases from $6\ln2-2$ to $K(1)$, where $6\ln2-2-\pi=-0.9827\ldots$ and $K(1)-\pi=1.1030\ldots$ (Lemma~\ref{lem:KJ}(d)).
\end{proof}

\begin{lemma}[the second derivative of $J_b$]\label{lem:Jb2}
Let $J_b$ be as in Lemma~\ref{lem:KJ}(d), i.e.\ $J_b(\theta)=F(s):=2(1-s^2)(M(s)-\pi)+8s^2\ln s$ with $s=\sin\frac\theta2$.
\begin{enumerate}
\item[(a)] $J_b\in C^1[0,\pi]\cap C^\infty(0,\pi]$, $J_b'(0)=J_b'(\pi)=0$, and $J_b'$ is absolutely continuous on $[0,\pi]$.
\item[(b)] For $\theta\in(0,\pi]$, $\ J_b''(\theta)=2\ln\sin\frac\theta2+g(\theta)$ with $g(\theta)=\frac14G(\sin\frac\theta2)$ and
\begin{equation}\label{eq:Gdef}
G(s)=-4(1-2s^2)(M-\pi)-40(1-s^2)B-8(1-s^2)^2(b-\ell)-(16s^2+8s^4)\ln s+24-32s^2 .
\end{equation}
$G$ extends continuously to $[0,1]$, with $G(0)=4\pi-12\ln2+20=24.2486\ldots$ and $G(1)=4(K(1)-\pi)-8=-3.5877\ldots$ .
\item[(c)] $G\in C^1(0,1]$ and, with $w=1+2b-4e$,
\begin{equation}\label{eq:Gprime}
\begin{aligned}G'(s)=\ &s\Big\{16(M-\pi)-96(1-2s^2)\frac{B}{s^2}+80(1-s^2)\,b+32(1-s^2)(b-\ell)\\&\quad-(32+32s^2)\ln s-80-8s^2\Big\}+\frac{8(1-s^2)^2\,w}{s}.\end{aligned}
\end{equation}
For $0<s\le0.002$: $|G'(s)|\le s\,(416.4+168.01\ln\frac1s)$.
\item[(d)] (computer-assisted) $\displaystyle\int_0^\pi|g'(\theta)|\,d\theta=\frac14\int_0^1|G'(s)|\,ds\le V_g:=6.97$; in particular $g$ is absolutely continuous on $[0,\pi]$. (In fact $G'<0$ on $[0.002,1]$, and $\frac14(G(0)-G(1))=6.9590\ldots$)
\end{enumerate}
\end{lemma}
\begin{proof}
(b) With $c=\cos\frac\theta2$: $ds/d\theta=\frac c2$, $d^2s/d\theta^2=-\frac s4$, so $J_b'=\frac c2F'(s)$ and $4J_b''=(1-s^2)F''(s)-sF'(s)$. By Lemma~\ref{lem:Kclosed}(b), $sM'=4B$ and $M''=8(\mathbf K-\mathbf E)-4B/s^2$, so
$F'=-4s(M-\pi)+\frac{8(1-s^2)B}{s}+16s\ln s+8s$ and $F''=-4(M-\pi)-32B+2(1-s^2)\big[8(\mathbf K-\mathbf E)-\frac{4B}{s^2}\big]+16\ln s+24$. Hence
\[
4J_b''=-4(1-2s^2)(M-\pi)-40(1-s^2)B-\frac{8(1-s^2)^2B}{s^2}+16(1-s^2)^2(\mathbf K-\mathbf E)+16(1-2s^2)\ln s+24-32s^2 .
\]
Insert $B/s^2=b+\ell-\ln s$ and $\mathbf K-\mathbf E=\ell-\ln s$: the terms with $\ln s$ add up to $[8(1-s^2)^2-16(1-s^2)^2+16(1-2s^2)]\ln s=8\ln s-(16s^2+8s^4)\ln s$, which gives $4J_b''=8\ln s+G(s)$ with $G$ as in \eqref{eq:Gdef}. The limits at $0$ follow from Lemma~\ref{lem:Kclosed}(c): $G(0)=-4(6\ln2-2-\pi)-8(\frac12-\frac32\ln2+1)+24$.

(a) $F$ is smooth on $(0,1]$. As $s\downarrow0$, $F'(s)\to0$ because $B/s\le s(\frac\pi2-\frac12+\ln\frac1s)$; hence $J_b'(\theta)\to0$ as $\theta\downarrow0$, and since $J_b$ is continuous on $[0,\pi]$ (Lemma~\ref{lem:KJ}) the mean value theorem gives $J_b\in C^1[0,\pi]$ with $J_b'(0)=0$. $J_b'(\pi)=0$ because $c=0$ at $\theta=\pi$ and $F'(1)$ is finite. By (b), $J_b''$ is integrable on $(0,\pi)$, so $J_b'$ is absolutely continuous.

(c) Differentiate \eqref{eq:Gdef} term by term with Lemma~\ref{lem:Kclosed}(b):
$\frac{d}{ds}[-4(1-2s^2)(M-\pi)]=16s(M-\pi)-16(1-2s^2)\frac Bs$;
$\frac{d}{ds}[-40(1-s^2)B]=80sB-80s(1-s^2)(\mathbf K-\mathbf E)=-80s(1-2s^2)\frac{B}{s^2}+80s(1-s^2)b$ (use $\mathbf K-\mathbf E=B/s^2-b$);
$\frac{d}{ds}[-8(1-s^2)^2(b-\ell)]=32s(1-s^2)(b-\ell)-\frac{8(1-s^2)^2}{s}\big(2(e-b)-(1-2e)\big)=32s(1-s^2)(b-\ell)+\frac{8(1-s^2)^2w}{s}$;
$\frac{d}{ds}[-(16s^2+8s^4)\ln s+24-32s^2]=-(32s+32s^3)\ln s-16s-8s^3-64s$. The sum is \eqref{eq:Gprime}.
For $s\le0.002$, by Lemma~\ref{lem:Kclosed}(d) (with $L=\ln\frac1s$): $|16(M-\pi)|\le17.67$; $96B/s^2\le96(1.0708+L)$; $80b\le62.84$; $32|b-\ell|\le32(\frac\pi4+0.72)\le48.18$; $(32+32s^2)|\ln s|\le32.01L$; $80+8s^2\le80.01$; and $8|w|/s\le s(104.9+40L)$. Adding: $|G'|\le s\,(416.4+168.01L)$.

(d) Since $\theta\mapsto s$ is increasing and smooth, $\int_0^\pi|g'(\theta)|d\theta=\frac14\int_0^1|G'(s)|ds$. By (c), $\int_0^{s_1}|G'|\le s_1^2\big[\frac{416.4}{2}+168.01(\frac12\ln\frac1{s_1}+\frac14)\big]<0.0031$ for $s_1=0.002$. On $[s_1,1]$ the script \texttt{16\_certify\_d3\_remainder.py} uses $90\,108$ cells $[s_i,s_{i+1}]$. On each cell the monotone functions $M$, $B$, $\mathbf E$ (increasing) and $\mathbf K-\mathbf E$ (decreasing) are enclosed by the hull of their end-point values (computed from Lemma~\ref{lem:Kclosed}(a) with Arb's complete elliptic integrals at $256$ bits), $s$ and $\ln s$ by the hull of the end points, and $w$ by $w(s_i)\pm(s_{i+1}-s_i)\sup|w'|$ with $w'=\frac4s(e-b)-4e'$, $e'=\frac{2s(\mathbf K-\mathbf E)}{(1+s^2)^2(2+s^2)}-\frac{\mathbf E\,(6s+4s^3)}{((1+s^2)(2+s^2))^2}$ evaluated on the cell; then \eqref{eq:Gprime} is evaluated in ball arithmetic. The sum of (cell width)$\times$(upper bound of $|G'|$), plus the bound for $[0,s_1]$, divided by $4$, is $<6.962$; the same computation certifies $G'<0$ on every cell.
\end{proof}

\begin{lemma}[trapezoidal sums for a logarithmically singular second derivative]\label{lem:traplog}
Let $f\in C^1[0,\pi]$ with $f'$ absolutely continuous, $f'(0)=f'(\pi)=0$, and suppose $f''(\theta)=2\ln\sin\frac\theta2+g(\theta)$ for a.e.\ $\theta$, where $g$ is absolutely continuous on $[0,\pi]$. Then for every $N\ge1$, in the notation of Lemma~\ref{lem:riemann},
\[
\mathcal E_N[f]=-\frac{f(0)+f(\pi)}2+\frac{\zeta(3)}{4N^2}+r_N,\qquad|r_N|\le\frac{\sqrt3\,\pi^2}{216\,N^2}\int_0^\pi|g'|,\qquad N^2r_N\to0\quad(N\to\infty).
\]
\end{lemma}
\begin{proof}
Let $h=\pi/N$, $T_N=h[\frac12f(0)+\sum_{k=1}^{N-1}f(\theta_k)+\frac12f(\pi)]$ and, on each cell $[\theta_{k-1},\theta_k]$, $q_h(\theta)=\frac12(\theta-\theta_{k-1})(\theta_k-\theta)$. Two integrations by parts on each cell give $T_N-\int_0^\pi f=\int_0^\pi q_hf''$. Put $\varkappa_h=q_h-\frac{h^2}{12}$; since $\int_0^\pi f''=f'(\pi)-f'(0)=0$, $T_N-\int_0^\pi f=\int_0^\pi\varkappa_hf''$. With $t=\{\theta/h\}$, $\varkappa_h=-\frac{h^2}{2}B_2(t)$, $B_2(t)=t^2-t+\frac16$, and (DLMF 24.8.1) $B_2(t)=\frac1{\pi^2}\sum_{j\ge1}\frac{\cos2\pi jt}{j^2}$, i.e.
$\varkappa_h(\theta)=-\frac{h^2}{2\pi^2}\sum_{j\ge1}\frac{\cos(2Nj\theta)}{j^2}$, uniformly.

\emph{The logarithmic part.} In $L^2(0,\pi)$, $\ln\sin\frac\theta2=-\ln2-\sum_{m\ge1}\frac{\cos m\theta}{m}$ (for $0<\theta<2\pi$, Abel's theorem applied to $-\ln(1-z)=\sum_{m\ge1}z^m/m$ at $z=e^{i\theta}$ gives, on taking real parts, $-\ln|1-e^{i\theta}|=-\ln(2\sin\frac\theta2)=\sum_{m\ge1}\frac{\cos m\theta}{m}$; the functions $\cos m\theta$, $m\ge0$, are orthogonal in $L^2(0,\pi)$ and the coefficients are square summable, so the series also converges in $L^2(0,\pi)$ to the same function). As $\int_0^\pi\varkappa_h=0$ and $\int_0^\pi\cos(2Nj\theta)\cos(m\theta)\,d\theta=\frac\pi2\delta_{m,2Nj}$ for $m,j\ge1$,
\[
\int_0^\pi\varkappa_h(\theta)\ln\sin\frac\theta2\,d\theta=\sum_{j\ge1}\frac1{2Nj}\cdot\frac{h^2}{2\pi^2j^2}\cdot\frac\pi2=\frac{h^2\zeta(3)}{8\pi N}=\frac{\pi\,\zeta(3)}{8N^3}.
\]
\emph{The regular part.} Let $\mathcal K_h(\theta)=\int_0^\theta\varkappa_h$. On a cell, with $t=(\theta-\theta_{k-1})/h$, $\mathcal K_h=h^3\mathcal P(t)$, $\mathcal P(t)=-\frac1{12}t(1-t)(1-2t)$; so $\mathcal K_h$ vanishes at all nodes, $|\mathcal K_h|\le h^3\max|\mathcal P|=\frac{\sqrt3}{216}h^3$, and $\int_0^1\mathcal P=0$. Integrating by parts, $\int_0^\pi\varkappa_hg=-\int_0^\pi\mathcal K_hg'$, hence $|\int_0^\pi\varkappa_hg|\le\frac{\sqrt3}{216}h^3\int_0^\pi|g'|$; moreover $h^{-3}\int_0^\pi\mathcal K_hg'=\int_0^\pi\mathcal P(\{N\theta/\pi\})g'(\theta)\,d\theta\to\int_0^1\mathcal P\cdot\int_0^\pi g'=0$ by the Riemann--Lebesgue (Fej\'er) lemma for $g'\in L^1$.
Finally $\mathcal E_N[f]=-\frac12(f(0)+f(\pi))+h^{-1}(T_N-\int_0^\pi f)$ and $h^{-1}\cdot2\cdot\frac{\pi\zeta(3)}{8N^3}=\frac{\zeta(3)}{4N^2}$; $r_N=h^{-1}\int_0^\pi\varkappa_hg$.
\end{proof}

\begin{lemma}[quantitative comparison with the continuum sums]\label{lem:quant}
Let $N\ge2$, $\kappa\in[0,N]^2$ with $r=|\kappa|\ge1$, $\theta_i=\pi\kappa_i/N$, $x=\frac{\pi\kappa}{2N}$, and $\rho,\phi$ as in Lemma~\ref{lem:small3}; put $\mathcal A=(1+\cos\theta_1)(1+\cos\theta_2)\coth\frac\phi2$, $y=N\phi$, $\mathcal A_\infty=\frac4{|x|}=\frac{8N}{\pi r}$, $P_\infty(\kappa)=\mathcal A_\infty(\coth(\pi r)-1)=\frac{16N}{\pi}F_*(r)$. Then
\begin{equation}\label{eq:compare}
|\mathcal A-\mathcal A_\infty|\le4|x|=\frac{2\pi r}{N},\qquad\kappa_3r\le y\le\pi r,\qquad0\le\pi r-y\le\frac{\pi r|x|^2}{3}=\frac{\pi^3r^3}{12N^2},
\end{equation}
\begin{equation}\label{eq:QQa}
\begin{aligned}\big|P_N(\kappa)-P_\infty(\kappa)\big|&\le\frac{Q(r)}{N},&Q(r)&=2\pi r\big(\coth(\kappa_3r)-1\big)+\frac{2\pi^2r^2}{3\sinh^2(\kappa_3r)},\\
\Big|\frac{\mathcal A}{\sinh y}-\frac{\mathcal A_\infty}{\sinh\pi r}\Big|&\le\frac{Q_a(r)}{N},&Q_a(r)&=\frac{2\pi r}{\sinh(\kappa_3r)}+\frac{2\pi^2r^2\cosh(\kappa_3r)}{3\sinh^2(\kappa_3r)},\end{aligned}
\end{equation}
and for fixed $\kappa$, $N(P_N(\kappa)-P_\infty(\kappa))\to p(\kappa)$ and $N\big(\frac{\mathcal A}{\sinh y}-\frac{\mathcal A_\infty}{\sinh\pi r}\big)\to p_a(\kappa)$ as $N\to\infty$, with $p,p_a$ from \eqref{eq:ppa}; in particular $|p|\le Q$, $|p_a|\le Q_a$.
Consequently, with $D_N$, $e_N$, $\mathrm{Alt}_N$, $E_\infty$ from Lemma~\ref{lem:small3} and $A_\infty:=\frac2\pi\sum_{k\in\Z^2\setminus0}\frac{(-1)^{k_1+k_2}}{|k|\sinh(\pi|k|)}=\frac{1-c_{m\tau3}}{3}$,
\begin{equation}\label{eq:deltaN}
\begin{gathered}
\delta_N:=D_N+\sum_{k=1}^{N-1}e_N(\theta_k)-E_\infty N,\qquad|\delta_N|\le\frac{c_\delta}{N},\qquad N\delta_N\to\delta_\infty;\\
\big|\mathrm{Alt}_N-A_\infty N\big|\le\frac{c_a}{N},\qquad N(\mathrm{Alt}_N-A_\infty N)\to\alpha_\infty,
\end{gathered}
\end{equation}
where (certified values) $c_\delta=\frac14\sum_{k\in\Z^2\setminus0}Q(|k|)+\sum_{j\ge1}\int_0^\infty Q(\sqrt{j^2+v^2})\,dv+2\cdot10^{-4}<4.864$ and $c_a=\frac14\sum_{k\in\Z^2\setminus0}Q_a(|k|)+0.08<27.13$.
\end{lemma}
\begin{proof}
With $s_i=\sin x_i$, $x_i\in[0,\frac\pi2]$: $x_i^2(1-\frac{x_i^2}{3})\le s_i^2\le x_i^2$ (as in Lemma~\ref{lem:elem}(a); $\sin v\ge v-\frac{v^3}6\ge0$ on $[0,\frac\pi2]$), hence $|x|^2(1-\frac{|x|^2}3)\le\rho^2\le|x|^2$, and $\rho\ge\frac2\pi|x|=\frac rN$.
\emph{$\mathcal A$:} $\mathcal A=4(1-s_1^2)(1-s_2^2)\sqrt{1+\rho^2}/\rho\le4/\rho$ (Lemma~\ref{lem:small3}), and $\frac1\rho-\frac1{|x|}=\frac{|x|^2-\rho^2}{\rho|x|(|x|+\rho)}\le\frac{|x|^2}{3\rho}\le\frac{\pi|x|}{6}$, so $\mathcal A-\mathcal A_\infty\le\frac{2\pi}3|x|$. Conversely $(1-s_1^2)(1-s_2^2)\ge1-\rho^2$, so $\mathcal A\ge\frac{4(1-\rho^2)}{\rho}$ (trivially if $\rho>1$), and $\frac4\rho-4\rho\ge\frac4{|x|}-4|x|$. Thus $|\mathcal A-\mathcal A_\infty|\le4|x|$.
\emph{$y$:} $y=2N\operatorname{arsinh}\rho\le2N|x|=\pi r$, and $y\ge\kappa_3N\rho\ge\kappa_3r$ by \eqref{eq:Pbound}. If $|x|^2\ge3$ the last bound in \eqref{eq:compare} is trivial. Otherwise $\operatorname{arsinh}\rho\ge\rho(1-\frac{\rho^2}6)$ (compare derivatives: $(1+\rho^2)^{-1/2}\ge1-\frac{\rho^2}{2}$), the function $\rho\mapsto\rho(1-\frac{\rho^2}6)$ is increasing on $[0,\sqrt2]$, and $\rho\ge\rho_-:=|x|(1-\frac{|x|^2}{3})^{1/2}$ with $1-\frac{\rho_-^2}{6}\ge1-\frac{|x|^2}{6}\ge(1-\frac{|x|^2}{3})^{1/2}$; so $\operatorname{arsinh}\rho\ge|x|(1-\frac{|x|^2}3)$, i.e.\ $\pi r-y\le\pi r|x|^2/3$.
\emph{\eqref{eq:QQa}:} With $F_c(z)=\coth z-1$ (positive, decreasing, $|F_c'|=\sinh^{-2}$ decreasing), $P_N-P_\infty=(\mathcal A-\mathcal A_\infty)F_c(y)+\mathcal A_\infty(F_c(y)-F_c(\pi r))$, so by \eqref{eq:compare} and the mean value theorem $|P_N-P_\infty|\le\frac{2\pi r}{N}F_c(\kappa_3r)+\frac{8N}{\pi r}\cdot\frac{\pi^3r^3}{12N^2}\sinh^{-2}(\kappa_3r)=Q(r)/N$. The second line is the same computation with $1/\sinh z$ in place of $F_c$ ($|\frac{d}{dz}\frac1{\sinh z}|=\frac{\cosh z}{\sinh^2z}$, decreasing).
\emph{Pointwise limits:} for fixed $\kappa$, as $N\to\infty$ ($x\to0$), with $q_4=x_1^4+x_2^4$: $\rho^2=|x|^2-\frac{q_4}3+O(|x|^6)$, $(1-s_1^2)(1-s_2^2)=1-|x|^2+O(|x|^4)$, $\sqrt{1+\rho^2}=1+\frac{|x|^2}2+O(|x|^4)$, hence
$\mathcal A-\mathcal A_\infty=\frac{4}{|x|}\big(\frac{q_4}{6|x|^2}-\frac{|x|^2}{2}\big)+O(|x|^3)=\frac{2\pi}{Nr}\chi_-(\kappa)+O(N^{-3})$ and
$\pi r-y=\pi r\big(\frac{q_4}{6|x|^2}+\frac{|x|^2}6\big)+O(N|x|^5)=\frac{\pi^3r}{4N^2}\chi_+(\kappa)+O(N^{-4})$. Therefore
$N(P_N-P_\infty)\to\frac{2\pi}{r}\chi_-F_c(\pi r)+\frac{8}{\pi r}\cdot\frac{\pi^3r}{4}\chi_+\sinh^{-2}(\pi r)=p(\kappa)$, and likewise for $p_a$.
\emph{\eqref{eq:deltaN}:} Recall $D_N=\sum'\varepsilon'_{k_1}\varepsilon'_{k_2}P_N(k)$ and $e_N(\theta_k)=\int_0^NP_N(k,v)\,dv$ (proof of Lemma~\ref{lem:small3}), and $E_\infty N=\sum'_{k\in\mathbb N_0^2\setminus0}\varepsilon'\varepsilon'P_\infty(k)+\sum_{j\ge1}\int_0^\infty P_\infty(j,v)\,dv$. Hence
\[
|\delta_N|\le\frac1N\Big[\frac14\sum_{k\in\Z^2\setminus0}Q(|k|)+\sum_{j\ge1}\int_0^\infty Q\big(\sqrt{j^2+v^2}\big)dv\Big]+T_1+T_2,
\]
where
\begin{align*}
T_1&=\frac{4N}\pi\sum_{|k|_\infty\ge N}F_*(|k|)\le\frac{4N}{\pi}\sum_{j\ge N}\frac{8}{e^{2\pi j}-1}\le\frac{32N\,e^{-2\pi N}}{\pi(1-e^{-2\pi})(1-e^{-4\pi})},\\
T_2&=\frac{16N}{\pi}\Big[\sum_{j\ge N}\int_0^\infty F_*\big(\sqrt{j^2+v^2}\big)dv+\sum_{j=1}^{N-1}\int_N^\infty F_*\big(\sqrt{j^2+v^2}\big)dv\Big]\\
&\le\frac{16N}{\pi}e^{-2\pi N}\Big[\frac{1}{2\sqrt N(1-e^{-2\pi})(1-e^{-4\pi})}+\frac{1}{2\pi(1-e^{-4\pi})}\Big]
\end{align*}
(for the first sum use the computation of Lemma~\ref{lem:small3} with $\pi$ in place of $\kappa_3$; for the second $F_*(\sqrt{j^2+v^2})\le F_*(v)\le\frac{e^{-2\pi v}}{N(1-e^{-4\pi})}$). As $N^2e^{-2\pi N}$ is decreasing, $N(T_1+T_2)\le1.8\cdot10^{-4}$ for $N\ge2$. The limit $N\delta_N\to\delta_\infty$ follows by dominated convergence (dominating functions $Q(|k|)$ and $Q(\sqrt{j^2+v^2})$, which do not depend on $N$; the tails $N(T_1+T_2)\to0$). The statements on $\mathrm{Alt}_N$ are proved in the same way; its tail is $\frac{2N}{\pi}\sum_{|k|_\infty\ge N}\frac{1}{|k|\sinh\pi|k|}\le\frac{32N\,e^{-\pi N}}{\pi(1-e^{-\pi})(1-e^{-4\pi})}$, at most $0.08/N$ for $N\ge2$. $A_\infty=(1-c_{m\tau3})/3$ is the definition \eqref{eq:c3consts} of $c_{m\tau3}$.
The numerical values: $Q$ and $Q_a$ are summed over $|k|_\infty\le30$ in ball arithmetic; for $r\ge1$, $Q(r)\le e^{-2\kappa_3r}(a_1r+a_2r^2)$ and $Q_a(r)\le e^{-\kappa_3r}(a_3r+a_4r^2)$ with $a_1=a_3=\frac{4\pi}{1-q}$, $a_2=\frac{8\pi^2}{3(1-q)^2}$, $a_4=\frac{4\pi^2(1+q)}{3(1-q)^2}$, $q=e^{-2\kappa_3}$, which bounds the tails, and with $v=j\sinh t$,
$\int_0^\infty e^{-2\kappa_3\sqrt{j^2+v^2}}(a_1r+a_2r^2)\,dv=a_1j^2\frac{K_0(z)+K_2(z)}{2}+a_2j^3\frac{3K_1(z)+K_3(z)}{4}$, $z=2\kappa_3j$ ($K_\nu$ the modified Bessel functions, $K_\nu(z)=\int_0^\infty e^{-z\cosh t}\cosh\nu t\,dt$, DLMF 10.32.9).
\end{proof}

\subsection{Proof of Theorem \ref{thm:d3}(c),(d)}
By \eqref{eq:tau3identity}, \eqref{eq:deltaN} and the algebra $A_0+\frac3\pi(\pi+2-6\ln2)+3E_\infty=2+\frac{12}\pi(\gamma-\ln\pi-2\ln2)+3E_\infty=c_{\tau3}$ (insert $I_\Psi=-\frac12+\frac32\ln2-\ln\pi$),
\begin{equation}\label{eq:R3}
R_3(N)=\tau_3-C_3N^3-c_{\tau3}N^2=\frac{3N^2}{\pi}\big(\varepsilon_b(N)-\varepsilon_b(\infty)\big)+3N\delta_N+\rho_N,\qquad\varepsilon_b(\infty):=\pi+2-6\ln2 .
\end{equation}
By Lemma~\ref{lem:Jb2}, $f=J_b$ satisfies the hypotheses of Lemma~\ref{lem:traplog} with $\int|g'|\le V_g$, and $-\frac12(J_b(0)+J_b(\pi))=-\frac12g_1(0)=\varepsilon_b(\infty)$; so
\[
\varepsilon_b(N)=\varepsilon_b(\infty)+\frac{\zeta(3)}{4N^2}+r_N,\qquad\frac{3N^2}{\pi}|r_N|\le\frac{\sqrt3\,\pi}{72}V_g<0.527 .
\]
With $|3N\delta_N|\le3c_\delta<14.60$ and $\rho_N\in(-2.536,3.900)$ this gives
$R_3\ge\frac{3\zeta(3)}{4\pi}-0.527-14.60-2.536>-17.4$ and $R_3\le\frac{3\zeta(3)}{4\pi}+0.527+14.60+3.900<19.4$ ($\frac{3\zeta(3)}{4\pi}=0.28696\ldots$).
By \eqref{eq:m3exact}, $m_3-\tau_3=N^2-1-3N\mathrm{Alt}_N=(1-3A_\infty)N^2-1-3N(\mathrm{Alt}_N-A_\infty N)$, so $R_3'=-1-3N(\mathrm{Alt}_N-A_\infty N)$ and $|R_3'+1|\le3c_a<81.4$. This proves (c).

For (d): $\frac{3N^2}{\pi}(\varepsilon_b(N)-\varepsilon_b(\infty))\to\frac{3\zeta(3)}{4\pi}$ (Lemma~\ref{lem:traplog}), $3N\delta_N\to3\delta_\infty$ and $R_3'\to-1-3\alpha_\infty$ (Lemma~\ref{lem:quant}). It remains to show $\rho_N=1-\frac{12N^2}{\pi}\eta_N+6Ne_1(N)\to1+\frac\pi{12}$. By Lemma~\ref{lem:harm}, $N^2\eta_N\to\frac1{12}$. For $e_1(N)$ (Lemma~\ref{lem:trap}; step $h=\frac{\pi}{2N}$ on $[0,\frac\pi2]$) the cell formula of Lemma~\ref{lem:traplog} gives $h\,e_1(N)=\int_0^{\pi/2}q_h\Psi''=\frac{h^2}{12}\big(\Psi'(\frac\pi2)-\Psi'(0)\big)+\int_0^{\pi/2}\varkappa_h\Psi''$, and $|\int\varkappa_h\Psi''|=|\sum_{\mathrm{cells}}\int\varkappa_h(\Psi''-\Psi''(\mathrm{mid}))|\le\frac{h^2}{12}\cdot\frac\pi2\cdot\omega_{\Psi''}(h)=o(h^2)$ because $\Psi''$ is continuous (Lemma~\ref{lem:Psi}(a)). By Lemma~\ref{lem:Psi}(a), $\Psi'(x)=c'(x)+Q'(\sin x)\cos x$ with $c'(0)=\frac16$, $Q'(0)=-\frac12$, $c'(\frac\pi2)=\frac{4}{\pi^2}$, so $\Psi'(\frac\pi2)-\Psi'(0)=\frac4{\pi^2}+\frac13$ and $6Ne_1(N)\to6\cdot\frac{\pi}{24}\big(\frac{4}{\pi^2}+\frac13\big)=\frac1\pi+\frac\pi{12}$. Hence $\rho_N\to1-\frac1\pi+\frac1\pi+\frac\pi{12}$, and $R_3(N)\to\frac{3\zeta(3)}{4\pi}+3\delta_\infty+1+\frac\pi{12}$.
The values of $\delta_\infty$, $\alpha_\infty$ are certified in \texttt{16\_certify\_d3\_remainder.py} (lattice sums over $|k|_\infty\le12$, rigorous integration of $p(j,\cdot)$, $j\le12$, on $[0,14]$, and tails bounded through $|p(\kappa)|\le\pi r(\coth\pi r-1)+\frac{2\pi^2r^2}{3\sinh^2\pi r}$, $|p_a(\kappa)|\le\frac{\pi r}{\sinh\pi r}+\frac{2\pi^2r^2\cosh\pi r}{3\sinh^2\pi r}$, which follow from $|\chi_-|\le\frac{r^2}{2}$, $0\le\chi_+\le\frac{r^2}3$). \qed

\begin{remark}[how sharp is (c)?]\label{rem:d3sharp}
The true remainders are much smaller than the bounds in (c): $R_3(N)$ increases from $1.6117$ ($N=2$) to its limit $1.75740$ and $R_3'(N)$ from $-0.6263$ to $-0.29892$ (script \texttt{17}; in agreement with (d) to all digits shown). The loss is in the crude comparison constants $c_\delta$, $c_a$ of Lemma~\ref{lem:quant} (the true values of $N|\delta_N|$ and $N|\mathrm{Alt}_N-A_\infty N|$ are below $0.07$ and $0.24$), which use the worst-case exponent $\kappa_3=1.62$ instead of $\pi$, and in the bound for $\rho_N$. The limits in (d) are proved without a rate.
\end{remark}


\section{Explicit consequences in two and three dimensions}\label{sec:explicit}

Theorem~\ref{thm:cc} expresses every bound through the single number $\Xb=\Lambda_1\tau$ (the ratio of the stationary-start hitting time to the relaxation time of the visible sector). Theorems~\ref{thm:d2} and~\ref{thm:d3} locate $\Xb$.

\begin{corollary}[$\Xb$ and $\Lambda_1(m-\tau)$ in $d=2$ and $d=3$]\label{cor:Xbar}
For every $N\ge2$, $X_-\le\Xb\le X_+$ and $M_-\le\Lambda_1(m-\tau)\le M_+$, where, with $a_N=1-\frac{\pi^2}{12N^2}$, $\mu_-=\mu_-(N)$ and $\ell_W=1+\ln W$ from Theorem~\ref{thm:cc}:
\begin{align*}
d=2:\quad&X_-=a_N\Big(2\pi\ln N+\tfrac{\pi^2}4c_\tau-\tfrac{\pi^2}{4}\tfrac{2.14}{N^2}\Big),&&X_+=2\pi\ln N+\tfrac{\pi^2}4c_\tau+\tfrac{\pi^2}4\tfrac{3.53}{N^2},\\
&M_-=\max\Big\{\mu_-,\ a_N\Big(\pi\ln2-\tfrac{\pi^2}{4}\tfrac{9.06}{N^2}\Big)\Big\},&&M_+=\min\Big\{\ell_W,\ \pi\ln2+\tfrac{\pi^2}{4}\tfrac{7.72}{N^2}\Big\},\\
d=3:\quad&X_-=a_N\,\tfrac{\pi^2}{6}\Big(C_3N+c_{\tau3}-\tfrac{17.4}{N^2}\Big),&&X_+=\tfrac{\pi^2}6\Big(C_3N+c_{\tau3}+\tfrac{19.4}{N^2}\Big),\\
&M_-=\max\Big\{\mu_-,\ a_N\,\tfrac{\pi^2}{6}\Big(c_{m\tau3}-\tfrac{82.4}{N^2}\Big)\Big\},&&M_+=\min\Big\{\ell_W,\ \tfrac{\pi^2}{6}\Big(c_{m\tau3}+\tfrac{80.4}{N^2}\Big)\Big\}.
\end{align*}
In particular, with the exact constants $\frac{\pi^2}{4}c_\tau=-1.7960326\ldots$, $\frac{\pi^2}6C_3=7.1068721\ldots$, $\frac{\pi^2}{6}c_{\tau3}=-8.9136329\ldots$, $\frac{\pi^2}6c_{m\tau3}=2.5193561\ldots$,
\[
\begin{aligned}d=2:\quad&\Big|\Xb-\big(2\pi\ln N+\tfrac{\pi^2}{4}c_\tau\big)\Big|\le\frac{8.8+5.2\ln N}{N^2},&&\big|\Lambda_1(m_2-\tau_2)-\pi\ln2\big|\le\frac{24.2}{N^2};\\
d=3:\quad&-\frac{5.85}{N}-\frac{21.3}{N^2}\le\Xb-\tfrac{\pi^2}{6}\big(C_3N+c_{\tau3}\big)\le\frac{31.92}{N^2},&&\big|\Lambda_1(m_3-\tau_3)-\tfrac{\pi^2}{6}c_{m\tau3}\big|\le\frac{137.7}{N^2}.\end{aligned}
\]
(The exact constants must not be replaced by truncated decimals if these bounds are used for very large $N$. For $d=3$ and small $N$ the elementary bounds \eqref{eq:tau3} may be sharper than $X_\pm$.)
\end{corollary}
\begin{proof}
$\Lambda_1=\frac2d\sin^2u$ with $u=\frac\pi{2N}\le\frac\pi4$, and $u^2(1-\frac{u^2}{3})\le\sin^2u\le u^2$ (Lemma~\ref{lem:elem}(a)), i.e.\ $a_N\frac{\pi^2}{2dN^2}\le\Lambda_1\le\frac{\pi^2}{2dN^2}$; multiply by the bounds \eqref{eq:tau2}, \eqref{eq:m2}, \eqref{eq:tau3sharp} for $\tau$ and $m-\tau$ (as $\tau>0$ and $m-\tau>0$, the lower bound for $\Lambda_1$ may be multiplied by a lower bound for $\tau$ or $m-\tau$ even if the latter is negative), and use $\frac{\pi^2}{4}c_{m\tau}=\pi\ln2$ (Lemma~\ref{lem:closed}) and \eqref{eq:muminus}, \eqref{eq:mtau}. The simplified forms follow from $(1-\frac{a}{N^2})(B-\frac b{N^2})\ge B-\frac{b+a\max(B,0)}{N^2}$ for $a,b\ge0$, with $a=\frac{\pi^2}{12}=0.8225$:
for $d=2$, $B=2\pi\ln N+\frac{\pi^2}4c_\tau\le2\pi\ln N$, $b=\frac{\pi^2}{4}\cdot2.14<5.281$, so $b+aB\le3.81+5.17\ln N$, and $\frac{\pi^2}{4}\cdot3.53<8.711$; for $\Lambda_1(m_2-\tau_2)$: $B=\pi\ln2$, $b=\frac{\pi^2}{4}\cdot9.06<22.36$, $b+aB<24.15$, and $\frac{\pi^2}4\cdot7.72<19.05$;
for $d=3$, $B=\frac{\pi^2}6(C_3N+c_{\tau3})>0$, $b=\frac{\pi^2}{6}\cdot17.4<28.63$, $\frac{aB}{N^2}\le\frac{5.8452}{N}-\frac{7.331}{N^2}$, and $\frac{\pi^2}6\cdot19.4<31.92$; for $\Lambda_1(m_3-\tau_3)$: $B=\frac{\pi^2}{6}c_{m\tau3}$, $b=\frac{\pi^2}6\cdot82.4<135.55$, $b+aB<137.62$, and $\frac{\pi^2}{6}\cdot80.4<132.26$ (all products certified in scripts \texttt{09}, \texttt{16}).
\end{proof}

\begin{corollary}[explicit bounds in terms of $N$]\label{cor:explicitN}
Let $d\in\{2,3\}$, let $Z$ be any upper bound for $Z_d$ (for instance $6.02682$ resp.\ $16.53232$ by \eqref{eq:Zvalues}, or the sharper certified values \eqref{eq:Zcert}, which are used in Tables~\ref{tab:cc2}--\ref{tab:cc3x}), and $X_\pm$, $M_\pm$ as in Corollary~\ref{cor:Xbar}; assume $X_->1$ and put $\hat\beta=X_-/(X_--1)$ (the hypothesis $X_->1$ holds for every $N\ge3$, since both factors of $X_-$ are positive and increasing in $N$ for $N\ge3$ and $X_-(3)=4.10\ldots$ resp.\ $8.38\ldots$; for $N=2$ it fails in both dimensions, $X_-=0.984\ldots$ for $d=2$ and $X_-<0$ for $d=3$, although $\Xb=\frac52$ resp.\ $\frac{29}6$, so the corollary is void for $N=2$, where Theorem~\ref{thm:cc} with the exact $\Xb$ applies). Then all statements of Theorem~\ref{thm:cc} hold with $\Xb$, $\mu_-$, $\ell_W$ replaced as follows:
\begin{enumerate}
\item[(i)] $t_-:=\dfrac{2}{1+\sqrt{1+4\hat\beta Z/X_-^2}}\le\sigma_0\tau\le t_+:=\dfrac{2}{1+\sqrt{1+4W/X_+^2}}$;
\item[(ii)] $\dfrac{M_-}{X_+}\le\mu\le\dfrac{M_+}{X_-}$ and $\Big(1+\dfrac{M_-}{X_+}\Big)t_-\le\sigma_0m\le\Big(1+\dfrac{M_+}{X_-}\Big)t_+$;
\item[(iii)] $\delta\le\delta_+:=W/(X_--W-1+\vartheta_0)$ if the denominator is positive; if $\delta_+<1$, $\omega:=2(1+\delta_+)/(1-\delta_+)$ and $\delta\ge\delta_-:=W/(X_+-W+\omega(Z-W)-\frac{1}{1+\delta_+})$;
\item[(iv)] $\Lambda_1(1+\delta_--\frac1{X_-})\le\sigma_1-\sigma_0\le\Lambda_1(1+\delta_+-\frac{1}{1+X_+})$;
\item[(v)] $\Big(1+\dfrac{M_-}{X_+}\Big)\Big(1-\dfrac{\hat\beta(1+\hat\beta)Z}{X_-^2}\Big)\le r_0\le1+\dfrac{M_+}{X_-}+\dfrac{\hat\beta Z}{X_-^2}$ (the lower bound if $\hat\beta(1+\hat\beta)Z\le X_-^2$);
\item[(vi)] the bounds of Theorem~\ref{thm:cc}(vi) for $r_1$, computed with these $\delta_\pm$, $\omega$ and with $y_-=W-M_+-\omega(Z-W)-1$, $y_+=W-M_-+\omega(Z-W)-\frac{1}{1+\delta_+}$.
\end{enumerate}
\end{corollary}
\begin{proof}
Each bound of Theorem~\ref{thm:cc} is monotone in $\Xb$: the lower bounds in (i), (v) increase with $\Xb$ ($\beta$ decreases), the upper bound in (i) increases, and $\mu=\Lambda_1(m-\tau)/\Xb\in[M_-/X_+,M_+/X_-]$, so that $\sigma_0m=(1+\mu)\sigma_0\tau$ and \eqref{eq:r0} give (ii) and (v) exactly as in the proof of Theorem~\ref{thm:cc} (there $\mu_-\le\Lambda_1(m-\tau)\le\ell_W$ was used; here $M_-\le\Lambda_1(m-\tau)\le M_+$, in particular $Y_{x_0}(0)=W-\Lambda_1(m-\tau)\in[W-M_+,W-M_-]$); $\delta_+$ decreases with $\Xb$, and by Remark~\ref{rem:monotone}(b),(c) the lower bound for $\delta$ and the bounds for $r_1$ remain valid when $\delta_+$ (hence $\omega$) is replaced by a larger number and $\Xb$ in $\delta_-$ by a larger number.
\end{proof}

\begin{theorem}[first-order pole structure; $d=2,3$]\label{thm:asym}
For the corner-to-corner geometry in $d\in\{2,3\}$ let
\[
Z_d=\sum_{k\in\Z^d\setminus0}\frac{1}{|k|^4},\qquad Z_d'=\sum_{k\in\Z^d,\ |k|^2\ge2}\frac{1}{|k|^2(|k|^2-1)},\qquad
\mu_\infty=\begin{cases}\pi\ln2=2.17758\ldots&(d=2),\\[2pt] \frac{\pi^2}{6}c_{m\tau3}=2.51935\ldots&(d=3).\end{cases}
\]
($Z_2=6.02681\ldots$, $Z_3=16.53231\ldots$, $Z_2'=3.22591\ldots$, $Z_3'=14.70229\ldots$, certified in Lemma~\ref{lem:epstein}.) Then, as $N\to\infty$ (so that $\Xb\to\infty$ by Corollary~\ref{cor:Xbar}):
\begin{enumerate}
\item[(a)] $\Lambda_1^2K_2\to Z_d$ and $\ \sigma_0\tau=1-\dfrac{Z_d}{\Xb^2}+o(\Xb^{-2})$.
\item[(b)] $\Lambda_1(m-\tau)\to\mu_\infty$, so that $X:=\Lambda_1m=\Xb+\mu_\infty+o(1)$ and
$\ \sigma_0\,m=1+\dfrac{\mu_\infty}{\Xb}+o(\Xb^{-1})$. In both dimensions $\Lambda_1(m-\tau)=\mu_\infty+O(N^{-2})$ with the explicit constants of Corollary~\ref{cor:Xbar}; for $d=2$, $\sigma_0m_2-1=\dfrac{\ln2}{2\ln N}\,(1+O(1/\ln N))$; for $d=3$, $\Xb=\frac{\pi^2}{6}(C_3N+c_{\tau3})+O(N^{-1})$ and $\sigma_0m_3-1=\dfrac{\mu_\infty}{\frac{\pi^2}{6}C_3N}\big(1+O(N^{-1})\big)\approx\dfrac{0.3545}{N}\big(1+O(N^{-1})\big)$.
\item[(c)] $\dfrac{\sigma_1}{\Lambda_1}=1+\delta$ with $\ \dfrac W\delta=\Xb-(2d+1-Z_d')+o(1)$; in particular $\delta=\dfrac{2d}{\Xb}\big(1+O(\Xb^{-1})\big)$.
\item[(d)] $\dfrac{\sigma_1-\sigma_0}{\Lambda_1}=1+\dfrac{2d-1}{\Xb}+O(\Xb^{-2})$.
\item[(e)] $r_0=1+\dfrac{\mu_\infty}{\Xb}+o(\Xb^{-1})$, $\quad r_1=-\dfrac{2d}{\Xb}\big(1+O(\Xb^{-1})\big)$, $\quad\dfrac{r_1\sigma_1}{r_0\sigma_0}\to-2d$.
\end{enumerate}
All $O(\cdot)$ terms are explicit through Theorem~\ref{thm:cc} and Corollary~\ref{cor:explicitN}; the $o(\cdot)$ terms come from dominated convergence.
\end{theorem}
\begin{proof}
(a) $\Lambda_1^2K_2=\sum_{k\ne0}w_k(\Lambda_1/\lambda_k)^2$, where for fixed $k$, $w_k\to\prod_i\varepsilon_{k_i}$ and $\lambda_k/\Lambda_1=\sum_i\sin^2(k_iu)/\sin^2u\to|k|^2$ as $N\to\infty$, and each term is bounded by $\prod_i\varepsilon_{k_i}|k|^{-4}$ (proof of Lemma~\ref{lem:K2}), which is summable for $d\le3$. So $\kappa\Xb^2\to Z_d$; now use \eqref{eq:s0sharp}, $1-\beta\kappa\le\sigma_0\tau\le1-\kappa+2\kappa^2$, with $\beta=1+O(\Xb^{-1})$.
(b) Corollary~\ref{cor:Xbar} (which rests on \eqref{eq:m2}, Lemma~\ref{lem:closed} and Theorem~\ref{thm:d3}(c)). Then $\sigma_0m=(1+\mu)\sigma_0\tau$ with $\mu=\Lambda_1(m-\tau)/\Xb$, and (a). For $d=2$, $\Xb=2\pi\ln N+O(1)$.
(c) By Theorem~\ref{thm:cc}(iii), $\delta\le\delta_+\to0$. By \eqref{eq:deltaexact}, $W/\delta=Y(\sigma_1)-\frac{1}{1+\delta}$, $Y(\sigma_1)=Y_0+[Y(\sigma_1)-Y_0]$, $Y_0=\Xb-W$, $W\to2d$, and
\[
Y(\sigma_1)-Y_0=\sum_{k:\ |k|^2\ge2}w_k\,\frac{\Lambda_1\sigma_1}{\lambda_k(\lambda_k-\sigma_1)} .
\]
For fixed $k$ the summand tends to $\prod_i\varepsilon_{k_i}/(|k|^2(|k|^2-1))$. Once $\delta_+\le\frac12$ we have $\sigma_1\le\frac32\Lambda_1\le\frac34\lambda_k$ for these $k$ (as $\lambda_k\ge2\Lambda_1$), so the summand is at most $6\,w_k(\Lambda_1/\lambda_k)^2\le6\prod_i\varepsilon_{k_i}|k|^{-4}$, and dominated convergence gives $Y(\sigma_1)-Y_0\to Z_d'$.
(d) $\sigma_0/\Lambda_1=\sigma_0\tau/\Xb=\Xb^{-1}+O(\Xb^{-3})$ by (a), and (c).
(e) Theorem~\ref{thm:cc}(v) gives $r_0=1+\mu+O(\Xb^{-2})$. In \eqref{eq:r1exact} with $A_{x_0}=W$, the quantities $Y_{x_0}(\sigma_1)-\varsigma$ and $\varsigma+(1+\delta)n\Lambda_1^2h'(\sigma_1)$ are bounded (Theorem~\ref{thm:r1size} and Remark~\ref{rem:monotone}), so $-r_1=\delta(1+O(\delta))$. Finally $\frac{r_1\sigma_1}{r_0\sigma_0}=\frac{r_1}{r_0}(1+\delta)\frac{\Xb}{\sigma_0\tau}\to-2d$.
\end{proof}

\begin{remark}[comparison with the formal perturbation theory of the computations of Sections~5--6 of the article]\label{rem:floatstudy}
In the notation of the computations of Sections~5--6 of the article $\nu_j=\sigma_j$, $a_j=r_j\sigma_j$, $\mu_1=\Lambda_1$, $X=\Lambda_1\cdot\mathrm{MFPT}=\Xb+\mu_\infty+o(1)$, $A=W$. Theorem~\ref{thm:asym} proves the four relations
``$\nu_0=\mathrm{MFPT}^{-1}(1+O(1/X))$, $\nu_1=\mu_1(1+W/X+O(X^{-2}))$, $a_0\approx\nu_0$, $a_1\approx-A\nu_0$'' and identifies the constants: $\nu_0\cdot\mathrm{MFPT}=1+\mu_\infty/X+\dots$ with $\mu_\infty=\pi\ln2$ ($d=2$) or $2.519$ ($d=3$), and $(\nu_1-\nu_0)/\mu_1=1+(2d-1)/X+\dots$, which is the rate ``$\mu_1(X+W-1)/X$'' underlying the closed form $\mathrm{mode}\approx\mathrm{MFPT}\,\ln(2dX)/(X+2d-1)$. The numerical values of Sections~5--6 of the article ``$\nu_0\cdot\mathrm{MFPT}=1.164\to1.022$ (2D)'', ``$1.076\to1.0003$ (3D)'' and ``$a_1/a_0\to-4.09,\,-5.998$'' agree with (b) and (e). What Theorem~\ref{thm:asym} does \emph{not} give is a theorem about the mode itself; see Section~\ref{sec:gaps}.
\end{remark}

\begin{proposition}[the constants of Corollary~\ref{cor:Xbar} may not be truncated]\label{prop:roundedfalse}
If the ``in particular'' part of Corollary~\ref{cor:Xbar} is stated with truncated decimal constants, the two statements become false:
\begin{enumerate}
\item[(a)] ($d=2$) The inequality $\big|\Xb-(2\pi\ln N-1.796034)\big|\le\dfrac{8.8+5.2\ln N}{N^2}$ is violated for every $N\ge8872$; more precisely
$\Xb>2\pi\ln N-1.796034+\dfrac{8.8+5.2\ln N}{N^2}$ for all $N\ge8872$. By direct evaluation it holds at $N=8539$ (and at $N=100$, $1000$, $7000$, $8000$) and is violated at $N=8540$, $10^4$, $2\cdot10^4$, $10^5$.
\item[(b)] ($d=3$) The inequality $7.1069\,N-14.0-\dfrac{5.85}{N}-\dfrac{4.18}{N^2}\le\Xb$ is violated for every $N\ge182\,420$.
\end{enumerate}
The admissible statements are those of Corollary~\ref{cor:Xbar} (exact constants $\frac{\pi^2}4c_\tau$, $\frac{\pi^2}{6}C_3$, $\frac{\pi^2}6c_{\tau3}$); they hold for every $N\ge2$.
\end{proposition}
\begin{proof}
(Computer-assisted through Corollary~\ref{cor:Xbar} and the certified constants; script~\texttt{20}.)
(a) Put $\beta=\frac{\pi^2}{4}c_\tau=-1.7960326424\ldots$, $\Delta=1.796034+\beta\in(1.3574\cdot10^{-6},1.3576\cdot10^{-6})$, $b=\frac{\pi^2}4\cdot2.14$ and $B=2\pi\ln N+\beta>0$. By Corollary~\ref{cor:Xbar}, $\Xb\ge X_-=(1-\frac{\pi^2}{12N^2})(B-\frac b{N^2})\ge B-\frac{b+\frac{\pi^2}{12}B}{N^2}$. Hence
\[
N^2\Big[\Xb-\Big(2\pi\ln N-1.796034+\frac{8.8+5.2\ln N}{N^2}\Big)\Big]\ \ge\ h_0(N):=\Delta N^2-b-\frac{\pi^2}{12}\beta-\Big(\frac{\pi^3}6+5.2\Big)\ln N-8.8 .
\]
$h_0'(N)=2\Delta N-(\frac{\pi^3}6+5.2)/N>0$ for $N\ge1955$, and $h_0(8872)>0$ (certified; $h_0(8871)<0$). So $h_0(N)>0$ for all $N\ge8872$. The direct evaluations use the single sum \eqref{eq:tau2exact} in ball arithmetic; at each of these $N$ the inequality of Corollary~\ref{cor:Xbar} is verified as well.
(b) By Corollary~\ref{cor:Xbar}, $\Xb\le X_+=\frac{\pi^2}{6}(C_3N+c_{\tau3}+\frac{19.4}{N^2})$, so the lower bound with truncated constants minus $\Xb$ is at least
\[
g(N):=\Big(7.1069-\frac{\pi^2}{6}C_3\Big)N-14.0-\frac{\pi^2}{6}c_{\tau3}-\frac{5.85}{N}-\frac{4.18+\frac{\pi^2}{6}19.4}{N^2},
\]
which is increasing in $N$ because $7.1069-\frac{\pi^2}6C_3=2.788\ldots\cdot10^{-5}>0$; and $g(182\,420)>0$ (certified; $g(182\,419)<0$).
\end{proof}
\begin{remark}
In $d=2$ the statement (a) is confirmed by explicit counter-examples that do not use Corollary~\ref{cor:Xbar} (e.g.\ $N=10^4$: $|\Xb-(2\pi\ln N-1.796034)|>9.18\cdot10^{-7}$, while $(8.8+5.2\ln N)/N^2<5.67\cdot10^{-7}$). In $d=3$ no direct counter-example was computed (at $N\approx1.8\cdot10^5$ the double sum \eqref{eq:tau3exact} has $3\cdot10^{10}$ terms); (b) is a consequence of the upper bound $\Xb\le X_+$, which is proved for all $N$ but was verified by direct evaluation only for $N\le800$ (script \texttt{17}; separately written code of a check, within the same project, reached $N=1600$). The lesson is the parenthetical warning in Corollary~\ref{cor:Xbar}: an $O(N^{-2})$ error term cannot absorb a constant that is wrong in the sixth decimal.
\end{remark}

\part*{Part IV. Neglected poles, certification, open points}
\section{The neglected poles: a rigorous tail bound}\label{sec:tail}

A mode theorem based on the two slowest poles needs a bound on the contribution of all the others. The following bound is valid for every finite $N$.

\begin{theorem}[tail of the spectral expansion]\label{thm:tail}
Assume \eqref{eq:A}, $x\in\Omega'$, and let $1\le J\le p-1$. Write $S_x(t)=\PP_x(T>t)$, $f_x=-S_x'$ for the unit-rate walk in continuous time, $p_t(x,y)$ for the transition kernel without target, and $\varrho_j$ as in \eqref{eq:varrho}. Then for all $t_0>0$:
\begin{align}
\Big|S_x(t)-\sum_{j<J}r_j(x)e^{-\sigma_jt}\Big|&\le\Big(n\,p_{t_0}(x,x)\sum_{j\ge J}\varrho_j\Big)^{1/2}e^{-\sigma_J(t-t_0/2)}&&(t\ge\tfrac{t_0}2),\label{eq:tailS}\\
\Big|f_x(t)-\sum_{j<J}r_j(x)\sigma_je^{-\sigma_jt}\Big|&\le\Big(n\,p_{t_0}(x,x)\sum_{j\ge J}\varrho_j\Big)^{1/2}\sigma_J\,e^{-\sigma_J(t-t_0/2)}&&(t\ge\tfrac{t_0}2+\tfrac1{\sigma_J}).\label{eq:tailf}
\end{align}
Here $\sum_{j\ge J}\varrho_j\le\sum_{j\ge1}\varrho_j\le\beta^2\kappa$ if $\Xb>1$ (Theorem~\ref{thm:r0}), and $\sigma_J$ may be replaced by any lower bound $\Lambda\le\sigma_J$ (in \eqref{eq:tailf}: for $t\ge\frac{t_0}2+\frac1\Lambda$), e.g.\ $\Lambda=\Lambda_J$.
In discrete time with activity $q$, for integers $0\le t_1\le t$,
\begin{equation}\label{eq:taildisc}
\Big|\PP_x(T>t)-\sum_{j<J}r_j(x)\nu_j^{\,t}\Big|\le\Big(n\,P^{2t_1}(x,x)\sum_{j\ge J}\varrho_j\Big)^{1/2}\nu_*^{\,t-t_1},\qquad\nu_*=\max_{j\ge J}|\nu_j| .
\end{equation}
For activity $q$ in continuous time replace $t,t_0$ by $qt,qt_0$.
\end{theorem}
\begin{proof}
Let $\Pi=\sum_{j\ge J}\psi_j\psi_j^{\mathsf T}/\|\psi_j\|^2$, the orthogonal projector onto $\spn\{\psi_j:j\ge J\}\subset K_0$; it commutes with $L_a$. By the expansion in the proof of Theorem~\ref{thm:fpt}, $\sum_{j\ge J}r_j(x)e^{-\sigma_jt}=\ip{e_x}{e^{-tL_a}\Pi\one_{\Omega'}}=\ip{e^{-\frac{t_0}2L_a}e_x}{e^{-(t-\frac{t_0}2)L_a}\Pi\one_{\Omega'}}$. Now $\|e^{-\frac{t_0}2L_a}e_x\|^2=\ip{e_x}{e^{-t_0L_a}e_x}=\PP_x(X_{t_0}=x,\ T>t_0)\le p_{t_0}(x,x)$; the operator $e^{-sL_a}$ restricted to the range of $\Pi$ has norm $e^{-\sigma_Js}$; and $\|\Pi\one_{\Omega'}\|^2=\sum_{j\ge J}\ip{\psi_j}{\one}^2/\|\psi_j\|^2=n\sum_{j\ge J}\varrho_j$ by Proposition~\ref{prop:uniform}. Cauchy--Schwarz gives \eqref{eq:tailS}. For \eqref{eq:tailf} replace $e^{-sL_a}$ by $L_ae^{-sL_a}$, whose norm on the range of $\Pi$ is $\max_{j\ge J}\sigma_je^{-\sigma_js}=\sigma_Je^{-\sigma_Js}$ when $s\ge1/\sigma_J$ (the function $\sigma\mapsto\sigma e^{-\sigma s}$ decreases on $[1/s,\infty)$); the same monotonicity allows replacing $\sigma_J$ by $\Lambda\le\sigma_J$ if $s\ge1/\Lambda$. In discrete time, $\sum_{j\ge J}r_j\nu_j^t=\ip{Q^{t_1}e_x}{Q^{t-t_1}\Pi\one_{\Omega'}}$, $\|Q^{t_1}e_x\|^2=Q^{2t_1}(x,x)\le P^{2t_1}(x,x)$, and $\|Q^{t-t_1}\Pi\|=\nu_*^{t-t_1}$.
\end{proof}

\begin{corollary}[corner to corner, $d\ge2$]\label{cor:tailcc}
For the box, CC, $d\ge2$, $N\ge2$, $\Xb>1$, with $Z\ge\Lambda_1^2K_2$ as in Theorem~\ref{thm:cc} and $c_*=1+\frac{\pi^{3/2}}{2}=3.7841\ldots$:
\begin{enumerate}
\item[(a)] $n\,p_t(x_0,x_0)\le\big(1+N\sqrt{\pi d/(2t)}\big)^d$ for $t>0$; in particular $n\,p_{1/\Lambda_1}(x_0,x_0)\le c_*^{\,d}$.
\item[(b)] (continuous time, unit rate) for $t\ge\frac{1}{2\Lambda_1}$:
\[
\big|S_{x_0}(t)-r_0e^{-\sigma_0t}\big|\le\sqrt e\,c_*^{d/2}\,\frac{\beta\sqrt Z}{\Xb}\,e^{-\Lambda_1t},\qquad
\big|S_{x_0}(t)-r_0e^{-\sigma_0t}-r_1e^{-\sigma_1t}\big|\le e\,c_*^{d/2}\,\frac{\beta\sqrt Z}{\Xb}\,e^{-2\Lambda_1t},
\]
and for $t\ge\frac1{\Lambda_1}$: $\ \big|f_{x_0}(t)-r_0\sigma_0e^{-\sigma_0t}-r_1\sigma_1e^{-\sigma_1t}\big|\le2e\,c_*^{d/2}\,\dfrac{\beta\sqrt Z}{\Xb}\,\Lambda_1e^{-2\Lambda_1t}$.
\item[(c)] (discrete time) For every $q\in(0,1]$ and $j\ge2$: $|\nu_j|<1-2q\Lambda_1$. If $q\le\frac12$ and $t_1=\lceil\frac{1}{2q\Lambda_1}\rceil$, then for $t\ge t_1$
\[
\Big|\PP_{x_0}(T>t)-r_0\nu_0^{\,t}-r_1\nu_1^{\,t}\Big|\le c_*^{d/2}\,\frac{\beta\sqrt Z}{\Xb}\,(1-2q\Lambda_1)^{t-t_1}.
\]
For $\frac12<q\le1$ the same holds with $c_*^{d}$ replaced by $n\,P^{2t_1}(x_0,x_0)=\sum_kw_k(1-q\lambda_k)^{2t_1}$.
\end{enumerate}
\end{corollary}
\begin{proof}
(a) As in Lemma~\ref{lem:arrival}, $n\,p_t(x_0,x_0)=a_N(t/d)^d$ with $a_N(s)=\sum_{k=0}^{N-1}\varepsilon_k\cos^2(ku)e^{-(1-c_k)s}\le1+2\sum_{k\ge1}e^{-2k^2s/N^2}\le1+2\int_0^\infty e^{-2x^2s/N^2}dx=1+N\sqrt{\pi/(2s)}$, using $1-c_k=2\sin^2(ku)\ge2k^2/N^2$. For $t=1/\Lambda_1$, $s=1/(d\Lambda_1)=1/(2\sin^2u)\ge2N^2/\pi^2$, so $N\sqrt{\pi/(2s)}\le\pi^{3/2}/2$.
(b) Theorem~\ref{thm:tail} with $t_0=1/\Lambda_1$, $J=1$ ($\sigma_1>\Lambda_1$) and $J=2$ ($\sigma_2>\Lambda_2=2\Lambda_1$, Corollary~\ref{cor:box}(c); if $p=2$ the left sides vanish), together with $\sum_{j\ge1}\varrho_j\le\beta^2\kappa\le\beta^2Z/\Xb^2$.
(c) For $j\ge2$, $\nu_j=1-q\sigma_j<1-2q\Lambda_1$, and $\nu_j>1-q\Lambda_p$ where $\Lambda_p=\lambda_{(N-1,\dots,N-1)}=1+\cos\frac\pi N=2-d\Lambda_1$ is the largest eigenvalue of $L$; so $\nu_j>-(1-2q\Lambda_1)+\big[2(1-q)+q(d-2)\Lambda_1\big]\ge-(1-2q\Lambda_1)$. If $q\le\frac12$ then $0\le1-q\lambda_k\le e^{-q\lambda_k}$, hence $n\,P^{2t_1}(x_0,x_0)\le n\,p_{2qt_1}(x_0,x_0)\le c_*^d$ by (a) and monotonicity in $t$, since $2qt_1\ge1/\Lambda_1$. Apply \eqref{eq:taildisc} with $J=2$.
\end{proof}

\begin{remark}[what the tail bound says near the mode]\label{rem:tailmode}
At times $t$ with $e^{-\Lambda_1t}\asymp1/(2d\Xb)$, where the two slowest terms of $f_{x_0}$ are of the same order $\Lambda_1\Xb^{-1}$ (Theorem~\ref{thm:asym}), the neglected part is at most $2e\,c_*^{d/2}\beta\sqrt Z\,\Lambda_1\Xb^{-1}e^{-2\Lambda_1t}=O(\Lambda_1\Xb^{-3})$, i.e.\ smaller by a factor $O(\Xb^{-2})$. For $t\ge1/\Lambda_1$ this bound is absolute: it is a relative error $O(\Xb^{-2})$ only where $e^{-\Lambda_1t}\asymp1/\Xb$, that is near the mode, whereas at $t=1/\Lambda_1$ it is of the same order $\Lambda_1\Xb^{-1}$ as the two slowest terms themselves (for large $\Xb$ their sum is even negative there, since $a_1/a_0\to-2d$ and $1-2d/e<0$); so the density is a two-exponential function up to a small relative error near the mode, not on the whole window $t\ge1/\Lambda_1$ ($\Xb\sim2\pi\ln N$ in $d=2$, $\Xb\sim7.1069N$ in $d=3$). This, with Theorems~\ref{thm:cc} and \ref{thm:asym}, is the complete list of spectral inputs for a mode theorem.
\end{remark}

\section{Computer-assisted certification}\label{sec:cert}

All certified computations use Arb ball arithmetic through \texttt{python-flint} $0.9.0$ (F.~Johansson, \emph{Arb: efficient arbitrary-precision midpoint-radius interval arithmetic}, IEEE Trans.\ Comput.\ 66 (2017) 1281--1292) at $200$--$256$ bits. A comparison of two balls is accepted only if it holds for all points of the balls. Scripts are in \texttt{scripts/}, their outputs in \texttt{data/} and \texttt{logs/}.

\paragraph{Constants} (scripts \texttt{06\_certify\_constants.py} and \texttt{09\_certify\_part3.py}).\enspace
$\sup|\Psi''|\le2.35$ (Lemma~\ref{lem:Psi}(d)); $\Lambda_*$, $\Sigma_*$ (series with explicit geometric tail bounds; the closed forms of Lemma~\ref{lem:closed} overlap); $c_\tau$, $c_{m\tau}$, $C_2$; $S_1,\dots,S_4$; elementary enclosures of $Z_2$, $Z_3$; $C_3$ (Lemma~\ref{lem:C3}, split point $S=10^6$, cross-checked with $S=200$); $K(1)$ (rigorous integration of the analytic integrand $4\cos^2\varphi\sqrt{(2+\sin^2\varphi)/(1+\sin^2\varphi)}$ on $[0,\frac\pi2]$); $B_D$ (partial sum over $|k|_\infty\le30$ plus the tail bound $8e^{-62\kappa_3}/(1-e^{-2\kappa_3})^2$); $B_e$; $E_\infty$, $c_{\tau3}$, $c_{m\tau3}$ (lattice sums with exponential tail bounds and rigorous integration). The script asserts that the rounded constants quoted in Theorems~\ref{thm:d2} and~\ref{thm:d3}(b) and in Corollary~\ref{cor:Xbar} are implied by the enclosures. Balls passed from one script to the next as decimal strings are re-inflated to cover the rounding of the strings.

\paragraph{Epstein sums (\texttt{14\_certify\_epstein.py}).} $Z_2$, $Z_3$, $Z_2'$, $Z_3'$ by Lemma~\ref{lem:epstein} ($300$ bits), with the three cross-checks described there.

\paragraph{Second order in $d=3$ (\texttt{16\_certify\_d3\_remainder.py}).} $K(1)=6\mathbf K(\frac12)-4\mathbf E(\frac12)$ (agrees with the quadrature value); $V_g$ (Lemma~\ref{lem:Jb2}(d)); $c_\delta$, $c_a$ (Lemma~\ref{lem:quant}); $\delta_\infty$, $\alpha_\infty$, $c^{(0)}_{\tau3}$, $c^{(0)}_{m\tau3}$, $C_3''$ (Theorem~\ref{thm:d3}(d)); the rounded constants of Theorem~\ref{thm:d3}(c) and of Corollary~\ref{cor:Xbar}. As a cross-check the script also bounds $\int_0^\pi|J_b''|\le2.651$, which gives the weaker first-order estimate $|\varepsilon_b(N)-\varepsilon_b(\infty)|\le0.694/N$.

\paragraph{Poles and residues (\texttt{07\_certify\_poles.py}, \texttt{09\_certify\_part3.py}).}
For CC, $d=2$ ($N\le300$) and $d=3$ ($N\le50$; $\tau$, $m$ also for $N=70,100$), the functionals $\tau,m,K_2$ are evaluated as exact spectral sums \eqref{eq:ccsums} in ball arithmetic. A pole is certified as follows: by Theorem~\ref{thm:structure}(i), $G$ is strictly increasing on $(\Lambda_j,\Lambda_{j+1})$ and has exactly one zero there; if $\Lambda_j<s_-<s_+<\Lambda_{j+1}$ and $G(s_-)<0<G(s_+)$ are verified with balls, then $\sigma_j\in(s_-,s_+)$. The brackets have relative width $2^{-149}$. The residues $r_j=-G_{x_0}(\sigma_j)/(\sigma_jG'(\sigma_j))$ are then evaluated on the bracket. For each $N$ the scripts verify with certified comparisons: the enclosures of Theorems~\ref{thm:d2}, \ref{thm:d3}(b),(c) and both forms of Corollary~\ref{cor:Xbar}; Lemmas~\ref{lem:arrival}, \ref{lem:K2}, \ref{lem:taugen}; $r_0>1$, $\sigma_0m>1$, $r_1<0$ (Theorem~\ref{thm:r0gt1}, Corollary~\ref{cor:r1cc}); every inequality of Theorems~\ref{thm:sigma0}--\ref{thm:r1size} with the exact functionals; and every inequality of Theorem~\ref{thm:cc} and Corollary~\ref{cor:explicitN} whose hypothesis holds ($10$ to $38$ certified comparisons per case, $37$ cases, all passed). Tables~\ref{tab:cc2}--\ref{tab:cc3x} show a selection.

\paragraph{Signs of the residues (\texttt{13\_certify\_residue\_signs.py}).} Proposition~\ref{prop:signs}; exact grouping of the eigenvalues modulo the cyclotomic polynomial, $300$-bit balls; as a consistency check the script also asserts the statements of Theorem~\ref{thm:noalt} (which is proved without computation) for each of its $20$ cases, including the exact purity and parity of the two eigenvalues used in part (b).

\paragraph{Limits in $d=2$ (\texttt{19\_certify\_d2\_limits.py}).} $S_5,\dots,S_8$, $c^{(0)}_\tau$, $c^{(0)}_{m\tau}$, $C_2^{(0)}$ (Theorem~\ref{thm:d2limits}); comparison with the closed forms of Remark~\ref{rem:d2closed} (agreement to $10^{-50}$); certified values of $R_\tau(N)$, $R_{m\tau}(N)$ for $22$ values of $N\le10^6$ (inside the proved windows, increasing, below the limits) and certified strict monotonicity for all consecutive $N\le300$; floating-point checks of every step of the proof (Lemma~\ref{lem:trap}, the majorants, the pointwise limits, the four sums).

\paragraph{Truncated constants (\texttt{20\_certify\_rounded\_constants.py}).} Proposition~\ref{prop:roundedfalse}: the thresholds $8872$ and $182\,420$, the monotonicity conditions, and the direct evaluations in $d=2$.

\paragraph{Printed decimals (\texttt{21\_check\_printed\_decimals.py}).} Every decimal string of the text that ends in dots is compared with a $40$-digit value of the constant and must be its truncation.

\paragraph{Sanity checks of every lemma (floating point).}
\texttt{01}: Theorems~\ref{thm:structure}--\ref{thm:residues} against brute-force diagonalisation of $L_a$ for $29$ geometries in $d=1,2,3$ (corner, centre, generic and degenerate targets), including the time-domain formulas \eqref{eq:cont}, \eqref{eq:disc}. \texttt{03}: Section~\ref{sec:bounds}. \texttt{02}, \texttt{04}, \texttt{05}: Propositions~\ref{prop:d2exact}, \ref{prop:d3exact}, Lemmas~\ref{lem:cycle}, \ref{lem:cot}, \ref{lem:Psi}, \ref{lem:trap} and the remainder pieces of Theorem~\ref{thm:d2}. \texttt{08}: Proposition~\ref{prop:uniform}, Lemmas~\ref{lem:elem}--\ref{lem:taugen}, Corollary~\ref{cor:alld} ($d=3,4,5$), Lemmas~\ref{lem:harm}, \ref{lem:KJ}--\ref{lem:small3}, \ref{lem:C3}, the enclosures of Theorems~\ref{thm:d2}, \ref{thm:d3} up to $N=5000$ ($d=2$) and $N=200$ ($d=3$), and Theorem~\ref{thm:tail}/Corollary~\ref{cor:tailcc} against brute force (the ratio remainder/bound is between $0.01$ and $0.15$; discrete time checked for $q=0.3,\,0.5,\,0.8,\,1$). \texttt{10}: the limits of Theorem~\ref{thm:asym}. \texttt{12}: Theorem~\ref{thm:r1} and Remark~\ref{rem:signrule} (formula \eqref{eq:rsign}, $\eta_1\le\Lambda_1$, $G^+(\eta_1)=0$, and the sign rule for every pole). \texttt{15}: Proposition~\ref{prop:qsd}, the lower bound of Lemma~\ref{lem:arrival}, Lemma~\ref{lem:coupling} (continuous and discrete time, $d=1,\dots,4$), Theorem~\ref{thm:r0gt1}, and the sharpened Theorem~\ref{thm:cc}(ii),(v),(vi) up to $N=640$ ($d=2$), $N=64$ ($d=3$). \texttt{17}: Lemmas~\ref{lem:Kclosed}--\ref{lem:quant} (closed form of $K$ against quadrature, the formulas for $M'$, $M''$, $J_b''$, the comparison inequalities at $24\,000$ random points) and Theorem~\ref{thm:d3}(c),(d) for $N\le800$. \texttt{18}: Proposition~\ref{prop:pure} ($3282$ pure eigenvalues of $352$ random chains with an involution and $2986$ of boxes with arbitrary symmetric pairs), Remark~\ref{rem:signpoly} ($198$ random start--target pairs: degree and leading coefficient of $\mathcal N_x$, both sign relations, the three counting consequences), and every spectral claim in the proof of Theorem~\ref{thm:noalt} (part (b) for $104$ pairs $(d,N)$ with $d\ge2$, $N\ge3$, part (c) for $91$ of them) together with floating-point sign patterns for $49$ pairs (the statistics of Observation~\ref{obs:signs} are asserted there).

\begin{table}[ht]\centering\footnotesize
\begin{tabular}{rrccccc}\toprule
$N$ & $\Xb$ & $[X_-,X_+]$ & $\sigma_0\tau$ & bounds (i) & $\sigma_0m$ & bounds (ii)\\\midrule
5 & 8.1304 & [7.839, 8.665] & 0.928025 & [0.907421, 0.955961] & 1.16428 & [1.06743, 1.23475]\\
10 & 12.5893 & [12.515, 12.759] & 0.964332 & [0.961349, 0.977113] & 1.12909 & [1.10737, 1.16149]\\
20 & 16.9972 & [16.979, 17.048] & 0.979370 & [0.978720, 0.986684] & 1.10446 & [1.10027, 1.11600]\\
50 & 22.7773 & [22.774, 22.787] & 0.988292 & [0.988133, 0.992421] & 1.08273 & [1.08214, 1.08764]\\
100 & 27.1371 & [27.136, 27.140] & 0.991723 & [0.991644, 0.994629] & 1.07129 & [1.07112, 1.07451]\\
200 & 31.4937 & [31.493, 31.494] & 0.993849 & [0.993802, 0.996000] & 1.06256 & [1.06250, 1.06488]\\
300 & 34.0416 & [34.042, 34.042] & 0.994735 & [0.994698, 0.996572] & 1.05837 & [1.05832, 1.06033]\\
\bottomrule\end{tabular}
\\[6pt]\begin{tabular}{rcccccc}\toprule
$N$ & $\delta=\sigma_1/\Lambda_1-1$ & bounds (iii) $[\delta_-,\delta_+]$ & $r_0$ & bounds (v) & $r_1$ & bounds (vi)\\\midrule
5 & 0.40891 & [0.03923, 0.89586] & 1.18370 & [0.89247, 1.40406] & -0.29622 & [-19.90502, --]\\
10 & 0.30187 & [0.19978, 0.45562] & 1.14042 & [1.05136, 1.23052] & -0.25470 & [-0.92575, --]\\
20 & 0.23498 & [0.19848, 0.30601] & 1.11172 & [1.07268, 1.15328] & -0.21108 & [-0.42176, -0.04962]\\
50 & 0.17929 & [0.16496, 0.21282] & 1.08723 & [1.06791, 1.10810] & -0.16718 & [-0.25043, -0.08339]\\
100 & 0.15128 & [0.14262, 0.17284] & 1.07461 & [1.06144, 1.08881] & -0.14318 & [-0.19350, -0.08739]\\
200 & 0.13057 & [0.12477, 0.14548] & 1.06511 & [1.05548, 1.07543] & -0.12477 & [-0.15828, -0.08525]\\
300 & 0.12082 & [0.11608, 0.13315] & 1.06058 & [1.05239, 1.06933] & -0.11594 & [-0.14321, -0.08288]\\
\bottomrule\end{tabular}
\caption{$d=2$, corner to corner: certified values (ball arithmetic, all digits shown are correct up to rounding) and the bounds of Corollary~\ref{cor:explicitN}, which depend on $N$ only ($X_\pm$, $M_\pm$ from Corollary~\ref{cor:Xbar}; $Z$ is the certified upper end $6.0268120396919401\ldots$ of $Z_2$ from Lemma~\ref{lem:epstein}, not the rounded value $6.02682$, with which the bounds change in the last digit shown at most, e.g.\ the lower bound for $r_1$ at $N=5$ becomes $-19.90509$). A dash means that the hypothesis of the corresponding bound is not satisfied.}\label{tab:cc2}\end{table}

\begin{table}[ht]\centering\footnotesize
\begin{tabular}{rrccccc}\toprule
$N$ & $\Xb$ & $[X_-,X_+]$ & $\sigma_0\tau$ & bounds (i) & $\sigma_0m$ & bounds (ii)\\\midrule
4 & 18.7009 & [16.814, 21.508] & 0.973816 & [0.944530, 0.989168] & 1.09643 & [1.00884, 1.14409]\\
6 & 33.0422 & [32.180, 34.614] & 0.989375 & [0.984045, 0.995371] & 1.06265 & [1.02848, 1.07958]\\
10 & 61.6742 & [61.360, 62.474] & 0.996421 & [0.995576, 0.998505] & 1.03671 & [1.02131, 1.04353]\\
20 & 132.9573 & [132.878, 133.304] & 0.999146 & [0.999058, 0.999665] & 1.01803 & [1.01537, 1.02062]\\
30 & 204.1091 & [204.074, 204.328] & 0.999626 & [0.999601, 0.999857] & 1.01195 & [1.01118, 1.01292]\\
50 & 346.3172 & [346.305, 346.443] & 0.999867 & [0.999862, 0.999950] & 1.00714 & [1.00697, 1.00738]\\
\bottomrule\end{tabular}
\\[6pt]\begin{tabular}{rcccccc}\toprule
$N$ & $\delta=\sigma_1/\Lambda_1-1$ & bounds (iii) $[\delta_-,\delta_+]$ & $r_0$ & bounds (v) & $r_1$ & bounds (vi)\\\midrule
4 & 0.22154 & [0.07462, 0.39767] & 1.10263 & [0.93107, 1.21880] & -0.20422 & [-1.97022, --]\\
6 & 0.14461 & [0.09194, 0.19856] & 1.06516 & [1.01016, 1.10108] & -0.14124 & [-0.41132, --]\\
10 & 0.08574 & [0.07144, 0.10204] & 1.03755 & [1.01662, 1.04956] & -0.08567 & [-0.14382, -0.03428]\\
20 & 0.04255 & [0.03987, 0.04627] & 1.01823 & [1.01440, 1.02191] & -0.04273 & [-0.05342, -0.03136]\\
30 & 0.02829 & [0.02723, 0.02991] & 1.01204 & [1.01077, 1.01346] & -0.02840 & [-0.03274, -0.02362]\\
50 & 0.01694 & [0.01659, 0.01751] & 1.00717 & [1.00683, 1.00757] & -0.01699 & [-0.01845, -0.01535]\\
\bottomrule\end{tabular}
\caption{$d=3$, corner to corner: certified values and the bounds of Corollary~\ref{cor:explicitN} (functions of $N$ only), with $X_\pm$, $M_\pm$ from Corollary~\ref{cor:Xbar}, i.e.\ from Theorem~\ref{thm:d3}(c), and $Z$ the certified upper end of $Z_3$ from Lemma~\ref{lem:epstein}. A dash means that the hypothesis of the corresponding bound is not satisfied.}\label{tab:cc3}\end{table}

\begin{table}[ht]\centering\footnotesize
\begin{tabular}{rrccccc}\toprule
$N$ & $\Xb$ & $1-\sigma_0\tau$ & bounds (i) on $1-\sigma_0\tau$ & $\delta$ & $[\delta_-,\delta_+]$ & $r_1$ and bounds (vi)\\\midrule
4 & 18.701 & 2.618e-02 & [1.423e-02, 4.550e-02] & 0.22154 & [0.08550, 0.34685] & -0.20422 [-1.39532, --]\\
6 & 33.042 & 1.063e-02 & [5.076e-03, 1.515e-02] & 0.14461 & [0.09502, 0.19267] & -0.14124 [-0.38882, --]\\
10 & 61.674 & 3.579e-03 & [1.534e-03, 4.379e-03] & 0.08574 & [0.07218, 0.10148] & -0.08567 [-0.14262, -0.03485]\\
20 & 132.957 & 8.544e-04 & [3.371e-04, 9.405e-04] & 0.04255 & [0.03996, 0.04624] & -0.04273 [-0.05357, -0.03144]\\
30 & 204.109 & 3.742e-04 & [1.436e-04, 3.985e-04] & 0.02829 & [0.02725, 0.02990] & -0.02840 [-0.03284, -0.02363]\\
50 & 346.317 & 1.332e-04 & [4.997e-05, 1.382e-04] & 0.01694 & [0.01660, 0.01751] & -0.01699 [-0.01849, -0.01534]\\
\bottomrule\end{tabular}
\caption{$d=3$: bounds of Theorem~\ref{thm:cc} evaluated with the exact (certified) $\Xb$ and with $Z$ the certified upper end of $Z_3$ (Lemma~\ref{lem:epstein}).}\label{tab:cc3x}\end{table}


\section{Numerical observations and what is not proved}\label{sec:gaps}

\begin{proposition}[certified sign patterns; the residues do not alternate for $d\ge2$, $N\ge3$]\label{prop:signs}
For the corner-to-corner geometry the signs of $(r_0,r_1,\dots,r_{p-1})$ are (computer-assisted, ball arithmetic; $S$ = number of sign changes, $D-1=d(N-1)-1$ the lower bound of Theorem~\ref{thm:residues}(d), $p-1$ the number needed for alternation):
\begin{center}\footnotesize
\begin{tabular}{llrrrl}
$(d,N)$ & signs of $r_0,\dots,r_{p-1}$ & $p$ & $S$ & $D-1$ & negative\\\hline
$(2,2)$ & \texttt{+-} & 2 & 1 & 1 & 1\\
$(2,3)$ & \texttt{+--+-} & 5 & 3 & 3 & 3\\
$(2,4)$ & \texttt{+-+++-+-} & 8 & 5 & 5 & 3\\
$(2,5)$ & \texttt{+-++---++--+-} & 13 & 7 & 7 & 7\\
$(2,6)$ & \texttt{+-++--++--++--+-} & 16 & 9 & 9 & 8\\
$(2,7)$ & \texttt{+-++-+-+-+--+-++---++--+-} & 25 & 17 & 11 & 13\\
$(2,8)$ & \texttt{+-++-+-+-+--+-++-+--+++---++--+-} & 32 & 21 & 13 & 16\\
$(2,10)$ & \texttt{+-++-+-+-+--+-++-+---+-+++-+--++++----+++---++--+-} & 50 & 29 & 17 & 25\\
$(2,11)$ & \texttt{+-++-+-+-+--+-++-++--+--++------+-++-----+++-+--+-++---++--+-} & 61 & 35 & 19 & 34\\
$(2,12)$ & \texttt{+-++-+-+-+--+-++-+-+--+--+--+++++-+--+++++--+--+++---++---++--+-} & 64 & 37 & 21 & 31\\
$(3,2)$ & \texttt{+-+} & 3 & 2 & 2 & 1\\
$(3,3)$ & \texttt{+-++--+-} & 8 & 5 & 5 & 4\\
$(3,4)$ & \texttt{+-+-+-+-++--+-+} & 15 & 12 & 8 & 7\\
$(3,5)$ & \texttt{+-+-+-+-+++-+---++++---++--+-} & 29 & 17 & 11 & 14\\
$(3,6)$ & \texttt{+-+-+-+++---+++---+++---++---++-+} & 33 & 16 & 14 & 16\\
$(4,2)$ & \texttt{+-+-} & 4 & 3 & 3 & 2\\
$(4,3)$ & \texttt{+-+--++--+-} & 11 & 7 & 7 & 6\\
$(4,4)$ & \texttt{+-+-++-++-+-+-+-+--++-+-} & 24 & 19 & 11 & 11\\
$(5,2)$ & \texttt{+-+-+} & 5 & 4 & 4 & 2\\
$(5,3)$ & \texttt{+-+--+--++--+-} & 14 & 9 & 9 & 8
\end{tabular}
\end{center}
In particular all residues are non-zero in these cases; for $N=2$ the sequence alternates and for $N\ge3$ it does not ($S<p-1$), in agreement with Theorem~\ref{thm:noalt} (which does not depend on this proposition). In every row with $N\ge3$ the three signs $(\operatorname{sign}r_{i-1},\operatorname{sign}r_i,\operatorname{sign}r_{i+1})$ named in Theorem~\ref{thm:noalt}(b) are \texttt{--+} when $\Lambda_i,\Lambda_{i+1}$ are purely even and \texttt{++-} when they are purely odd. The lower bound $D-1$ of Theorem~\ref{thm:residues}(d) is attained in the rows with $N\le3$, for $d=2$, $N\le6$, but not in general; about one half of the residues are negative.
\end{proposition}
\begin{proof}
Script \texttt{13\_certify\_residue\_signs.py}. The distinct eigenvalues are determined exactly: with $\zeta=e^{i\pi/N}$, $2\sum_i\cos(\pi k_i/N)=\sum_i(\zeta^{k_i}+\zeta^{2N-k_i})$, and two such sums are equal if and only if the corresponding integer polynomials agree modulo the cyclotomic polynomial $\Phi_{2N}$ (exact arithmetic). The distinct values are then separated by ball arithmetic, every zero $\sigma_j$ of $G$ is bracketed by a certified sign change of $G$ (Theorem~\ref{thm:structure}(i)), and the sign of $r_j=-G_{x_0}(\sigma_j)/(\sigma_jG'(\sigma_j))$ is certified on the bracket. The script also asserts, for every row with $N\ge3$, the purity and parity of $\Lambda_i,\Lambda_{i+1}$ (exact) and the monotonicity of the three signs of Theorem~\ref{thm:noalt}(b), and for every row covered by Theorem~\ref{thm:noalt}(c) that the pair named there does not have the excluded signs.
\end{proof}

\begin{observation}[signs of the residues, larger $N$]\label{obs:signs}
In floating-point arithmetic (script \texttt{18}; $49$ pairs $(d,N)$ with $p\le140$, among them $d=2$, $N\le16$ and $d=3$, $N\le9$) the number of negative residues stays between $0.33\,p$ and $0.6\,p$, the number $S$ of sign changes stays strictly between $D-1$ and $p-1$ for $7\le N\le16$ ($d=2$) and $4\le N\le9$ ($d=3$), and for $N=3$ the lower bound is attained, $S=D-1=2d-1$, in every dimension examined ($2\le d\le11$). That the residues do not alternate for \emph{every} $d\ge2$, $N\ge3$ is Theorem~\ref{thm:noalt}; the exact pattern for a given $(d,N)$ is given by Remark~\ref{rem:signrule} and Proposition~\ref{prop:pure}, and no closed description of it is known to us.
\end{observation}

\paragraph{What is not proved in this note.}
\begin{enumerate}
\item[(G1)] \emph{The mode.} No statement about the location of the maximum of $f$ or about unimodality is proved here. The spectral inputs are complete (Theorems~\ref{thm:cc}, \ref{thm:asym}, \ref{thm:tail}): for $t\ge1/\Lambda_1$ the density is an explicit two-exponential function up to an absolute error $O(\Lambda_1\Xb^{-1}e^{-2\Lambda_1t})$, a relative error $O(\Xb^{-2})$ near the mode (Remark~\ref{rem:tailmode}). What is missing for a mode theorem is (i) a quantitative analysis of the critical point of the two-exponential function under this perturbation, and (ii) control of $f$ on $[0,1/\Lambda_1]$, where all poles contribute (Theorem~\ref{thm:tail} does not apply there, and the elementary bound $f(t)\le\frac1n(\frac12n\,p_t(x_0,a)+\partial_t\,n\,p_t(x_0,a))$ is too weak by a factor of order $\ln N$ in $d=2$). In $d=2$ every rate of convergence is only logarithmic in $N$ ($\Xb\sim2\pi\ln N$), so explicit thresholds for $N$ would be astronomically large; this is intrinsic, not an artefact.
\item[(G2)] \emph{Sharpness in $d=3$.} For every $N\ge2$, $R_3(N)\in[-17.4,19.4]$ and $R_3'(N)\in[-82.4,80.4]$ are proved, and the limits $1.7574$ and $-0.2989$ are proved; the true remainders stay within $[1.61,1.76]$ and $[-0.63,-0.29]$ (Remark~\ref{rem:d3sharp}). The loss is in the comparison constants $c_\delta$, $c_a$ (worst-case exponent $\kappa_3$) and in the range of $\rho_N$; no rate is proved for the convergence to the limits (this would need a third-order analysis of $P_N-P_\infty$ and of $g'$).
\item[(G3)] \emph{Sharpness in $d=2$.} $R_\tau\in[-2.14,3.53]$ is proved for every $N$ and the limit $R_\tau\to0.8830352\ldots$ is proved (Theorem~\ref{thm:d2limits}); $R_\tau(N)\in[0.8512,0.8831]$ for all $N$, the monotonicity in $N$ and the rate $O(N^{-2})$ are only observed. An explicit rate would follow from the same proof with second-order remainders (a certified bound for $\Psi'''$ and second-order versions of \eqref{eq:yk}); it was not carried out.
\item[(G4)] \emph{Computer assistance.} The numerical constants listed in Section~\ref{sec:cert} are certified by interval arithmetic, not by hand. The structure theorems (Sections~\ref{sec:structure}--\ref{sec:bounds}), Lemmas~\ref{lem:arrival}--\ref{lem:K2}, \ref{lem:taugen}, \ref{lem:coupling}, Theorem~\ref{thm:r0gt1}, Theorem~\ref{thm:cc} with $Z=Z_d(N)$, the limit statements of Theorem~\ref{thm:asym} and Theorem~\ref{thm:tail} do not depend on any computation.
\item[(G5)] \emph{Other geometries.} Sections~\ref{sec:structure}--\ref{sec:bounds} and Theorem~\ref{thm:tail} hold for every start and target (and every finite symmetric chain). The lattice-sum evaluations of Sections~\ref{sec:corner}--\ref{sec:explicit} are carried out for opposite corners only. For a centre target the first visible eigenvalue is $\Lambda_1=\frac1d(1-\cos\frac{2\pi}N)$ and for a centre start with corner target $A_{x_0}=0$ and the residue $r_1$ changes sign (Supplementary Section~S7 of the article); the general theorems apply but the functionals were not evaluated.
\item[(G6)] \emph{Quality of the lower bound for $\delta$ and of the bounds for $r_1$.} They contain the factor $\omega=2(1+\delta_+)/(1-\delta_+)$ and are informative only when $\delta_+$ is small ($N\ge20$ in $d=2$, $N\ge10$ in $d=3$; see the tables).
\item[(G7)] \emph{Discrete time, $q>\frac12$.} Corollary~\ref{cor:tailcc}(c) is explicit for $q\le\frac12$; for $\frac12<q\le1$ the prefactor $n\,P^{2t_1}(x_0,x_0)$ is left as a finite sum (near-bipartite oscillations of the non-lazy walk).
\item[(G8)] \emph{Classical identities.} The closed forms of Lemma~\ref{lem:closed} rest on the cited DLMF formulas (theta-function and elliptic-integral identities at the lemniscatic point); the bounds of Theorem~\ref{thm:d2} do not. Theorem~\ref{thm:d3}(c),(d) rests on the classical derivative formulas and the inequality $\ln4\le\mathbf K+\ln k'\le\frac\pi2$ for complete elliptic integrals \eqref{eq:ellfacts}, on the Fourier series of $B_2$ and of $\ln\sin\frac\theta2$, and on Arb's evaluation of $\mathbf K$, $\mathbf E$, $K_\nu$. The identification of $C_3$ with lattice Green's function values is elementary, but the closed form of Watson's integral $u(0)$ is quoted from the literature and not used. The values of $Z_3$, $Z_2'$, $Z_3'$ rest on Jacobi's theta transformation (cited) and on Arb's incomplete gamma function.
\end{enumerate}

\begin{thebibliography}{99}\small
\bibitem{AS} M.~Abramowitz and I.~A.~Stegun, \emph{Handbook of Mathematical Functions}, National Bureau of Standards, 1964 (formula 4.3.68).
\bibitem{AB} D.~J.~Aldous and M.~Brown, Inequalities for rare events in time-reversible Markov chains I, in: \emph{Stochastic Inequalities}, IMS Lecture Notes--Monograph Series 22 (1992), 1--16.
\bibitem{DLMF} \emph{NIST Digital Library of Mathematical Functions}, \texttt{https://dlmf.nist.gov}, formulas 10.32.3, 10.32.9, 19.4.1, 19.4.2, 19.9.1, 19.12.1, 20.4.6, 20.5.3, 20.7.5, 20.7.32, 20.9.1, 20.9.2, 23.5.5, 23.15.9, 23.17.8, 24.8.1 (all checked against the online edition on 1 October 2026).
\bibitem{EW} J.~W.~Essam and F.~Y.~Wu, The exact evaluation of the corner-to-corner resistance of an $M\times N$ resistor network: asymptotic expansion, J.~Phys.~A 42 (2009) 025205.
\bibitem{GZ} M.~L.~Glasser and I.~J.~Zucker, Extended Watson integrals for the cubic lattices, Proc.\ Natl.\ Acad.\ Sci.\ USA 74 (1977) 1800--1801.
\bibitem{HJ} R.~A.~Horn and C.~R.~Johnson, \emph{Matrix Analysis}, 2nd ed., Cambridge University Press, 2013 (Theorems 4.3.17, 8.3.1, 8.4.4).
\bibitem{Jo1} F.~Johansson, Arb: efficient arbitrary-precision midpoint-radius interval arithmetic, IEEE Trans.\ Comput.\ 66 (2017) 1281--1292.
\bibitem{Jo2} F.~Johansson, Numerical integration in arbitrary-precision ball arithmetic, in: Mathematical Software -- ICMS 2018, LNCS 10931, 255--263.
\bibitem{Wa} G.~N.~Watson, Three triple integrals, Quart.\ J.\ Math.\ Oxford 10 (1939) 266--276.
\bibitem{ZR} I.~J.~Zucker and M.~M.~Robertson, Exact values of some two-dimensional lattice sums, J.~Phys.~A 8 (1975) 874--881.
\end{thebibliography}

\end{document}
